% =============================================================================
% draft_v3_onlineappendix.tex -- Online Appendix to
% "Liquidity, Activism Disclosure, and Takeover Premia"
%
% Compilation (from the repo root, after draft_v3.tex has been built, since
% cross-references into the main text resolve through draft_v3.aux):
%   xelatex -interaction=nonstopmode draft_v3_onlineappendix.tex
%   xelatex -interaction=nonstopmode draft_v3_onlineappendix.tex
%
% Content: the extended discussion of the standing assumptions and the
% numerical records that measure them against the implemented calibration
% (both moved from draft_v3.tex on 2026-08-30), the August 2026 empirical
% record preserved exactly as viewed (added 2026-09-01; the record itself was
% deleted from the live tree the same day and lives in history at 9b98089),
% then the proofs of every result stated in the main text (moved verbatim
% from draft_v3.tex on 2026-08-30). Numbered objects cited without a section
% letter (Assumption 3, Lemma 2, equation (8)) refer to the main text.
% =============================================================================

\documentclass[11pt]{article}

\usepackage[margin=1in]{geometry}
\usepackage{amsmath,amssymb,amsthm}
\usepackage{setspace}
\usepackage{enumitem}
\usepackage{array}
\usepackage{booktabs}
\usepackage{graphicx}
\usepackage{float}
\usepackage{xr-hyper}
\usepackage[hidelinks]{hyperref}
\externaldocument{draft_v3}

\setstretch{1.1}

% Theorem environments (statement numbers of results live in the main text;
% the numerical-record remarks of the appendix are numbered locally)
\newtheorem{theorem}{Theorem}
\newtheorem{proposition}{Proposition}
\newtheorem{lemma}{Lemma}
\newtheorem{corollary}{Corollary}
\newtheorem{definition}{Definition}
\newtheorem{assumption}{Assumption}
\newtheorem{remark}{Remark}

% Math macros (identical to draft_v3.tex; \Tmap is always the calligraphic
% outer best-response map, never the filing window T)
\newcommand{\Eop}{\mathbb{E}}
\newcommand{\Prb}{\mathbb{P}}
\newcommand{\ind}{\mathbf{1}}
\newcommand{\Dact}{\Delta^{\mathrm{act}}}
\newcommand{\Tmap}{\mathcal{T}}
\newcommand{\Atau}{\mathrm{A}(\tau)}
\newcommand{\Abr}{\mathrm{A}(\mathrm{br})}
\newcommand{\Akap}{A'_{\kappa}}
\newcommand{\ND}{\mathrm{ND}}
\newcommand{\Sfl}{\mathsf{S}_F}

\title{\textbf{Online Appendix to ``Liquidity, Activism Disclosure, and Takeover
Premia''}\\[0.3em]
\large Discussion of assumptions, numerical records, and proofs}
\author{Austin Li\\[0.25em]}
\date{This draft: 30 August 2026}

\begin{document}
\maketitle

\noindent This appendix carries the material the main text compresses out of the reader's path.
Section~\ref{sec:app-assumptions} discusses the standing assumptions; Section~\ref{sec:app-records}
carries the numerical records that measure them against the implemented calibration;
Section~\ref{app:honesty} preserves the August 2026 empirical record --- a registered design the
paper has superseded --- exactly as viewed, with its estimates, intervals, MDEs and gate verdicts;
Sections~\ref{sec:proofs-core}, \ref{sec:proofs-existence} and~\ref{sec:proofs-ge} give the proofs
of every result stated in the main text. A numbered object cited without a section letter, such as
an assumption, a lemma or an equation, refers to the main text.

\bigskip

\appendix
%==============================================================================
\section{Discussion of assumptions}
\label{sec:app-assumptions}
%==============================================================================

\noindent This section collects the extended discussion of Assumptions 3--6 of the main text that
the body of the paper compresses: what each clause is for, what is derived rather than assumed,
where the satisfiability boundary sits, and which executed checks bear on it. References to
sections, results and equations without a lettered section are to the main text.

\paragraph{The inner-pricing regularity clause.}
Most of Assumption~\ref{asm:A5} is a theorem rather than an
assumption. Under \eqref{eq:m0}, existence, uniqueness and continuity of the \emph{inner} fixed point
are proved, not assumed. With the belief summaries $\bar y(\mathcal I) = \Eop_\mu[v] + \pi\Delta_V$
and $\bar m(\mathcal I) = m_0 + \pi\Delta_m$, the pricing map reduces to the scalar equation
$P \mapsto \bar y + \bar m\,p(P)/(1-p(P))$, which is strictly decreasing in $P$ wherever
$\bar m \ge 0$ and therefore crosses the identity exactly once. Three independent calculations
support this, one of them also recording two counterexamples: one producing zero roots and another
producing three, once $m_0 < 0$. What Assumption~\ref{asm:A5} is retained for is its \emph{continuity} clause: the
pricing family is continuous in the cutoff vector and the parameters, and measurable in the flagged
tuple. Because the flagged information sets are continuum-indexed, ``unique fixed point'' must be
read as a measurably selected family, not a finite list. Where a result below cites
Assumption~\ref{asm:A5} for existence or uniqueness, it may cite \eqref{eq:m0} instead.

Two continuities must be kept apart here, since only one of them is proved. Continuity of the inner
root \emph{in its belief summaries} $(\hat v,\pi)$ follows from \eqref{eq:m0}. Continuity of the
\emph{composition} $k \mapsto (\hat v,\pi) \mapsto P$ in the cutoff vector is what the clause above
retains as an assumption: the $k$-dependence runs through the conditioning rather than through the
pricing map, and no step derives it. Remark~\ref{rem:A6record} records that composition measured
discontinuous at the implemented calibration, at the same loci and by the same checks, so this
clause carries adverse applicability evidence of its own. None of that moves any result:
Assumption~\ref{asm:A5} is a hypothesis, and no result below consumes its cutoff clause.

\paragraph{Identification on flagged histories.}
The weak identification wording of Assumption~\ref{asm:A7} is not enough for Lemma~\ref{lem:L2}: it permits two pairs $(j,s)$
with different pooled paths, which is exactly that lemma's first failure case. Two injective forms
are therefore named separately, and the form is named every time either is used.

Assumption~\ref{asm:A7J} is strictly stronger than Assumption~\ref{asm:A7p}, and it is the form the
existence argument consumes. Strictness of $B^F$ alone is neither necessary (it fails at
crossing-date jumps on the pro-rata menu) nor sufficient, because of multi-Voice backtracking.
Under Assumption~\ref{asm:A7p} the flagged tuple is continuum-valued as a tuple, since injectivity
forces $(B^F,Q^F)$ to be continuum-valued, though the coordinates may trade the burden between them;
injectivity plus measurability already gives the measurable inverse on standard Borel spaces, so no
separate assumption is needed for that. Satisfiability is resolved for Assumption~\ref{asm:A7p}:
together with a fixed cutoff policy and $\Omega > 0$ it delivers the on-path injective form on
positive-probability flagged tuples with an explicit inverse, and a satisfying menu exists: the
pro-rata single-Voice menu with terminal target strictly increasing on all of $\mathbb R$, which also
satisfies Assumption~\ref{asm:A7J}. Assumption~\ref{asm:A7J} additionally needs $b^*$ strictly
increasing off the Voice region: a target flat below the Voice cutoff breaks it, in an executed
check producing forty collisions, while leaving Assumption~\ref{asm:A7p} intact. The failure boundary
is a binding stake cap, quantized stakes, a composed target repeating values
across Voice-plan switches, $\Omega = 0$, and policy dependence when the condition is stated only at
one equilibrium's cutoffs. Menus satisfying Assumption~\ref{asm:A7p} are fully separating on the
flagged set, so the burden moves to the incentive compatibility of Proposition~\ref{prop:P1}, not
away.

\paragraph{Where the chord restriction bites.}
Where the bite of Assumption~\ref{asm:Atau} actually sits is worth stating, because it is narrower
than the display suggests. Given a $\kappa$-invariant three-point support, the derivative
restrictions $A_0' = A_1' = \Akap$ and $A_{1/2}' = -2\Akap$ are not an extra assumption at all: they
are \emph{equivalent} to $\kappa$-invariance of the pooled block's total mass and of its
unnormalised engagement moment, both of which the model delivers at fixed policies. The entire
remaining content of Assumption~\ref{asm:Atau} is the support condition. A one-round ternary-noise
market with informed mark $2\bar z$ and pre-order engagement share $\tfrac12$ satisfies it; a
no-disclosure structure with informed mark $\bar z$ does not, since its pooled law has four atoms,
two of which move with $\kappa$. Whether the two-round pooled cell of Section~\ref{sec:timing}
satisfies the support condition is \textbf{open}, and every result conditional on
Assumption~\ref{asm:Atau} (Lemma~\ref{lem:L3}, leg 3 of Lemma~\ref{lem:L4}, and Part~(B) of
Theorem~\ref{thm:T1}) inherits that conditionality. Remark~\ref{rem:AtauRecord} records what the implemented calibration does to the support condition.

\paragraph{The bridge clauses.}
Assumption~\ref{asm:Abr} is consumed by leg 3 of Lemma~\ref{lem:L4} and by Part~(B) of
Theorem~\ref{thm:T1}, and by nothing else. Clause (br-v) is needed in addition to (br-i)--(br-iv),
because those four clauses do not imply it. One
sharpening is recorded rather than assumed: with $\rho_{\mathrm{ch}}(\tau) := \tfrac12 A_{1/2} + A_1$, which is provably
$\kappa$-free, $\bar\pi = \bar\pi_{\mathrm{pr}}/\rho_{\mathrm{ch}}(\tau)$, so (br-iv) is equivalent to
$\rho_{\mathrm{ch}}(\tau')/\rho_{\mathrm{ch}}(\tau) \ge \bar\pi_{\mathrm{pr}}(\tau')/\bar\pi_{\mathrm{pr}}(\tau)$; under the
level-symmetric reading $\rho_{\mathrm{ch}} = \tfrac12$ and $\bar\pi = 2\bar\pi_{\mathrm{pr}}$, which forces
$\bar\pi_{\mathrm{pr}} \le 1/2$, an inherited restriction on the domain of
Assumption~\ref{asm:Atau} that Lemma~\ref{lem:L4} does not resolve. Clause (br-iii) is the clause
with the least justification behind it, and it is the one to attack first. Because (br-i) carries the
representation \eqref{eq:atau} at both thresholds, the record in Remark~\ref{rem:AtauRecord} bears
directly on Assumption~\ref{asm:Abr} as well.

%==============================================================================
\section{Numerical records}
\label{sec:app-records}
%==============================================================================

\noindent This section carries the full numerical records that Section~\ref{sec:calibration} of the main
text summarises in prose. Each record was verified on the numerical grid at the implemented
calibration; none of them moves any result, because every hypothesis they measure is carried as a
hypothesis. The window-margin record's table is Table~\ref{tab:window-record} of the main text.

\begin{remark}[Numerical record for Assumption~\ref{asm:A3}]
\label{rem:A3record}
At the implemented calibration Assumption~\ref{asm:A3} itself fails, at two independently found
loci, upstream of Assumption~\ref{asm:A6}. First, at $(\kappa = 0.5,\ \tau_{50},\ T = 5)$ with $k_2$
on an \emph{open set} above a cell edge (verified at offsets $10^{-9}$ through
$2\times10^{-2}$), the difference $U_V - U_H$ has three strict sign changes, at
$s = 1.5754434$, $1.5833333$ and $1.5902426$, with middle excursions of $2.4$--$2.8\times10^{-4}$
against a $10^{-9}$ payoff tolerance. Those excursions are the panel's grid maxima; refined
per-interval maxima are larger, at $3.07$ and $2.69\times10^{-4}$, so the failure is if anything
understated here. The pointwise argmax runs $H,V,H,V$, single-valued on each
interval, so no weakly increasing selection exists: the selection set is empty and the outer map
$\Tmap$ is \emph{undefined} there, not merely discontinuous. Second, at
$(\kappa = 0.15,\ 0.05,\ 5)$, one of the four unresolved nodes of Remark~\ref{rem:P1record},
the argmax reverses from Voice to Hold across cell edge $s = 1.659062163$ at both located fixed
points, so the preferred plan decreases in $s$. The route is the $s$-direction step of the Voice
payoff, whose interior count $n(s)$ is integer-valued, interacting with the off-path price snap. A
candidate mechanical account of the four unresolved nodes, that the trouble sits at a cell edge
where $U_H - U_V$ jumps through zero without crossing it, was probed at one node and then swept
over the other three under a pre-registered rule; it holds at all three, though in different forms.
At two of them a located fixed point sits on such an edge. At the third no fixed point was found on
any candidate edge, and the account holds instead through the achieving basin's worst deviation,
which sits in the cell immediately above one. Under the alternative, mass-proportional off-path
family the first locus's failure relocates rather than disappears: the weakly-increasing-selection
set is empty on an open set of $k_2$ under both families, so the verdict does not depend on the
family. This is diagnostic evidence at one calibration, and existence at these nodes stays neither
claimed nor denied. None of this moves any result: Assumption~\ref{asm:A3} is a hypothesis, and
Proposition~\ref{prop:P1} remains proved as a conditional.
\end{remark}

\begin{remark}[Numerical record for Assumption~\ref{asm:A6}]
\label{rem:A6record}
The continuity clause of Assumption~\ref{asm:A6} fails for the declared construction of $\Theta$,
and the locus is known. All $k$-dependence of $U_j$ runs through the pooled price vector, the
flagged layer being $k$-free under Assumption~\ref{asm:A7J}; Bayes applies where the reaching weight
is positive but a $k$-free, plan-uniform posterior is imposed on the frontier, so the price system
can be discontinuous exactly on the boundary of the reaching set of each pooled history. That set
lies inside the union of the finitely many \emph{cell-edge hyperplanes} $\{k_i = a\}$ and the
\emph{collapse faces whose dying plan is the sole generator} of some reachable pooled history,
not the collapsed cutoff vectors as such. The jump reaches $\Tmap$ with non-vanishing weight,
because $U_j$ integrates those prices against the deviator's own noise law, with weight at least
$\min(\kappa/2,\,1-\kappa)^{d+1}$, independent of the dying plan's population mass; the
vanishing-mass defusal is therefore refuted. That analytic bound is not itself computed by any
probe; what is measured is its counterpart, the jump entering the adjacent-plan payoff difference
undiminished. The largest-weakly-increasing-selection tie-break is
pointwise in $k$ and passes the jump through, and no $k$-independent perturbation family reconciles
the limits: at fixed perturbation size the system is continuous in $k$, and the discontinuity is
created only in the limit, which the choice of family cannot fix. The implemented menu's
Hold-collapse face is measured clean, with pooled prices within $4.441\times10^{-16}$ as $k_1$ sweeps
to full collapse, so the implemented menu is not in the class (middle plans owning exclusive
pooled histories) on which no family works at all; that class result is a single-pass derivation
and has not been independently checked. At the implemented calibration the failure is
live at the interior $n(s)$ cell edges instead: measured $\Tmap_2$ jumps of $6.33\times10^{-3}$,
$1.09\times10^{-2}$ and $2.83\times10^{-2}$ across steps in $k_2$ of at most $2\times10^{-9}$ at
$(\kappa = 0.5,\ \tau_{50},\ T = 5)$, measured independently by two separate scripts and agreeing to
three significant figures; at $(\kappa = 0.15,\ 0.05,\ 5)$ the jumps reach $0.16$ and a diagonal
crossing of $\Tmap_2$ is destroyed. A chamber interior $\Theta^+ = [1.23,\,1.245]\times[1.5253,\,
1.5506]$ has been exhibited that is compact, self-mapping and jump-free at the baseline node
(Brouwer runs on it verbatim and it contains $k^\star$), but it is not the $\Theta$ the existence
argument constructs, it cannot be exhibited without approximately locating the fixed point first,
and no such chamber exists at the $\kappa = 0.15$ node, where a fixed point sits exactly on the edge
$k_2 = 1.659062163$. None of this moves any result: Assumption~\ref{asm:A6} is a hypothesis, and
Proposition~\ref{prop:P1} remains proved as a conditional. Nonexistence is neither claimed nor shown:
$23$ of $27$ sweep nodes converge, and a discontinuous self-map may still have fixed points. Two
repairs are identified and neither is carried out here: a $t$-constrained game with a Kakutani
argument and $t \downarrow 0$, and a cutoff-indexed concentration family, constructible but with an
unresolved corner. What the implementation ships instead is a hard switch, so the
$\Theta$-continuity either repair would buy is absent from the implementation at every resolvable
scale; that last is a finding about the off-path family this implementation ships and not about
every family. The coverage behind all of this should be read for what it is: probes at one node per
claim class, together with the twenty-seven-node census, and no sweep over $(\kappa,\tau,T)$.
\end{remark}

\begin{remark}[Applicability, and the numerical status, of Proposition~\ref{prop:P1}]
\label{rem:P1record}
This proposition is a conditional, and its two main hypotheses about the equilibrium object are
measured to fail at the calibration the accompanying implementation runs. Assumption~\ref{asm:A6}'s
continuity clause fails there for the $\Theta$ the proof constructs (Remark~\ref{rem:A6record}), and
Assumption~\ref{asm:A3} fails there at two independently found loci (Remark~\ref{rem:A3record}). In
plain words: the proof needs the outer best-response map to be continuous and single-valued, and at
the implemented calibration that map has measured jumps and stretches where no weakly increasing
selection exists at all. The proposition stays proved as a conditional, in the same pattern as
$\Atau$: what is on record is that its antecedent, read with the $\Theta$ the proof constructs,
is not satisfied by the implemented calibration, so at that calibration the proposition asserts
nothing about the implemented model. Nonexistence is neither claimed nor shown anywhere. The
conditional stands on the proof, not on a grid. The grid evidence is stated separately: twenty-three of
twenty-seven equilibrium-sweep nodes converge, and the
remaining four ($\kappa \in \{0.15,\,0.85\}$ crossed with $(\tau,T) \in \{(0.05,\,5),\,(0.075,\,
1)\}$) stay unresolved after thirty seeds each, with best payoff-scale residuals of
$3.1\times10^{-4}$ to $1.5\times10^{-3}$ against a $10^{-9}$ criterion. Existence at those four nodes is neither claimed nor denied
by that evidence.
\end{remark}

\begin{remark}[Numerical record for Assumption~\ref{asm:Atau}]
\label{rem:AtauRecord}
At the implemented calibration Assumption~\ref{asm:Atau} \textbf{fails}. The decisive failure is
already established by the support condition alone; the derivative pattern also fails, and
independently. The pooled cell's engagement-posterior law was enumerated exactly (all
$4^{H+1} = 4{,}194{,}304$ order-flow paths, the same law the implementation integrates) at $200$
nodes: $\kappa \in \{0.05,\dots,0.95\}$ crossed with the five frozen $\tau$ percentiles and
$T \in \{1,2,5,10\}$, at frozen policies with $H = 10$. Two consistency checks establish that the
object measured is the one the restriction is about: an independent re-enumeration reproduces the pooled mass to
$0.0$ exactly, and the enumerated mean equals the pooled share $\bar\pi_{\mathrm{pr}}$ to
$1.7\times10^{-16}$. Twenty nodes are degenerate: no engaging atom survives into the pooled cell,
the law is the point mass at $0$, and the node decides nothing. At all $180$ non-degenerate nodes
$\Atau$ fails; at none does it hold. The support carries $23$ to $767$ distinct posterior values,
never three, and there is no mass at $\bar\pi/2$ at any node. Between $0.57\%$ and $91.8\%$ of the
pooled mass sits off $\{0,\bar\pi/2,\bar\pi\}$ ($13.9\%$ at the median node), and the interior
atoms move with $\kappa$: the two-sided Hausdorff distance between adjacent-$\kappa$ support sets
reaches $0.4608$ against a predicted $<10^{-12}$, at $0$ of $18$ series. The derivative restriction
$A_0' = A_1'$ holds at $0$ of $180$ nodes, and the chord identity's residual reaches $0.0717$, up
to $7.17$ premium percentage points, against a predicted $<10^{-10}$. One half of clause
($\tau$-ii) does hold: $\bar\pi = 1$ to $1.5\times10^{-13}$ at every non-degenerate node and is
$\kappa$-free to the same order, which is a separate finding and not a partial rescue. This is
numerical applicability evidence at one calibration; none of it moves any result,
because $\Atau$ is an assumption rather than a claim. Lemma~\ref{lem:L3}, leg 3 of
Lemma~\ref{lem:L4} and Part~(B) of Theorem~\ref{thm:T1} remain proved as conditionals; at this
calibration they say nothing about the implemented pooled cell. The open
question is a question about the restriction's \emph{domain}: a different menu, a different $H$, or a
different calibration could still satisfy the support condition. One coverage caveat travels with
the record: the $18$ non-degenerate series are only \emph{six distinct pooled cells}, and all six
fail.
\end{remark}

\begin{remark}[The window-margin record, verified on the numerical grid]
\label{rem:T1record}
The numerical check runs the two-round model at frozen policies, cutoffs pinned at the baseline
equilibrium $k = (1.240576,\,1.531022)$, where the flagged weight is $\Omega = 13.84$ percent, the
disclosed share of engagements is $\omega_a = 61.15$ percent, and the cell averages are
$M_F = 0.553$ and $M_P = 0.223$ premium percentage points, over a $\kappa$ grid of $71$ nodes on
$[0.15,\,0.85]$, at five $\tau$ percentiles and $T \in \{5,10\}$. Table~\ref{tab:window-record}
reports the weight leg, the composition leg and their product for the window cut
$(T',T) = (5,10)$. Attenuation holds at every node: $W_T C_T \le 1$ throughout, with no node above
one, and the composition leg carries the effect: $C_T$ runs from $0.21$ to $0.79$ while $W_T$ shaves
only $3$ to $14$ percent.


Three supporting results from the same record are worth carrying. The threshold margin gives the
same direction (zero of eight compared pairs have $W_\tau C_\tau > 1$), with the three pairs
that reclassify real mass at $0.5566$, $0.4780$ and $0.8846$ and the other five pairs reclassifying
no mass at all and returning $1.000$ to within $3\times 10^{-16}$. The factorisation of Part~(A) holds to machine noise,
which is wiring rather than evidence: at frozen policy the weights are $\kappa$-free by
construction. And Lemma~\ref{lem:L2}'s conclusion is confirmed \emph{exactly}: every flagged object
(posterior, price, entry probability, $M_F$) is invariant in $\kappa$ to machine zero over the
grid, while the pooled $M_P$ moves $0.0722$ premium percentage points and the filing-day jump $J$
moves $5{,}090$ basis points over the same range, so the target of the invariance is not vacuous.
One caveat is the record's own: at $H = 10$ the long window is the corner $T = H$, which drives
$\Omega(T{=}10)$ to $6.81\times10^{-4}$ and makes the headline comparison corner-versus-interior;
the $H = 12$ column of Table~\ref{tab:window-record}, where $T = 10$ is strictly interior, points
the same way. That column is computed by a different route, and the difference matters. At $H = 12$
the enumerated pooled-support object is computationally unavailable, so the column's composition
ratio $C_T$ is obtained from the chord closed form rather than from independently enumerated pooled
levels as at $H = 10$, and that closed form is the one whose support condition
Remark~\ref{rem:AtauRecord} measures to fail at this calibration. The $H = 12$ column is therefore
directional corroboration of the $H = 10$ comparison rather than a second independent magnitude.
What the record does not do is supply a sign theorem: it is node evidence at one
calibration, and Part~(C) claims none. At this calibration the model's directional prediction for
the February 2024 acceleration is attenuation, with the composition leg carrying it, a
prediction the superseded registered design preserved in Section~\ref{app:honesty} treats as one
of two live branches, not as a derived sign, because the antecedent behind the composition leg's closed form is the very $\Atau$
that Remark~\ref{rem:AtauRecord} measures to fail here.
\end{remark}

\begin{remark}[Dominance-and-contraction nodes, verified on the numerical grid]
\label{rem:C1record}
Eighteen of eighty grid nodes are pointwise \emph{dominance-and-contraction nodes}. The largest
contiguous block is $T = 5$ at $\tau$-percentiles $\{50,70,90\}$ with
$\kappa \in \{0.65,\,0.75,\,0.85\}$; the slack $\eta_r$ has minimum $0.0595$ and median $0.3467$, and
$L_{\mathcal R} \in [0.264,\,0.501]$ everywhere on the grid. These values come from a numerical
calculation reproduced independently. What these nodes establish is narrow and
worth stating exactly: they verify the two pointwise inequalities $L_{\mathcal R} < 1$ and
$\eta_r > 0$ together with supporting diagnostics. They do \emph{not} verify the full antecedent
(C-1)--(C-7) of Proposition~\ref{prop:C1}, and they do \emph{not} exhibit a named nonempty region; a
dominance-and-contraction node is a pointwise numerical fact, not evidence of a region. A
named-region promotion is a separate step and is not claimed. The earlier aspiration, this result
proved on a named nonempty region and numerical off it, is retired as structurally undeliverable
as worded; what is deliverable is the three objects above: the implication proved with the region as
a hypothesis, the dominance-and-contraction nodes verified on the grid, and a named-region promotion
left open, for which an $\varepsilon$-ball plus integral-control pattern is the template. The
proof also gives the bound $\mathcal B_r^{GE} = O\bigl((1-L_{\mathcal R})^{-3}\bigr)$, so
the dominance-and-contraction nodes are cubically bounded away from $L_{\mathcal R} = 1$.
\end{remark}

%==============================================================================
\section{The August 2026 empirical record, exactly as viewed}
\label{app:honesty}
%==============================================================================

\noindent This section preserves, as a record, an empirical design the paper no longer uses. The
design was registered as \path{research/empirics_v4/SPEC.md} on 2026-08-20, before any
estimation: seven pre-specified legs around the February 2024 acceleration of the Schedule 13D
filing deadline --- the run-up/jump partition test (H1), the post-reform slope change (H2), the
bindingness dose, the stake at filing, a bounded null computed from the Commission's own tables,
the hand-audited bid-arrival coding, and the matched difference-in-differences on the twelve-month
bid hazard --- plus two supporting legs, the pseudo-trigger placebo and the bidder-entry-by-liquidity
triple difference. Four dated corrigenda (2026-08-23; two of 2026-08-30; 2026-08-31) updated
counts, power inputs and the theory-lane sign as they landed, each marked ``no test changed''. The
design was executed in late August 2026 and its outcomes were viewed. The headline interpretation
the design was built for was rejected: H1's calendar-return windows are not the theory's
information cells, H2's detection thresholds are far above any economically relevant sign, and the
matched control-outcome leg is not estimated. Under the position this paper adopts, the registered
results therefore remain here, in an honesty appendix, rather than in the headline: every viewed
result and every design failure is disclosed with its intervals, minimum detectable effects (MDEs)
and gate verdicts. The record was deleted from the live tree on 2026-09-01 (deletion commit
\texttt{235de22}) and is preserved in the repository's history at commit \texttt{9b98089};
everything below is quoted from that record exactly as viewed, without recomputation and without
re-interpretation. Each leg names its registered location in the specification and the file that
carries its realised numbers; paths are those the record carried before the deletion. Throughout,
``pp'' abbreviates percentage points.

The viewed-result lock that bound the record binds it here: the outcomes may be reproduced,
tabulated and diagnosed, but not searched for significance. Nothing in this section is cited as
evidence anywhere in the main text; the main text's empirical content is the descriptive
filing-delay result of Section~\ref{sec:empirics}.

\paragraph{H1 --- the run-up/jump partition test (specification \S3.4). Label: ESTIMATED.}
Source: \path{empirics/output/h1_estimate.json} (block \path{h1_main}). The registered
estimand is the partition coefficient $\hat\beta_2$ on $\mathrm{LIQ} \times \mathrm{Flagged}$ in
the stacked run-up/jump regression: the pooled (run-up) cell should carry the liquidity derivative,
the flagged (jump) cell none, and the two slopes should differ. The realised sample is $979$
filings ($446$ pre-reform, $533$ post; $1{,}958$ stacked rows, $891$ firm clusters, $47$ month
clusters). The realised partition coefficient is $+0.58$pp with a two-way clustered standard error
of $1.56$pp; the quoted $p$-value is the conservative wild month-bootstrap value $0.741$; the
realised MDE is $4.38$pp. The pooled slope is $+0.67$pp (standard error $2.84$pp) and the flagged
slope $+1.24$pp (standard error $1.93$pp); the record's corroboration audit derives, and flags as
derived, the flagged slope's $95$\% interval, $-2.5$ to $+5.0$pp
(\path{research/empirics_v4/theory_corroboration_2026-08-31.md}). Under the registered
decision table (\S10, row~1) the leg is uninformative: a supportive reading needs the partition
coefficient significant with the pooled slope significantly non-zero and the flagged slope inside
its MDE, and the second condition fails. Robustness rows in the same file: the
$\mathrm{FD}^* \equiv \mathrm{FD}$ variant ($p = 0.725$, $N = 979$), the
one-to-thirteen-calendar-day screen ($p = 0.913$, $N = 779$), and the full-funnel variant
($p = 0.834$, $N = 1{,}075$). The within-sample standard deviation of log illiquidity is $2.926$
(interquartile range $4.303$), the reporting companion the specification attaches to every slope.
Two record-keeping facts travel with the leg, both documented in the verification note named below.
The first-run estimate ($+0.95$pp, $p = 0.586$, $N = 976$) was disclosed alongside the corrected
one under the specification's own no-change rule, with the cause named (a calendar-date bracket
where a trading-day bracket was registered). And the liquidity standardisation ran over the
$1{,}093$-row full-funnel frame rather than the estimation sample --- the largest of the seven
minor script-versus-specification deviations the verification records; re-standardising moves
$\hat\beta_2$ from $+0.58$ to $+0.48$pp, against the $4.38$pp MDE.

\paragraph{H2 --- the reform slope change (specification \S3.5). Label: ESTIMATED.}
Source: \path{empirics/output/h2_estimate.json}. The theory lane's directional sign landed on
2026-08-30 (attenuation, the specification's Branch~A; both branches remained live). The registered
contamination check runs first: a slope change on the filing-day jump of the same sign and
comparable size as on the run-up would refute the partition itself. It recorded
$\hat\delta(\mathrm{JUMP}) = +0.54$pp ($p = 0.928$) against
$\hat\delta(\mathrm{RUNUP5}) = +0.02$pp ($p = 0.997$), \path{partition_refuted: false} --- and
the record's own caption is honest about it: the check fires only on a significant jump
interaction, so with an $11.62$pp MDE against a $0.54$pp coefficient it had no power to fire; an
untaken test, not a passed one (\path{theory_corroboration_2026-08-31.md}, \S2). The leg
proper: $\hat\delta(\mathrm{RUNUP5}) = +0.02$pp, standard error $4.42$pp, quoted $p = 0.997$,
realised MDE $12.39$pp; $\hat\delta(\mathrm{JUMP}) = +0.54$pp, standard error $4.15$pp, quoted
$p = 0.928$, realised MDE $11.62$pp. Under the decision table (\S10, row~2) both are uninformative
bounded nulls: the principal coefficient is roughly one seven-hundredth of its own MDE, and the
design cannot discriminate economically relevant signs. The pre-period level prediction both
branches shared (a negative liquidity slope on the run-up, \S3.5.1) did not come through
point-wise either: the realised pre-reform slope on the five-day run-up is $+4.14$pp with standard
error $3.95$pp. The level shift $\gamma$(Post) is not identified under the year-quarter fixed
effects and was correctly not reported (\S3.5's collinearity note). Two registered defence rows:
the per-day run-up variant, $\hat\delta = -1.39$pp ($p = 0.341$, $N = 882$), and the
identical-window variant in which every run-up is exactly five trading days, $\hat\delta = +4.18$pp
($p = 0.478$, $N = 242$, MDE $14.71$pp).

\paragraph{The bindingness dose (specification \S4). Label: ESTIMATED.}
Source: \path{empirics/output/dose_estimate.json}. The dose --- a filer's pre-period propensity
to file past the five-business-day deadline, built from pre-period behaviour only --- is directly
measured for $127$ filers covering $249$ filings; the imputed-dose robustness row carries $731$,
and the specification's realised-counts corrigendum records how thin the directly measured
coverage is ($189$ of $1{,}710$ pre-period filers, $27.5$ percent of pre-period filings, on the
dose window without the script's $90$-day band). Every dose interaction is a null. On the primary
dose: run-up $\hat\varphi = +3.83$pp ($p = 0.620$, MDE $17.01$pp), jump $+3.14$pp ($p = 0.362$,
MDE $8.62$pp), stake $+0.14$pp of stake ($p = 0.966$, MDE $7.51$pp). The alternative-dose,
imputed-dose and within-liquidity-tercile rows run from $p = 0.31$ to $0.87$ with MDEs from $9.45$
to $62.36$pp. The registered confound diagnostic records $\mathrm{corr}(D, \mathrm{LIQ}) = -0.027$.
The first stage --- high-dose filers cut their realised delay by $7.25$ business days more than
low-dose filers, standard error $1.44$, $t = -5.03$, $N = 249$ --- is labelled in the artifact
itself as a mechanical validity check, not a finding. Under the decision table (\S10, row~3) the
leg is uninformative: its discriminating cells cannot fire when the H2 interaction they key on is
a zero.

\paragraph{Stake at filing (specification \S5). Label: ESTIMATED.}
Source: \path{empirics/output/stake_estimate.json}. The sample is $962$ filings with a usable
parsed stake (seven rows outside $[0,100]$ dropped; $81$ rows below five percent carried as a
count; winsorised at $[1.66, 93.53]$, winsorised standard deviation $18.57$pp; raw median $9.43$
percent against the activist-hedge-fund benchmark median of $6.3$ percent, carried as a parser
validation). The registered object is the level shift, identified by dropping the year-quarter
effects for a linear trend and quarter-of-year dummies: $\hat\gamma(\mathrm{Post}) = -0.98$pp of
stake, standard error $2.34$pp, wild-bootstrap $p = 0.714$, realised MDE $6.55$pp. The interaction
is $\hat\delta = -2.16$pp (standard error $1.27$pp, wild $p = 0.120$, realised MDE $3.57$pp), with
the year-quarter-fixed-effects variant at $-2.19$pp ($p = 0.117$) and $\gamma$ absorbed there by
construction. The liquidity main effect, $-6.65$pp of stake per standard deviation of illiquidity
(standard error $1.37$pp), is the strongest outcome coefficient in the package, and the record
notes that the model signs nothing about it. The specification's registered power warning was borne
out: the historical days-to-stake gradient implies a mechanical change of about $0.12$pp for the
two-business-day cut in the median delay, so the realised MDE is about fifty-five times what a
purely mechanical channel could produce --- the design registered in advance that it could not see
that channel, and it did not.

\paragraph{The bounded null (specification \S6). Arithmetic, not an estimate.}
Sources: \path{research/empirics_v4/SPEC.md}, \S6, re-verified against the release text in the
verification note named below; the ladder as stored in
\path{empirics/output/did_estimate.json} (block \path{bounded_null_ladder_pp}). The
Commission's Table~3 classifies non-corporate-action initial 13Ds, 2011--2021, by how much of the
reported stake was already accumulated by the amended deadline: $80$ percent of campaigns were
unconstrained, $20$ percent would have had some accumulation cut, $3$ percent would have lost at
least a tenth of the stake, $1$ percent at least a quarter. Even if every constrained campaign's
bid outcome flipped, the accumulation channel's aggregate effect on the twelve-month bid hazard is
therefore capped by a ladder: at most $20$pp at the loosest rung, at most $3$pp at the headline
rung, at most $1$pp at the tightest --- a ceiling, not an estimate. Three registered restrictions
travel with it: it covers the accumulation channel only, not the rule's aggregate footprint; a
reduced-form matched estimate may legitimately exceed it; and the dollar figures that accompany it
are the release's own, which the Commission itself disclaims. Every figure was re-grepped clean
against the release text on 2026-08-30.

\paragraph{The bid-arrival coding (specification \S8.3): gates and audit.}
Sources: \path{empirics/output/bid12_gates.json};
\path{research/empirics_v4/bid12_audit_result_2026-08-30.md}. Of the fourteen recorded coding
gates, thirteen pass; the one failure is the treated base-rate gate: the realised bid-arrival rate
is $12.61$ percent on $3{,}418$ coded rows of $3{,}710$, against the $18.1$ percent
acquired-within-twelve-months anchor borrowed from the literature for every power calculation in
\S8.6 and \S9 --- the coder counts failed and withdrawn bids, so the realised rate was expected to
run above the anchor and runs below it, and the artifact's own summary line is
\path{all_passed: false}. The caveat travels with every bid-outcome MDE the record prints: each
is built on the borrowed $18.1$ percent. The registered blind hand audit --- thirty filings,
independent fresh readers who had not written the coder, readings sealed before the keys were
opened, the control half's full-text searches re-run live --- found $0$ disagreements of $30$
against a blocking threshold above $10$ percent, and does not block the leg. Two rulebook-coverage
gaps the audit surfaced were flagged to the author with the registered rule frozen.

\begingroup\sloppy
\paragraph{The matched difference-in-differences (specification \S8). Label: NOT ESTIMATED.}
Sources: \path{empirics/output/did_estimate.json};
\path{empirics/output/did_diagnostics.json};
\path{research/empirics_v4/did_matching_2026-08-31.md}. The matched control-outcome design is
not estimated, and the artifact says so: label \texttt{NOT ESTIMATED}, status
\path{design_failure}, \path{quote_as_result: false}. The registered funnel cuts $3{,}710$
treated rows to a matched treated sample of $465$ ($325$ pre / $140$ post). The \S8.2 balance gate
fails: post-match standardised differences on the past twelve-month return ($-0.131$) and
idiosyncratic volatility ($0.122$) remain above $0.10$ after the predeclared tighter-caliper rerun
($839$ pairs at the $0.20$ caliper; $932$ at $0.25$), so no coefficient exists. The \S8.8
pre-trend test independently prohibits causal language on its own terms: joint $F = 3.25$ on six
pre-quarters, $p = 0.0214$, below the registered $0.10$ block --- nothing in this leg may be read
as an effect, whatever the balance gate had done. The \S8.7 placebo is recorded as blocked: of
$568$ registered pseudo-reform dates none was estimable, with four blockers named in the artifact.
There is no realised MDE --- no regression was fitted, and the artifact labels the power numbers
``not a realised MDE''. What exists is design arithmetic on the registered anchors: on the realised
matched counts, standard error $4.17$pp, MDE $11.70$pp, $15.32$pp clustered; the specification's
printed $9.09$pp (clustered) computes one funnel stage higher, on $569$/$543$. Against the $3$pp
headline rung of the bounded null, the design was three to five times too coarse: it could only
ever have ruled out a large effect. The registered S2 sample (corporate-action Item-4 coding) does
not exist, and no S2 row was synthesised. The survivorship treatment is recorded with its bias
directions (the control group is conditioned on survival for $701$ unresolved delisted firms, $419$
recovered and validated; the control bid rate biased down, the level coefficient biased up). One
registered bar is cleared: the hand audit above.\par\endgroup

\paragraph{Bidder entry by liquidity (specification \S9). Label: ESTIMATED, inheriting a failed
gate.} Source: \path{empirics/output/bidder_entry_estimate.json}. Within 13D targets
($N = 465$), the liquidity-by-post interaction is $-0.65$pp (standard error $3.93$pp, quoted
$p = 0.878$, realised MDE $11.01$pp), with the liquidity main effect at $-1.63$pp (standard error
$2.65$pp). The triple difference against the matched controls ($N = 1{,}133$; $326$ treated, $807$
controls) is $\hat\tau = -5.94$pp (standard error $6.59$pp, quoted $p = 0.398$, realised MDE
$18.45$pp) --- estimated on the unbalanced matched sample, so the artifact carries
\path{design_status: failed_balance} and the standing instruction that $\hat\tau$ is reported
as an interval and a sign, never as a test. The specification's printed rule-of-thumb MDEs
($8.8$pp on S1, $19.8$pp on S2) were rescaled on realised counts to $18.2$ and $40.6$pp ($15.8$pp
on the CRSP-matched variant), with the registered reading unchanged: the leg is not powered to
detect anything economically plausible, and an uninformative landing was the registered
expectation. The S2 row and the $2025$ extension are NOT ESTIMATED. Every MDE in this leg sits
downstream of the failed base-rate gate above.

\begingroup\sloppy
\paragraph{The pseudo-trigger placebo (specification \S3.7). Label: ESTIMATED.}
Source: \path{empirics/output/placebo_h1h2_estimate.json}. Trigger and filing dates are shifted
back $63$ trading days with windows identical by construction, on the full-funnel basis ($1{,}071$
rows; the basis deviation from the main-sample registration is disclosed in the record). The
pseudo-partition coefficient is $-0.19$pp (quoted $p = 0.754$, MDE $1.58$pp) and the pseudo-jump
slope change $+0.21$pp ($p = 0.727$, MDE $1.66$pp), both clean. The pseudo-run-up slope change is
$-2.01$pp ($p = 0.081$, MDE $3.22$pp): a borderline flag, recorded as one with its reasoning (the
real run-up slope change is null at $p = 0.997$), and no demotion, since the registered rule
demotes only on a significant pseudo-slope. Two registered checks produced no artifact: the 13G
descriptive placebo, pending its own fetch and parse pass and recorded in the artifact's
\path{not_run} field, and the quarterly liquidity-slope pre-trend $F$-test for the timing split,
unrun and unflagged, as the verification note records. The registered run-up path (\S7) ---
descriptive by design --- was never executed; it is listed here so that the register of registered
legs is complete.\par\endgroup

\paragraph{The independent verification of 2026-08-30.}
Source: \path{quality_reports/verification/2026-08-30_estimates_verify.md} (which survives in
the live tree). A fresh session that had written none of the scripts re-ran the four committed
estimation sets of 2026-08-30 --- the timing split, the dose, the stake, the placebo, and the
bounded-null sources --- from the manifest inputs and diffed against the committed outputs: five of
six artifacts identical up to the timestamp field, the sample file bit-identical, two spot
quantities recomputed with independent code agreeing to fifteen significant digits, and every
release figure behind the bounded null re-grepped clean. Every claimed number matched. The named
caveats: three documentation imprecisions (full-frame figures quoted where the estimation sample
was meant), seven minor script-versus-specification deviations (the largest the standardisation
population noted under H1), and the unrun pre-trend bullet noted above. The matched-DiD and
bidder-entry artifacts landed on 2026-08-31, after that verification; they carry their own gates
and audits, recorded above.

\paragraph{What the record does and does not establish.}
H1 does not establish the pooled slope: the registered supportive cell needs the run-up's liquidity
slope significantly non-zero, and it is $+0.67$pp against a $2.84$pp standard error. H2's principal
coefficient is roughly one seven-hundredth of its own minimum detectable effect; its nominal sign
is not evidence of the branch it points to. The dose is null at every row; the stake's level shift
is null at fifty-five times the mechanical magnitude it was hunting; the matched control-outcome
design is not estimated, its pre-trend test bars causal language in any case, and the bidder-entry
triple difference inherits the failed match. The most the record honestly provides is labelled
estimates, intervals, MDEs and failures --- plus one arithmetic ceiling, the Commission's Table~3
ladder, which caps the accumulation channel at about three percentage points of the twelve-month
bid hazard and is a ceiling, not an estimate. The record is preserved because it was registered,
executed, viewed and verified: a pre-specified design that is superseded belongs in the open,
disclosed in full and relied on for nothing. That is the position this paper adopts, and it is the
whole of what this section is for.

%==============================================================================
\section{Proofs of the core lemmas}
\label{sec:proofs-core}\label{sec:proofs}
%==============================================================================


\begin{proof}[Proof of Lemma~\ref{lem:D1} (D1)]
Throughout this proof, $j(\cdot)$ denotes the cutoff selection map of
Definition~\ref{def:cpbe}\,(i), the map carrying the signal into a plan at the cutoff vector in
force, and $\sigma(\cdot)$ the $\sigma$-algebra generated by its argument. Parts~(a) and~(b) are established first, under no restriction whatever on
the probability of either cell; part~(c) follows.

\emph{Step 1 (one probability space, and a Borel selection map).} By Assumption~\ref{asm:A1} the
vector $(v,\varepsilon,\xi,z_{0:H})$ has a joint law on a finite product of Polish spaces, the noise
coordinates taking values in the finite set $\{-\bar z,0,+\bar z\}$ of
Definition~\ref{def:primitives}, so every object below lives on one probability space. The signal
$s = v+\varepsilon$ is a sum of two coordinates and is Borel. By Definition~\ref{def:cpbe}\,(i),
$j(\cdot)$ is a step function with breakpoints $k_1 \le \dots \le k_{J-1}$, finitely many because the
plan menu is finite by Assumption~\ref{asm:A2p}. A step function with finitely many breakpoints is a
finite sum of indicators of half-lines, so $j(\cdot)$ is Borel and $\{s : j(s) = j\}$ is Borel for
every $j \in \mathcal J$.

\emph{Step 2 (the engagement flag is measurable).} The realised flag $a = a_{j(s)}$ composes the
Borel map of Step~1 with the map $j \mapsto a_j$ of Definition~\ref{def:plans}, defined on the finite
set $\mathcal J$. A map on a finite set carrying the discrete $\sigma$-algebra is measurable, and
$\mathcal J$ is finite by Assumption~\ref{asm:A2p}; hence $\{a=1\}$ is Borel and lies in $\sigma(s)$.

\emph{Step 3 (the crossing date is attained, unique and measurable on the Voice region).} Fix $j$
with $a_j = 1$. By clause (TR-ii) of Assumption~\ref{asm:TR}, $s \mapsto B_j(s,d)$ is weakly
increasing for Voice plans, and a weakly monotone real function is Borel because the preimage of
$[\tau,\infty)$ is an interval; so $\{s : B_j(s,d) \ge \tau\}$ is Borel at each date $d$. The
calendar $\{0,\dots,H\}$ is finite by Assumption~\ref{asm:A2p}, so
$\{c_j \le d\} = \bigcup_{d' \le d}\{s : B_j(s,d') \ge \tau\}$ is a finite union of Borel sets, and
$c_j$ is a Borel map into $\{0,\dots,H\} \cup \{+\infty\}$. Its value is unique: a nonempty subset of
the finite well-ordered set $\{0,\dots,H\}$ has exactly one minimum, and Assumption~\ref{asm:A4}
makes $c$ that first date; if the subset is empty then $c_j = +\infty$ by Definition~\ref{def:plans}.

\emph{Step 4 (the disclosure indicator is measurable).} By Definition~\ref{def:plans}, $D=1$ requires
$a=1$, so
\begin{equation}
\label{eq:D1-cellunion}
\{D = 1\} \;=\; \bigcup_{j \in \mathcal J\,:\,a_j = 1}
  \Bigl( \{s : j(s) = j\} \cap \{s : c_j(s) \le H-T\} \Bigr),
\end{equation}
using $f_j = c_j + T \le H$ if and only if $c_j \le H-T$ together with the convention
$(+\infty)+T = +\infty > H$, which is available because $H$ is finite
(Definition~\ref{def:primitives}). Each set in \eqref{eq:D1-cellunion} is Borel by Steps~1 and~3, and
the union is finite by Assumption~\ref{asm:A2p}. Hence $\{D=1\}$ is Borel and $D$ is a measurable
$\{0,1\}$-valued map. Two facts are recorded for later use: at a fixed cutoff policy every ingredient
of \eqref{eq:D1-cellunion} is a function of $s$ alone, so $\{D=1\} \in \sigma(s)$; and Borel
regularity of $s \mapsto B_j(s,d)$ for non-Voice plans was nowhere invoked, the conjunct $a=1$ having
removed those plans from the union.

\emph{Step 5 (exactly one cell, at the level of states).} Put $\mathcal C_F := \{D=1\}$ and
$\mathcal C_P := \{D=0\}$. By Step~4, $D$ is the indicator of a Borel event, so it takes exactly one
of the values $0$ and $1$ at every sample point; therefore $\mathcal C_F \cap \mathcal C_P =
\varnothing$ and $\mathcal C_F \cup \mathcal C_P$ is the whole space. This derives, rather than
asserts, the exclusivity and exhaustiveness recorded in Definition~\ref{def:prices}.

\emph{Step 6 (the assignment is a function of the control-node history, which is part~(a)).} Step~5
partitions the space of states, while part~(a) is a statement about control-node histories. By
Definition~\ref{def:prices} each pooled history $\mathcal H_d^P$ carries the coordinate ``flag landed
by $d$'', so the flag is public by construction, and the same definition supplies the content of the
control-node information set $\mathcal I_H$, which contains the pooled history up to the control
node. By Assumption~\ref{asm:A4} the filing lands exactly at $f = c+T$ and only through the
disclosure node of Section~\ref{sec:timing}, so the coordinate ``flag landed by $H$'' equals
$\ind\{f \le H\}$; combining this with $D = \ind\{a=1,\ c<\infty,\ f \le H\}$ of
Definition~\ref{def:plans} and with Assumption~\ref{asm:A4}'s restriction that only Voice plans
cross, that coordinate equals $D$. Hence $D$ is measurable with respect to the control-node history,
and the map from a history to its cell is single-valued and onto
$\{\mathcal C_F,\mathcal C_P\}$. That is part~(a). The public-flag bridge is what part~(a) turns on:
were the flag coordinate absent from the public history, $D$ would remain a measurable function of
the state by Step~4 while the cell assignment would cease to be a function of what is observed at the
control node.

\emph{Step 7 (part~(b), the forward direction).} Let $j$ be a Voice plan and suppose $f_j \le H$.
Then $c_j = f_j - T \le H-T$ and $c_j < \infty$. By Definition~\ref{def:plans} and
Assumption~\ref{asm:A4}, $c_j$ is the first date at which the path reaches $\tau$, so
$B_j(s,c_j) \ge \tau$. By clause (TR-ii) of Assumption~\ref{asm:TR}, $\partial_d B_j \ge 0$ for
Voice plans, and $c_j \le H-T$, so $B_j(s,H-T) \ge B_j(s,c_j) \ge \tau$.

\emph{Step 8 (part~(b), the converse).} Suppose $B_j(s,H-T) \ge \tau$. Since
$T \in \{1,\dots,H\}$ by Definition~\ref{def:primitives}, the date $H-T$ lies in $\{0,\dots,H-1\}$
and is a calendar date of the model. It therefore belongs to
$\{d \in \{0,\dots,H\} : B_j(s,d) \ge \tau\}$, which is thus nonempty, so its minimum satisfies
$c_j \le H-T < \infty$ by Step~3 and $f_j = c_j+T \le H$. Steps~7 and~8 together give the equivalence
of part~(b) for every Voice plan. The converse direction used no monotonicity in the calendar date;
only the forward direction did. Nor does either direction require a strict crossing: a path sitting
exactly at $\tau$ from its first hitting date through $H-T$ satisfies both sides, the whole argument
running on weak inequalities.

\emph{Step 9 (why the core carries $b_0 < \tau$).} By Definition~\ref{def:plans},
$B_j(s,-1) = b_0$ for every plan, while Hold paths are constant and Exit paths weakly decreasing in
the calendar date by clause (TR-ii) of Assumption~\ref{asm:TR}. If $b_0 \ge \tau$ then a Hold plan
has $B_j(s,0) = b_0 \ge \tau$ and hence $c_j = 0 < \infty$ at a plan with $a_j = 0$, contradicting
Assumption~\ref{asm:A4}'s restriction that only Voice plans cross. Assumption~\ref{asm:A4} and the
path table of clause (TR-ii) are therefore jointly satisfiable only under the maintained restriction
$b_0 < \tau$ of clause (TR-i), which is why a pre-existing crossing is placed outside the core model
rather than treated as a special case.

\emph{Step 10 (a flagged history determines the objects of part~(c)).} On a flagged history the
filing date $f$ is observed, being the first date at which the ``flag landed by $d$'' coordinate of
Definition~\ref{def:prices} equals one (Step~6). By Assumption~\ref{asm:A4} the filing lands exactly
at $c+T$, and $T$ is a known policy parameter, so $c = f-T$ is identified from the history and with
it the baseline date $c^- = c-1$ entering $R_d$ and $R$. By Assumption~\ref{asm:A4} the filing
truthfully reveals the stake, so the reported $B^F$ is $B_j(s,f_j)$ as defined in
Definition~\ref{def:plans}; by the timing of Section~\ref{sec:timing} the flagged order $Q^F$ is
submitted and priced and so belongs to the control-node history. Measurability of these two objects
is not inherited from finiteness and has to be argued, because clause (TR-ii) of
Assumption~\ref{asm:TR} makes stake levels continuum-valued: the finiteness of
Assumption~\ref{asm:A2p} covers the plan menu, the image of $\Gamma$, the noise support and the
calendar, not the stake level. The argument is that $f_j$ is finite-valued and measurable by
Steps~3--8, so that, decomposing over the filing date across its whole admissible range,
\begin{equation}
\label{eq:D1-BFmeas}
B_j\bigl(s,f_j(s)\bigr) \;=\; \sum_{\ell=T}^{H} \ind\{f_j(s) = \ell\}\;B_j(s,\ell)
\end{equation}
is a finite sum of measurable functions, each $B_j(\cdot,\ell)$ measurable by Step~3, and likewise
$Q_j^F = b_j^*(s) - B_j^F$.

\emph{Step 11 (every price in part~(c) is a well-defined finite real number).} For
$d \in \{0,\dots,H\}$ the pooled price $P_d^P$ is the fixed point of the pricing map at
$\mathcal H_d^P$, unique by Assumption~\ref{asm:A5}. That map is well defined only if the joint law
of the observed flow and the terminal payoff exists, and this is where clause (TR-ii) of
Assumption~\ref{asm:TR} is consumed in its full form: $X_d = q_{jd} + z_d$ with
$q_{jd} = \Gamma\bigl(B_j(s,d)-B_j(s,d-1)\bigr)$ must be a random variable for \emph{every} plan in
the support, Exit plans included, because pooled pricing integrates over every type; $\Gamma$ is
Borel because it is ordered in the increment. Borel regularity of $s \mapsto B_j(s,d)$ is automatic
for Voice and for Hold and is a genuine addition for Exit, and it is needed here and only here.
Finiteness of the prices is the local boundedness and integrability of Assumption~\ref{asm:A2p}. At
$d = -1$ the price is $P_{-1}^P = \Eop[Y]$ by the first convention of
Definition~\ref{def:prices}, finite for the same reason. The flagged price $P^F$ is the fixed point
at the flagged information set, again by Assumption~\ref{asm:A5}. On the dates: a flagged history has
$c \ge 0$, hence $c^- \ge -1$, and $f^- = c+T-1 \ge 0$ because $T \ge 1$. The convention at $d=-1$ is
not decorative. When $T = H$, Step~8 forces $B_j(s,0) \ge \tau$ and hence $c = 0$ on every flagged
history, so $c^- = -1$ there and the entire run-up path is measured from that base. One reading of
Assumption~\ref{asm:A5} is fixed here as well: since $B^F$ is continuum-valued by Step~10, the
flagged information sets are continuum-indexed, so ``unique fixed point'' is to be read as a
measurably selected family of fixed points rather than a finite list, which is what the continuity
clause of Assumption~\ref{asm:A5} supplies.

\emph{Step 12 (the identity).} By the second convention of Definition~\ref{def:prices},
$P_{\ND}(\mathcal H_{f^-}^P) = P_{f^-}^P$: the not-yet-disclosed price is the last pre-filing pooled
price at the same realised order flow. With the definitions $R = P_{f^-}^P - P_{c^-}^P$ and
$J = P^F - P_{\ND}$ of the same definition,
\begin{equation}
\label{eq:D1-telescope}
R + J \;=\; \bigl(P_{f^-}^P - P_{c^-}^P\bigr) + \bigl(P^F - P_{\ND}\bigr)
       \;=\; \bigl(P_{f^-}^P - P_{c^-}^P\bigr) + \bigl(P^F - P_{f^-}^P\bigr)
       \;=\; P^F - P_{c^-}^P ,
\end{equation}
which is \eqref{eq:runup-jump}. The cancellation of $P_{f^-}^P$ is legitimate because all three
prices are finite real numbers by Step~11. The same-order-flow reading of $P_{\ND}$ is load-bearing
and not cosmetic: under a never-disclosed counterfactual reading the middle terms would not cancel
and the identity would carry the residual $P_{f^-}^P - P_{\ND} \ne 0$.

\emph{Step 13 (the run-up path, and what the identity is).} For $d \in \{c-1,\dots,f-1\}$ the run-up
path $R_d = P_d^P - P_{c^-}^P$ is well defined by Step~11 and computable at the control node by
Step~10, with $R_{c-1} = 0$ and $R_{f-1} = R$. Together with Steps~5, 6, 8, 10 and 12 this
establishes every conjunct of the claim. Two limits on what has been shown are worth stating. The
identity \eqref{eq:D1-telescope} is definitional bookkeeping produced by the same-order-flow
definition of $P_{\ND}$, and this step does not claim any sign for $R$, for $R_d$ or for $J$, nor
$\kappa$-invariance of $J$. And no probability restriction entered Steps~4 through~9: the partition
exists even when $\Prb(D=1)$ is zero or one, so Assumption~\ref{asm:A8} is needed for positive cell
mass and never for the structural partition.
\end{proof}

\begin{proof}[Proof of Lemma~\ref{lem:L1} (L1)]
\emph{Step 1 (the kernel is bounded).} By Definition~\ref{def:premium} the engagement-premium kernel
is $h(\mathcal I) = \pi(\mathcal I)p(\mathcal I)$, and by Definition~\ref{def:prices} the engagement
posterior satisfies $\pi \in [0,1]$ while the entry probability of \eqref{eq:entry} satisfies
$p \in (0,1)$. Their product satisfies $0 \le h \le 1$ pointwise.

\emph{Step 2 (the kernel is a random variable).} The posterior
$\pi(\mathcal I_H) = \Prb(a=1 \mid \mathcal I_H)$ is a version of a conditional probability and is
therefore $\mathcal I_H$-measurable; so is $P(\mathcal I_H) = \Eop[Y \mid \mathcal I_H]$, and
Assumption~\ref{asm:A5} pins \emph{the} version up to a null set, which is what makes $h(\mathcal
I_H)$ a single well-defined random variable rather than an equivalence class from which the
decomposition would have to choose. The entry probability is the composition of these with the
continuous map displayed in \eqref{eq:entry}, hence $\mathcal I_H$-measurable, and a product of
measurable functions is measurable. With Step~1 the kernel is integrable.

\emph{Step 3 (the partition and the weight).} By Lemma~\ref{lem:D1}, $\{D=1\}$ and $\{D=0\}$ are
measurable, disjoint and cover the space, so $\Omega = \Prb(D=1)$ of Definition~\ref{def:premium} is
defined and $\Prb(D=0) = 1-\Omega$.

\emph{Step 4 (the pointwise split).} By Step~3, at every sample point exactly one of
$\ind\{D=1\}$ and $\ind\{D=0\}$ equals one and the other zero, so
$h = h\,\ind\{D=1\} + h\,\ind\{D=0\}$ pointwise.

\emph{Step 5 (integration).} Each summand in Step~4 is bounded in absolute value by $1$ by Step~1 and
is therefore integrable on the single probability space that Assumption~\ref{asm:A1} supplies.
Linearity of the integral gives
$\Eop[h] = \Eop\bigl[h\ind\{D=1\}\bigr] + \Eop\bigl[h\ind\{D=0\}\bigr]$.

\emph{Step 6 (the interior case).} Suppose $0 < \Omega < 1$. Both conditioning events then have
strictly positive probability, so the elementary conditional expectations
$\Eop[h \mid D=1] := \Eop[h\ind\{D=1\}]/\Omega$ and
$\Eop[h \mid D=0] := \Eop[h\ind\{D=0\}]/(1-\Omega)$ are defined and finite, Step~1 bounding both by
$1$. This is the defining property of conditioning on an event of positive probability, not a theorem
being leaned on. Substituting into Step~5,
$\Eop[h] = \Omega\,\Eop[h \mid D=1] + (1-\Omega)\,\Eop[h \mid D=0]$.

\emph{Step 7 (scaling to premia).} The premium wedge $\Delta_m$ is a finite strictly positive
constant, by clause (TR-i) of Assumption~\ref{asm:TR} together with the integrability of
Assumption~\ref{asm:A2p}. Multiplying Step~6 by $\Delta_m$ and reading off the definitions
$\Dact = \Delta_m\Eop[h(\mathcal I_H)]$, $M_F = \Delta_m\Eop[h \mid D=1]$ and
$M_P = \Delta_m\Eop[h \mid D=0]$ of Definition~\ref{def:premium} gives \eqref{eq:L1}.

\emph{Step 8 (the degenerate case $\Omega = 1$).} Then $\Prb(D=0) = 0$, and by Step~1
$\bigl|\Eop[h\ind\{D=0\}]\bigr| \le \Prb(D=0) = 0$, so that term of Step~5 vanishes and
$\Eop[h] = \Eop[h\ind\{D=1\}] = \Omega\,\Eop[h \mid D=1] = \Eop[h \mid D=1]$, the middle equality
being Step~6's definition, available because $\Omega = 1 > 0$. Multiplying by $\Delta_m$ gives
$\Dact = M_F$.

\emph{Step 9 (at $\Omega=1$ the pooled average is not merely unavailable but unidentified).} The
statement's clause that the null-cell average is undefined rather than imputed is a
non-identification claim, and it is proved here rather than asserted. The elementary definition of
Step~6 does not apply to $M_P$, since it would divide by $\Prb(D=0)=0$. The route through the
generated $\sigma$-algebra does not rescue it either. Let $\mathcal D := \sigma(D)$, whose atoms are
$\{D=1\}$ and $\{D=0\}$, and let $x$ be any real number. The function
\begin{equation}
\label{eq:L1-version}
\Delta_m\,\Eop[h \mid D=1]\,\ind\{D=1\} \;+\; x\,\ind\{D=0\}
\end{equation}
is $\mathcal D$-measurable, and it integrates to $\Delta_m\Eop[h\ind_A]$ over every $A \in \mathcal
D$, because two such functions with different values of $x$ differ only on the null set $\{D=0\}$.
Every $x$ therefore yields a version of $\Delta_m\Eop[h \mid \mathcal D]$, so no value of $M_P$ is
determined by the law. The correct statement at $\Omega=1$ is Step~8's one-term equation, not
\eqref{eq:L1} evaluated with a term stipulated to be zero.

\emph{Step 10 (the degenerate case $\Omega = 0$).} Mirror Steps~8 and~9 with the cells exchanged:
$\bigl|\Eop[h\ind\{D=1\}]\bigr| \le \Prb(D=1) = 0$, so $\Eop[h] = \Eop[h \mid D=0]$ and
$\Dact = M_P$, while $M_F$ is unidentified by the argument of Step~9.

Two limits on the reach of this identity are worth recording, because both are easy to overshoot.
First, the identity holds for the disclosure partition that Lemma~\ref{lem:D1} constructs and for no
other: Definition~\ref{def:premium} factors $\Omega = \Prb(a=1)\,\omega_a$, so
$\Omega \ne \Prb(a=1)$ whenever the disclosed share of engagements is below one, and averaging the
kernel over $\{a=1\}$ and $\{a=0\}$ instead gives a different and generally unequal decomposition,
the engaged but undisclosed types sitting in the pooled cell. Second, the factorisation
$\mathcal S = (1-\Omega)\mathcal S_P$ of Theorem~\ref{thm:T1}\,(A) does not follow from
\eqref{eq:L1} alone. Differentiating it in $\kappa$ at fixed policies needs three further
ingredients: $\partial_\kappa M_F = 0$, which is Lemma~\ref{lem:L2}; $\partial_\kappa\Omega = 0$,
which is derived in the proof of Lemma~\ref{lem:L2} below and which Theorem~\ref{thm:T1} carries as
its hypothesis (T-7); and $\kappa$-differentiability of $M_P$, which no standing hypothesis of
Section~\ref{sec:hypotheses} supplies and which Theorem~\ref{thm:T1} therefore carries as (T-8).
\end{proof}

\begin{proof}[Proof of Lemma~\ref{lem:L2} (L2)]
Write $\Xi := (v,s,\xi)$ for the triple whose conditional independence is at issue, $z_{0:H}$ for the
noise vector, $\mathcal H_{f^-}^P$ for the pre-filing pooled history on a flagged path, and $\Sfl$
for the flagged tuple $(B^F,Q^F,a{=}1)$ of Definition~\ref{def:prices}, that is, the filing message
$F$ augmented by the flagged order. The functions $u_1,u_2$ below are bounded and measurable and are
local to this proof.

\emph{Step 1 (where $\kappa$ lives).} At fixed cutoff and execution policies, and under the
no-feedback timing of Assumption~\ref{asm:nofeedback}, every policy object (the selected plan
$j = j(s)$, the path $B_j(s,\cdot)$, the order marks $q_{jd}(s)$, the terminal target $b_j^*(s)$, the
crossing date $c_j(s)$, the filing date $f_j(s)$, the disclosure indicator $D_j(s)$, the filing stake
$B_j^F(s)$ and the flagged order $Q_j^F(s)$) is a function of $(s;\tau,T)$ alone.
Assumption~\ref{asm:nofeedback} removes any dependence on realised order flow or prices, and fixed
policies remove any dependence on $\kappa$ through re-optimised cutoffs. Inspecting the primitives of
Definition~\ref{def:primitives}, noise-trading intensity $\kappa$ appears in exactly one place, the
law of the ternary noise mark; the laws of $v$, $\varepsilon$ and $\xi$ and the remaining constants
carry no $\kappa$. Write $Q_\kappa$ for the law of $z_{0:H}$. At fixed policies, then, $\kappa$
enters the model only through $Q_\kappa$.

\emph{Step 2 (the flagged set and its mass are $\kappa$-free).} By Lemma~\ref{lem:D1} the flagged
cell $\mathcal C_F = \{D=1\}$ is measurable, and by Step~4 of its proof it lies in $\sigma(s)$; by
Step~1 it is one fixed Borel subset of the signal line, the same subset at every $\kappa$. By
Assumption~\ref{asm:A1} and Definition~\ref{def:primitives},
$s = v+\varepsilon \sim N(\mu_v,\sigma_v^2+\sigma_\varepsilon^2)$, a law free of $\kappa$ because the
noise enters neither $v$ nor $\varepsilon$. Hence $\Omega = \Prb(s \in \mathcal C_F)$ is the same
number at every $\kappa$. This yields, as a by-product used nowhere below but consumed by
Theorem~\ref{thm:T1} as its hypothesis (T-7), that $\partial_\kappa\Omega = 0$ at fixed policies.

\emph{Step 3 (the flagged tuple pins the signal, with a measurable inverse).} On $\mathcal C_F$ the
engagement flag is $a=1$, by clause (TR-iii) of Assumption~\ref{asm:TR}, which records
$D = 1 \Rightarrow a = 1$, together with Assumption~\ref{asm:A4}'s restriction that only Voice plans
cross; the third coordinate of $\Sfl$ is therefore uninformative there. By the definition
$Q_j^F = b_j^*(s) - B_j^F$ of Definition~\ref{def:plans},
\begin{equation}
\label{eq:L2-target}
B^F_{j(s)}(s) \;+\; Q^F_{j(s)}(s) \;=\; b^*_{j(s)}(s).
\end{equation}
By Assumption~\ref{asm:A7p} in its on-path injective form A7$'$, the composed terminal target on the
right of \eqref{eq:L2-target} is strictly increasing in the signal on the flagged region and is
therefore injective there. That form is consumed almost surely on the flagged set, which is the
permissive reading and the only coherent one when $B^F$ is continuum-valued, as Step~10 of the proof
of Lemma~\ref{lem:D1} establishes it is: with a continuum-valued flagged tuple no individual tuple
carries positive probability, so on-path injectivity must be read as holding almost surely on a set
of positive probability rather than on a set of atoms. Consequently the first two coordinates of
$\Sfl$ determine $b^*_{j(s)}(s)$, which determines $s$, which at fixed policies determines
$j = j(s)$. Write $\iota_F$ for the resulting inverse. It is strictly increasing on the image of the
composed target and a monotone real function on a Borel set is Borel, so $\iota_F$ is measurable; no
appeal to a measurable-selection theorem is needed here, though injectivity together with
measurability would supply one on standard Borel spaces. Since $\Omega > 0$, the set on which this
holds carries strictly positive probability and the statement is not vacuous. Restricted to
$\mathcal C_F$, therefore, $\sigma(\Sfl) = \sigma(s)$ up to null sets. This is the step at which the
weak identification wording of Assumption~\ref{asm:A7} would not do: it permits two pairs $(j,s)$
with different pooled paths, and \eqref{eq:L2-target} then fails to be injective.

\emph{Step 4 (the pre-filing pooled history is a noise transform indexed by plan and signal).}
Define, for a plan $j$ and a signal $s$,
\begin{equation}
\label{eq:L2-upsilon}
\Upsilon_{j,s}(z_{0:H}) \;:=\;
  \Bigl( \bigl(q_{j0}(s)+z_0,\dots,q_{j,f_j(s)-1}(s)+z_{f_j(s)-1}\bigr);\ \text{flag not landed by }
  f_j(s)-1 \Bigr),
\end{equation}
so that $\mathcal H_{f^-}^P = \Upsilon_{j(s),s}(z_{0:H})$ pointwise on $\mathcal C_F$, by the pooled
history of Definition~\ref{def:prices}. The indexing by $(j,s)$ alone is exactly
Assumption~\ref{asm:nofeedback}: with within-window re-optimisation the marks would depend on
realised flow, hence on past noise, and no such indexing would exist. The map
\eqref{eq:L2-upsilon} is jointly measurable in $(s,z_{0:H})$: the filing date takes finitely many
values because the calendar is finite by Assumption~\ref{asm:A2p}, and on each of the finitely many
level sets of $f_{j(s)}(s)$ (measurable by Lemma~\ref{lem:D1}), the map is
$(s,z) \mapsto (q_{j(s)d}(s)+z_d)_{d<t}$, measurable because
$q_{jd} = \Gamma\bigl(B_j(\cdot,d)-B_j(\cdot,d-1)\bigr)$ composes a Borel function, clause (TR-ii) of
Assumption~\ref{asm:TR} giving Voice monotonicity in the signal and only Voice plans reaching
$\mathcal C_F$, with the coarsening $\Gamma$, itself Borel because it is ordered in the increment.
Patching finitely many measurable restrictions gives a measurable map. The second component of
\eqref{eq:L2-upsilon} is a deterministic function of $s$ by Step~1 and carries no information beyond
$s$.

\emph{Step 5 (observed prices enlarge nothing).} By Assumption~\ref{asm:A5} each pooled price is the
unique fixed point of the pricing map at its pooled history and is therefore a function of that
history, and likewise the flagged price is a function of the flagged information set. Adjoining the
value of a function of a $\sigma$-algebra's generator to that $\sigma$-algebra does not enlarge it,
so prices may be ignored when the information content of $\mathcal I_H$ is computed.

\emph{Step 6 (an information sandwich).} On $\mathcal C_F$,
\begin{equation}
\label{eq:L2-sandwich}
\sigma(s) \cap \mathcal C_F \;\subseteq\; \mathcal I_H \;\subseteq\;
  \sigma(s,z_{0:H}) \cap \mathcal C_F .
\end{equation}
For the left inclusion: by the timing of Section~\ref{sec:timing} the filing reveals
$F = (B^F,a{=}1)$ and the flagged order is submitted and priced, so $\Sfl$ belongs to $\mathcal I_H$,
which Definition~\ref{def:prices} sets equal to $\Sfl$ on the flagged cell; by Step~3 the tuple
recovers $s$; and $\mathcal C_F \in \sigma(s)$ by Step~2, so restricting to the flagged cell is
itself an operation within $\sigma(s)$. For the right inclusion: every component of $\mathcal I_H$
(pooled order flow, flag status, filing content, flagged order and prices) is a measurable
function of $(s,z_{0:H})$ by Steps~1, 4 and~5. Everything below uses only
\eqref{eq:L2-sandwich}, so the conclusion is robust to the exact content assigned to the
control-node information set, provided it satisfies the sandwich.

\emph{Step 7 (a freezing lemma).} Let $\mathcal G$ be a $\sigma$-algebra with $z_{0:H}$ independent
of $\mathcal G$, let $S$ be $\mathcal G$-measurable, let $w$ be bounded and jointly measurable, and
put $\bar w(x) := \int w(x,\zeta)\,Q_\kappa(\mathrm d\zeta)$. Then
$\Eop[w(S,z_{0:H}) \mid \mathcal G] = \bar w(S)$. For a product form $w(x,\zeta) = u_1(x)u_2(\zeta)$
this is
$\Eop[u_1(S)u_2(z_{0:H}) \mid \mathcal G] = u_1(S)\Eop[u_2(z_{0:H})] = \bar w(S)$, pulling out the
$\mathcal G$-measurable factor and using independence. The product forms generate the product
$\sigma$-algebra and are closed under multiplication; both sides of the asserted equality are linear
in $w$ and stable under bounded monotone limits, so the monotone class theorem extends the equality
to every bounded jointly measurable $w$.

\emph{Step 8 (the conditional independence).} By Step~3, conditioning on $\sigma(\Sfl)$ on the
flagged set is conditioning on $\sigma(s)$, so what has to be shown is
$\mathcal H_{f^-}^P \perp \Xi \mid s$ on $\mathcal C_F$. Given $s$, the middle coordinate of $\Xi$ is
degenerate, so it suffices to show, for bounded measurable $u_1$ on the pre-filing history space and
bounded measurable $u_2$ on $\mathbb R^2$,
\begin{equation}
\label{eq:L2-factor}
\Eop\bigl[u_1(\mathcal H_{f^-}^P)\,u_2(v,\xi) \mid s\bigr]
 \;=\; \Eop\bigl[u_1(\mathcal H_{f^-}^P) \mid s\bigr]\cdot\Eop\bigl[u_2(v,\xi) \mid s\bigr]
 \qquad \text{a.s.\ on } \mathcal C_F,
\end{equation}
which is the standard characterisation of conditional independence given a $\sigma$-algebra. Put
$w(x,\zeta) := u_1\bigl(\Upsilon_{j(x),x}(\zeta)\bigr)$, bounded and jointly measurable by Step~4, so
that $u_1(\mathcal H_{f^-}^P) = w(s,z_{0:H})$. By Assumption~\ref{asm:A1} the noise vector is
independent of $(v,\varepsilon,\xi)$ and hence of $\Xi$, which is a measurable function of that
triple. Apply Step~7 with $\mathcal G = \sigma(v,s,\xi)$ and $S = s$:
$\Eop[w(s,z_{0:H})u_2(v,\xi) \mid v,s,\xi] = u_2(v,\xi)\,\bar w(s)$. Taking
$\Eop[\,\cdot \mid s]$ of both sides and pulling out the $\sigma(s)$-measurable factor $\bar w(s)$
gives $\Eop[u_1u_2 \mid s] = \bar w(s)\,\Eop[u_2 \mid s]$; setting $u_2 \equiv 1$ in the same display
gives $\Eop[u_1 \mid s] = \bar w(s)$, and substituting yields \eqref{eq:L2-factor}. This is the first
conjunct of the statement. The apparent paradox dissolves at exactly this point: unconditionally the
pre-filing pooled history depends strongly on the signal through the order marks, and the
conditioning event removes that dependence by pinning the signal, leaving the noise as the only
remaining randomness in \eqref{eq:L2-upsilon}.

\emph{Step 9 (the flagged posterior).} Let $u_3$ be bounded, $\mathcal I_H$-measurable and supported
on $\mathcal C_F$. By the right inclusion of \eqref{eq:L2-sandwich} and the Doob--Dynkin lemma,
$u_3 = u_5(s,z_{0:H})$ for some bounded jointly measurable $u_5$. Running the computation of Step~8
with $u_5$ in place of $w$ gives $\Eop[u_2u_3] = \Eop\bigl[\bar u_5(s)\,\Eop[u_2 \mid s]\bigr]$,
while the tower property together with $\Eop[u_3 \mid s] = \bar u_5(s)$ from Step~7 gives
$\Eop\bigl[u_3\,\Eop[u_2 \mid s]\bigr] = \Eop\bigl[\bar u_5(s)\,\Eop[u_2 \mid s]\bigr]$. Hence
$\Eop[u_2u_3] = \Eop\bigl[\Eop[u_2 \mid s]\,u_3\bigr]$ for every such $u_3$; since
$\Eop[u_2 \mid s]$ is $\sigma(s)$-measurable and the left inclusion of \eqref{eq:L2-sandwich} places
$\sigma(s) \cap \mathcal C_F$ inside $\mathcal I_H$, it is a version of $\Eop[u_2 \mid \mathcal I_H]$
on the flagged set:
\begin{equation}
\label{eq:L2-posterior}
\Eop[u_2(v,\xi) \mid \mathcal I_H] \;=\; \Eop[u_2(v,\xi) \mid s]
 \qquad \text{a.s.\ on } \mathcal C_F .
\end{equation}
The extension from bounded to integrable $u_2$, needed at Step~11 for $u_2(v,\xi) = v$, follows by
applying \eqref{eq:L2-posterior} to the truncations $\max(-n,\min(n,u_2))$ and passing to the limit
under conditional dominated convergence, the dominating variable being $|u_2|$, integrable because
$v$ is Gaussian by Definition~\ref{def:primitives}. Separately, $\pi(\mathcal I_H) = 1$ on
$\mathcal C_F$, because the flagged cell lies inside $\{a=1\}$ by clause (TR-iii) of
Assumption~\ref{asm:TR} and by Assumption~\ref{asm:A4}'s truthful revelation of purpose, and because
the cell is itself $\mathcal I_H$-measurable by Step~6 of the proof of Lemma~\ref{lem:D1}.

\emph{Step 10 (the posterior is $\kappa$-free).} By \eqref{eq:L2-posterior} the conditional law of
$(v,\xi)$ given $\mathcal I_H$ on the flagged set equals its conditional law given $s$. By
Assumption~\ref{asm:A1} and Definition~\ref{def:primitives} that law is
$v \mid s \sim N\bigl(\mu_v+\beta(s-\mu_v),(1-\beta)\sigma_v^2\bigr)$ with
$\beta = \sigma_v^2/(\sigma_v^2+\sigma_\varepsilon^2)$, and $\xi \mid s \sim N(0,\sigma_\xi^2)$
independent of $v$. None of these depends on $\kappa$, which by Step~1 lives only in $Q_\kappa$ and
enters no coordinate of $\Xi$. With $\pi = 1$ from Step~9, the flagged posterior over $(v,a,\xi)$ is
a $\kappa$-free function of the signal, hence by Step~3 a $\kappa$-free function of $\Sfl$.

\emph{Step 11 (the flagged price).} By clause (TR-iv) of Assumption~\ref{asm:TR} the flagged price
solves the inner fixed point $P = \Eop[Y \mid \mathcal I_H]$ with $Y$ as in \eqref{eq:Y}. On the
flagged cell $a=1$ and $\pi=1$, and $m_0 + \Delta_m = m_1$ by Definition~\ref{def:primitives}. Only
two properties of the bidder-entry rule are used, and they are the two the statement carries as
bookkeeping: (i) $\Prb(\mathsf B = 1 \mid \mathcal I_H) = p(\mathcal I_H)$, and (ii)
$\mathsf B \perp v \mid \mathcal I_H$. The rule \eqref{eq:entry} has both, since under
Assumption~\ref{asm:A1} and Step~10 the bidder's shock satisfies
$\xi \mid \mathcal I_H \sim N(0,\sigma_\xi^2)$ independently of $v$ and the entry indicator is a
function of $\xi$ and of $\mathcal I_H$-measurable quantities; any other rule with these two
properties works as well. Under them,
\begin{equation}
\label{eq:L2-innerFP}
P \;=\; \bigl(1-p(P)\bigr)\Bigl(\mu_v+\beta(s-\mu_v)+\Delta_V\Bigr) \;+\; p(P)\bigl(P+m_1\bigr),
\qquad
p(P) \;=\; 1-\Phi\!\left(\frac{P+K+m_1-\bar S}{\sigma_\xi}\right),
\end{equation}
where $\Eop[v \mid \mathcal I_H] = \Eop[v \mid s] = \mu_v+\beta(s-\mu_v)$ by
\eqref{eq:L2-posterior}. The data of \eqref{eq:L2-innerFP} are the single number $\Eop[v \mid s]$ and
the parameters of Definition~\ref{def:primitives}, none of which depends on $\kappa$ by Step~10.
Assumption~\ref{asm:A5} makes the fixed point unique, so the solution is a function of the data
alone and equal data give equal solutions; uniqueness is load-bearing rather than decorative here,
since without it an extraneous fixed-point selection could vary with $\kappa$ even though the map
does not. Hence the flagged price is a $\kappa$-free function of the signal and, by Step~3, of
$\Sfl$.

\emph{Step 12 (the entry probability).} By \eqref{eq:entry} with $\pi = 1$ from Step~9, the flagged
entry probability is $p = 1-\Phi\bigl((P^F+K+m_1-\bar S)/\sigma_\xi\bigr)$ on the flagged cell. Every
argument is $\kappa$-free, by Step~11 for the price and by Definition~\ref{def:primitives} for the
constants. Write $p_F(s)$ for it; it is measurable as a composition of the monotone map
$s \mapsto \Eop[v \mid s]$, the fixed-point solution and the normal distribution function, and it
takes values in $(0,1)$.

\emph{Step 13 (the flagged cell average).} By Definition~\ref{def:premium} the kernel is
$h = \pi p$, so $h = p_F(s)$ on the flagged cell by Steps~9 and~12. Since $\Omega > 0$, Step~6 of the
proof of Lemma~\ref{lem:L1} gives
\begin{equation}
\label{eq:L2-MF}
M_F \;=\; \Delta_m\,\Eop[h \mid D=1]
     \;=\; \frac{\Delta_m\,\Eop\bigl[p_F(s)\,\ind\{s \in \mathcal C_F\}\bigr]}{\Omega} .
\end{equation}
The numerator of \eqref{eq:L2-MF} integrates a $\kappa$-free bounded measurable function (Step~12)
over a $\kappa$-free Borel set (Step~2) against the $\kappa$-free law of the signal (Step~2), and the
denominator is $\kappa$-free and strictly positive. Hence $M_F$ is invariant to $\kappa$, which
completes the four invariance conclusions.

Two boundaries of the result are part of it. Nothing above asserts $\kappa$-invariance of the
pre-filing pooled price: that price conditions on order flow whose law is $Q_\kappa$-dependent and
moves with $\kappa$ in general, so with $J = P^F - P_{\ND}$ and $\partial_\kappa P^F = 0$ on the
flagged set, wherever the derivative exists $\partial_\kappa J = -\partial_\kappa P_{f^-}^P$. That
reproduces Definition~\ref{def:prices}'s refusal to claim $\kappa$-invariance of the filing-day jump.
And the result is a fixed-policy statement: if cutoffs or execution paths move with $\kappa$, Step~1
is unavailable and the resulting policy and composition responses are general-equilibrium channels,
which is the territory of Proposition~\ref{prop:C1} rather than of this lemma.
\end{proof}

\begin{proof}[Proof of Lemma~\ref{lem:L3} (L3)]
The order taken here is part~(b) first, since it is a statement about one fixed $\bar\pi$ and
consumes none of the chord restriction, then part~(a), then part~(c), which combines them.
Throughout, $\bar\pi > 0$ and $\kappa$ ranges over an open interval on which the weights of
\eqref{eq:atau} are differentiable.

\emph{Step 1 (what the minimal regularity gives).} By hypothesis $h$ is continuous on $[0,\bar\pi]$
and twice differentiable on the open interval $(0,\bar\pi)$. Twice differentiable on $(0,\bar\pi)$
means that $h'$ exists at every point of that interval and is itself differentiable at every point of
it; a differentiable function is continuous, so $h'$ is continuous on $(0,\bar\pi)$. Nothing is
assumed about $h'$ or $h''$ at either endpoint, and no continuity of $h''$ is assumed anywhere.

\emph{Step 2 (a first difference, and the second difference it produces).} For
$t \in [0,\bar\pi/2]$ set $\delta_h(t) := h\bigl(t+\tfrac{\bar\pi}{2}\bigr) - h(t)$. Evaluating at the
two endpoints of that interval and subtracting,
\begin{equation}
\label{eq:L3-seconddiff}
\delta_h\!\left(\tfrac{\bar\pi}{2}\right) - \delta_h(0)
 = \Bigl[h(\bar\pi)-h\!\left(\tfrac{\bar\pi}{2}\right)\Bigr]
 - \Bigl[h\!\left(\tfrac{\bar\pi}{2}\right)-h(0)\Bigr]
 = h(0) - 2h\!\left(\tfrac{\bar\pi}{2}\right) + h(\bar\pi) = C_h(\bar\pi),
\end{equation}
the last equality being the definition \eqref{eq:chord}. All four evaluations are at points of
$[0,\bar\pi]$, where $h$ is defined by Step~1, so each term is a real number and the cancellation is
arithmetic: the value $h(\bar\pi/2)$ appears once with a minus sign from each bracket, leaving the
coefficient $-2$.

\emph{Step 3 (the regularity of that first difference).} The function $\delta_h$ is continuous on
$[0,\bar\pi/2]$, since for $t$ there both $t$ and $t+\bar\pi/2$ lie in $[0,\bar\pi]$, where $h$ is
continuous. It is differentiable on the open interval $(0,\bar\pi/2)$, since for $t$ there both
$t \in (0,\bar\pi/2) \subset (0,\bar\pi)$ and
$t+\bar\pi/2 \in (\bar\pi/2,\bar\pi) \subset (0,\bar\pi)$ are points at which $h'$ exists by Step~1,
and the derivative of a difference is the difference of derivatives, so
$\delta_h'(t) = h'\bigl(t+\bar\pi/2\bigr) - h'(t)$.

\emph{Step 4 (the first application of the mean value theorem).} Step~3 supplies both hypotheses of
the mean value theorem for $\delta_h$ on $[0,\bar\pi/2]$, so there is $t_1 \in (0,\bar\pi/2)$ with
$\delta_h(\bar\pi/2) - \delta_h(0)
 = \tfrac{\bar\pi}{2}\,\delta_h'(t_1)
 = \tfrac{\bar\pi}{2}\bigl[h'(t_1+\tfrac{\bar\pi}{2}) - h'(t_1)\bigr]$.

\emph{Step 5 (the second application).} Consider the closed interval
$[\,t_1,\ t_1+\bar\pi/2\,]$. Since $0 < t_1 < \bar\pi/2$ by Step~4, its left endpoint satisfies
$t_1 > 0$ and its right endpoint satisfies $t_1+\bar\pi/2 < \bar\pi$, so the interval is contained in
$(0,\bar\pi)$, where $h'$ is continuous and differentiable by Step~1. The mean value theorem applied
to $h'$ on that interval yields $\zeta \in (t_1,\ t_1+\bar\pi/2)$ with
$h'(t_1+\tfrac{\bar\pi}{2}) - h'(t_1) = \tfrac{\bar\pi}{2}h''(\zeta)$.

\emph{Step 6 (part~(b)).} Chaining \eqref{eq:L3-seconddiff} with Steps~4 and~5,
\begin{equation}
\label{eq:L3-mvt}
C_h(\bar\pi) \;=\; \tfrac{\bar\pi}{2}\cdot\tfrac{\bar\pi}{2}\,h''(\zeta)
 \;=\; \tfrac14\,\bar\pi^2\,h''(\zeta),
\qquad \zeta \in \bigl(t_1,\ t_1+\tfrac{\bar\pi}{2}\bigr) \subset (0,\bar\pi),
\end{equation}
which is part~(b). Two features of this derivation are worth recording, because the corollary at
Step~12 does not share them. The identity is exact: no term was discarded and there is no remainder.
And every use of $h$ beyond continuity was at points of the \emph{open} interval $(0,\bar\pi)$, so
neither $h'$ nor $h''$ was invoked at $0$ or at $\bar\pi$; the minimal regularity of Step~1 is
genuinely all that is consumed, and differentiability at the endpoints is not. The statement's gloss
that Darboux's theorem does the rest names the alternative route to \eqref{eq:L3-mvt}, which reaches
the same point from a symmetric two-point average: $h''$ is a derivative on $(0,\bar\pi)$ and
therefore has the intermediate value property, so the arithmetic mean of its values at two points of
a compact subinterval is itself a value of $h''$ at some point between them. Neither route uses
continuity of $h''$, which is why the hypothesis does not carry it.

\emph{Step 7 (differentiating the representation).} Fix $\kappa$. By Assumption~\ref{asm:Atau} the
pooled block's expectation of the kernel is the three-term sum \eqref{eq:atau}. Clause ($\tau$-i) of
that assumption makes the kernel a function of the engagement posterior alone, so the three numbers
$h(0)$, $h(\bar\pi/2)$ and $h(\bar\pi)$ are well defined; clause ($\tau$-ii) fixes the three
evaluation points and $\bar\pi$ itself free of $\kappa$, so those three numbers do not vary with
$\kappa$. A finite sum of products of a constant with a differentiable function of $\kappa$ is
differentiable, with derivative the corresponding sum, so
\begin{equation}
\label{eq:L3-diff}
\partial_\kappa \Eop_\kappa[h] \;=\; A_0'(\kappa)h(0)
 + A_{1/2}'(\kappa)h\!\left(\tfrac{\bar\pi}{2}\right) + A_1'(\kappa)h(\bar\pi).
\end{equation}
Both clauses are load-bearing at exactly this step and neither is free. Without ($\tau$-i) the
right-hand side of \eqref{eq:L3-diff} carries the extra term
$\sum_i A_i(\kappa)\,\partial_\kappa h(\pi_i)$ and part~(a) is false; in the model
$h = \pi\,p$ is a function of two scalars, since \eqref{eq:entry} makes entry depend on the price as
well as on the posterior, so the clause is a restriction and not a reading. Without the $\kappa$-free
$\bar\pi$ half of ($\tau$-ii) it carries the further term
$\bigl[\tfrac12 A_{1/2}h'(\bar\pi/2) + A_1 h'(\bar\pi)\bigr]\partial_\kappa\bar\pi$, which is first
order in $\bar\pi$ against the quadratic vanishing of part~(c), so part~(c) is false without that
half of the clause.

\emph{Step 8 (part~(a)).} Substitute Assumption~\ref{asm:Atau}'s restrictions
$A_0' = A_1' = \Akap$ and $A_{1/2}' = -2\Akap$ into \eqref{eq:L3-diff} and factor out the common
coefficient:
\begin{equation}
\label{eq:L3-parta}
\partial_\kappa \Eop_\kappa[h]
 \;=\; \Akap\Bigl[h(0) - 2h\!\left(\tfrac{\bar\pi}{2}\right) + h(\bar\pi)\Bigr]
 \;=\; \Akap\,C_h(\bar\pi),
\end{equation}
the last equality being \eqref{eq:chord}. This is part~(a), and it is exact. The three restrictions
on the weight derivatives are used at this factorisation and nowhere else: they are precisely the
coefficient pattern $(+1,-2,+1)$ that the second difference carries, which is why the constant of
proportionality is a scalar rather than a triple.

\emph{Step 9 (the same conclusion through the chord gap, and where the symmetry is used).} Let
$\ell_h$ be the affine function on $[0,\bar\pi]$ with $\ell_h(0) = h(0)$ and
$\ell_h(\bar\pi) = h(\bar\pi)$; it is well defined because $\bar\pi > 0$. Since $\ell_h$ and
$h-\ell_h$ sum to $h$ pointwise, $\Eop_\kappa[h] = \Eop_\kappa[\ell_h] + \Eop_\kappa[h-\ell_h]$. The
second term collapses to a single product: $h-\ell_h$ vanishes at $0$ and at $\bar\pi$ by
construction, so under the three-point support of Assumption~\ref{asm:Atau} only the middle atom
survives, and affinity gives $\ell_h(\bar\pi/2) = \tfrac12\bigl(h(0)+h(\bar\pi)\bigr)$, whence
\begin{equation}
\label{eq:L3-gap}
\Eop_\kappa[h-\ell_h]
 \;=\; A_{1/2}(\kappa)\Bigl[h\!\left(\tfrac{\bar\pi}{2}\right)
   - \tfrac12\bigl(h(0)+h(\bar\pi)\bigr)\Bigr]
 \;=\; -\tfrac12\,A_{1/2}(\kappa)\,C_h(\bar\pi).
\end{equation}
Differentiating \eqref{eq:L3-gap} and using $A_{1/2}' = -2\Akap$ returns
$\partial_\kappa\Eop_\kappa[h-\ell_h] = \Akap C_h(\bar\pi)$. The first term contributes nothing:
writing the affine function through its intercept and slope,
$\Eop_\kappa[\ell_h] = \ell_h(0)\,\mathrm{m}(\kappa) + \bigl(h(\bar\pi)-h(0)\bigr)\bar\pi^{-1}
\mathrm{r}(\kappa)$, where $\mathrm{m}$ is the pooled block's total mass and $\mathrm{r}$ its
unnormalised engagement moment, and both are $\kappa$-free by hypothesis, so the derivative vanishes.
The hypothesis $h(0)=0$ makes $\ell_h(0) = 0$ and kills the mass term outright; the moment term needs
its own invariance in either case. This route agrees with Step~8 and displays the mechanism: the
affine part of the kernel contributes no interior motion at all, and the whole motion is carried by
the mass of the single middle atom. It also locates where the three-point symmetry is used, since it
is what collapses the chord-gap sum to one term. For a support with more than one interior atom the
same argument gives $\partial_\kappa\Eop_\kappa[h] = \sum_i A_i'(\kappa)\,(h-\ell_h)(\pi_i)$, a
weighted sum of chord gaps at distinct interior points, which is not a scalar multiple of the second
difference at the midpoint. The two routes are not independent derivations and must not be counted as
two: Step~8 consumes the three weight restrictions, this step consumes one of them together with the
two conservation laws, and Step~15 below shows those input sets are equivalent.

\emph{Step 10 (which normalisation, and the bridge to the pooled sensitivity).} By
Lemma~\ref{lem:D1} the pooled cell is measurable and the two cells are exclusive and exhaustive, so
the pooled block is a genuine cell of the partition and the weights of \eqref{eq:atau} are its atom
masses under either normalisation, conditional on $\{D=0\}$ and summing to one, or unnormalised and
summing to $1-\Omega$. Under the conditional normalisation $\Eop_\kappa[h] = \Eop[h \mid D=0]$, so
Definition~\ref{def:premium} and \eqref{eq:L3-parta} give
\begin{equation}
\label{eq:L3-SP}
\partial_\kappa M_P \;=\; \Delta_m\,\Akap\,C_h(\bar\pi),
\qquad
\mathcal S_P \;=\; \Delta_m\,\lvert \Akap\rvert\,\lvert C_h(\bar\pi)\rvert .
\end{equation}
Under the unnormalised normalisation the weights are $(1-\Omega)$ times the conditional ones, and
$1-\Omega$ is the pooled block's total mass, $\kappa$-free by hypothesis, so the two versions of
\eqref{eq:L3-parta} differ by a $\kappa$-free factor and the identity holds verbatim in both.

\emph{Step 11 (combining parts (a) and (b)).} Applying Step~6 to the kernel and substituting
\eqref{eq:L3-mvt} into \eqref{eq:L3-parta},
\begin{equation}
\label{eq:L3-exact}
\partial_\kappa\Eop_\kappa[h] \;=\; \tfrac14\,\Akap\,\bar\pi^2\,h''(\zeta_{\bar\pi}),
\qquad \zeta_{\bar\pi} \in (0,\bar\pi),
\end{equation}
exactly, at every admissible $\bar\pi$ and $\kappa$. No limit has been taken and no term discarded.

\emph{Step 12 (part~(c), the small-$\bar\pi$ corollary).} Assume in addition the second-order Peano
expansion of the kernel at $0+$, $h(\pi) = h(0)+h'(0)\pi+\tfrac12h''(0)\pi^2+o(\pi^2)$, and that one
and the same kernel serves the whole shrinking family. Substituting the expansion at the three
evaluation points of \eqref{eq:chord}, the constant terms cancel because $1-2+1 = 0$, the linear
terms cancel because $-2\cdot\tfrac12+1 = 0$, the quadratic terms leave
$\bigl(-2\cdot\tfrac18+\tfrac12\bigr)h''(0)\bar\pi^2 = \tfrac14h''(0)\bar\pi^2$, and the three
remainder terms are absorbed into a single $o(\bar\pi^2)$. Hence
$C_h(\bar\pi) = \tfrac14h''(0)\bar\pi^2 + o(\bar\pi^2)$ as $\bar\pi \downarrow 0$. The second clause
is not a technicality: without one kernel along the family, the three evaluations being compared at
different $\bar\pi$ belong to different functions and no limit statement is available at all. A
different sufficient condition, which does not run through Step~6, is that $h''$ extend continuously
to $0$ from the right, since then \eqref{eq:L3-mvt} gives
$\bigl|C_h(\bar\pi)-\tfrac14h''(0)\bar\pi^2\bigr|
= \tfrac14\bar\pi^2\bigl|h''(\zeta_{\bar\pi})-h''(0)\bigr|$ with
$\zeta_{\bar\pi} \to 0$. That route assumes less about the kernel near $0$ and more at $0$; neither
route implies the other, and neither is needed for part~(b).

\emph{Step 13 (the rate at which the interior motion vanishes).} Part~(c) asserts the vanishing of
the interior motion at rate $\bar\pi^2$, and the motion is the product \eqref{eq:L3-parta} rather
than the chord alone. Along the shrinking family,
$\bigl|\partial_\kappa\Eop_\kappa[h]\bigr| = \lvert\Akap\rvert\,\lvert C_h(\bar\pi)\rvert$, so the
rate transfers from the chord to the motion only if $\lvert\Akap\rvert$ is bounded \emph{uniformly in
$\bar\pi$} along the limit, which is the clause the statement carries for the seam at which
Lemma~\ref{lem:L4} consumes this lemma. Boundedness of $\lvert\Akap\rvert$ on $[0,1]$ at a fixed
policy is boundedness in $\kappa$ at one margin and does not deliver this. Under the uniform bound,
\begin{equation}
\label{eq:L3-rate}
\bigl|\partial_\kappa\Eop_\kappa[h]\bigr|
 \;\le\; \tfrac14\,\Bigl(\sup_{\bar\pi}\lvert\Akap\rvert\Bigr)\,\bar\pi^2\,
   \bigl|h''(0)+o(1)\bigr| \;=\; O(\bar\pi^2)
 \qquad (\bar\pi \downarrow 0),
\end{equation}
which is part~(c). The vanishing itself is more robust than the quadratic rate: by
\eqref{eq:L3-parta} the motion is $\Akap C_h(\bar\pi)$, and continuity of $h$ at $0$ alone gives
$C_h(\bar\pi) \to h(0)-2h(0)+h(0) = 0$, so the motion vanishes under continuity at $0$, the uniform
bound and the one-kernel clause, with the two remaining clauses of Step~12 buying the rate rather
than the vanishing.

\emph{Step 14 (the vanishing chord, and why the statement is an ``if'').} Suppose
$C_h(\bar\pi) = 0$. Then \eqref{eq:L3-parta} gives $\partial_\kappa\Eop_\kappa[h] = \Akap\cdot 0 = 0$
at every $\kappa$ in the interval, so the pooled block's expectation is constant there and, by
\eqref{eq:L3-SP}, $\partial_\kappa M_P = 0$ and $\mathcal S_P = 0$. The interior motion is exactly
zero, not small and not merely of indefinite sign. Three consequences belong to the statement rather
than to this case. First, the implication runs one way only. Zero interior motion does not imply a
vanishing chord: take any pooled law satisfying \eqref{eq:atau} whose weights happen to be locally
constant in $\kappa$ on a subinterval, where $\Akap = 0$ and \eqref{eq:L3-parta} returns zero motion
for \emph{every} kernel, including kernels with $C_h(\bar\pi)$ strictly negative. Writing the lemma
with ``if and only if'' would therefore assert something false, which is why it is stated as an
implication and never as an equivalence. Second, a vanishing chord does not make the kernel affine:
by \eqref{eq:L3-mvt} it forces $h''(\zeta_{\bar\pi}) = 0$ at the one point the mean value theorem
produced and at no other, and a kernel strictly concave on part of $[0,\bar\pi]$ and strictly convex
on the rest can have $C_h(\bar\pi) = 0$ exactly. Third, under the hypothesis $h(0)=0$ the vanishing
chord is the testable identity $h(\bar\pi) = 2h(\bar\pi/2)$: the top of the chord is exactly twice
the midpoint.

\emph{Step 15 (what the hypothesis set does, recorded rather than claimed).} The hypotheses of this
lemma are not overdetermined, and it is worth saying why, since two of them look like separate
restrictions and are one. Suppose the pooled block's posterior law is supported on the
$\kappa$-invariant points $\{0,\bar\pi/2,\bar\pi\}$ with weights differentiable in $\kappa$, summing
to a total mass $\mathrm{m}(\kappa)$ and carrying the unnormalised engagement moment
$\mathrm{r}(\kappa) = A_{1/2}(\kappa)\tfrac{\bar\pi}{2}+A_1(\kappa)\bar\pi$. Then
\begin{equation}
\label{eq:L3-equiv}
\bigl[\,\mathrm{m}'(\kappa) = 0 \ \text{ and }\ \mathrm{r}'(\kappa) = 0\,\bigr]
\qquad \Longleftrightarrow \qquad
\bigl[\,A_0' = A_1' \ \text{ and }\ A_{1/2}' = -2A_1'\,\bigr].
\end{equation}
In one direction, $\mathrm{r}' = 0$ reads $A_{1/2}'\tfrac{\bar\pi}{2}+A_1'\bar\pi = 0$, and dividing
by $\bar\pi/2 > 0$ gives $A_{1/2}' = -2A_1'$; substituting that into
$\mathrm{m}' = A_0'+A_{1/2}'+A_1' = 0$ gives $A_0' = 2A_1'-A_1' = A_1'$. In the other, summing the
two restrictions gives $A_1'-2A_1'+A_1' = 0$, so $\mathrm{m}' = 0$, and
$-2A_1'\tfrac{\bar\pi}{2}+A_1'\bar\pi = 0$, so $\mathrm{r}' = 0$. Given the $\kappa$-invariant
three-point support, therefore, the derivative pattern of Assumption~\ref{asm:Atau} and the two
conservation laws assumed here are the same restriction stated twice, and the entire remaining
content of the assumption is its support condition. One further consequence of \eqref{eq:atau} is
recorded because Lemma~\ref{lem:L4} consumes it: reading $\bar\pi$ as the \emph{mean} of the pooled
law rather than as its upper support point is degenerate. With conditional weights summing to one and
mean equal to $\bar\pi$, the moment equation is $\tfrac12 A_{1/2}+A_1 = 1$, which with
$A_0+A_{1/2}+A_1 = 1$ gives $A_0 = A_1-1 \le 0$; nonnegativity then forces $A_0 = 0$, $A_1 = 1$ and
$A_{1/2} = 0$, a point mass at $\bar\pi$ with $\Akap = 0$ and, by \eqref{eq:L3-parta}, zero interior
motion for every kernel. A mean cannot equal the maximum of its own support unless the law is
degenerate, which is why Definition~\ref{def:premium} reads $\bar\pi$ as the upper support point and
excludes the other reading throughout.

The three steps that follow ask what the support condition excludes and what would suffice to
establish it. They are recorded as part of the lemma because they fix the assumption's domain, on
which every conditional statement downstream of it depends; none of them is used in the derivation
of parts~(a)--(c) above.

\emph{Step 16 (a structure that satisfies the support condition, built from the model's own
primitives).} Take a one-round market with a single trading date, engagement binary with
$\Prb(a=1) = \varpi = \tfrac12$, the engaging plan's pooled order mark $q_0 = 2\bar z$, and the
non-engaging plan's mark $q_0 = 0$. The engaging mark is admissible because
Definition~\ref{def:plans} sets $q_{jd} = \Gamma\bigl(B_j(s,d)-B_j(s,d-1)\bigr)$ with $\Gamma$ a
finite ordered coarsening and places no ceiling of $\bar z$ on its image, and the non-engaging mark
is the constant Hold path of clause (TR-ii) of Assumption~\ref{asm:TR}. Noise is the ternary mark of
Definition~\ref{def:primitives} and observed flow is $X_0 = q_0+z_0$, so the conditional noise
probabilities allocate as
\begin{equation}
\label{eq:L3-exampleA}
\begin{array}{c|c|ccccc}
a\ (\text{prior}) & q_0 & X_0 = -\bar z & X_0 = 0 & X_0 = \bar z & X_0 = 2\bar z & X_0 = 3\bar z\\
\hline
1\ (\tfrac12) & 2\bar z & \text{---} & \text{---} & \kappa/2 & 1-\kappa & \kappa/2\\
0\ (\tfrac12) & 0 & \kappa/2 & 1-\kappa & \kappa/2 & \text{---} & \text{---}
\end{array}
\end{equation}
Reading \eqref{eq:L3-exampleA} cell by cell: at $X_0 \in \{2\bar z,3\bar z\}$ only the engaging type
contributes, so $\pi = 1$ at joint mass
$\tfrac12(1-\kappa)+\tfrac12\cdot\tfrac{\kappa}{2} = \tfrac{2-\kappa}{4}$; at
$X_0 \in \{-\bar z,0\}$ only the non-engaging type contributes, so $\pi = 0$ at the same joint mass;
and at $X_0 = \bar z$ both contribute, each with $\tfrac12\cdot\tfrac{\kappa}{2}$, so
$\pi = \tfrac12$ at mass $\tfrac{\kappa}{2}$. The support is $\{0,\tfrac12,1\}$, which is
$\{0,\bar\pi/2,\bar\pi\}$ with $\bar\pi = 1$, and every one of the three points is free of $\kappa$:
the two extreme cells are fully revealing at every $\kappa$, and the middle cell's posterior is
$\tfrac12$ because the two types reach $X_0 = \bar z$ through noise realisations of equal probability
$\kappa/2$, which cancels. The weights are $A_1(\kappa) = A_0(\kappa) = \tfrac{2-\kappa}{4}$ and
$A_{1/2}(\kappa) = \tfrac{\kappa}{2}$, summing to one, with $A_0' = A_1' = -\tfrac14$ and
$A_{1/2}' = +\tfrac12 = -2\cdot(-\tfrac14)$: Assumption~\ref{asm:Atau}'s derivative pattern holds
exactly with $\Akap = -\tfrac14$. The moment check of Step~15 confirms it, since
$A_{1/2}\tfrac12+A_1\cdot 1 = \tfrac{\kappa}{4}+\tfrac{2-\kappa}{4} = \tfrac12 = \varpi$, free of
$\kappa$. What raising $\kappa$ does here is move mass out of the two revealing end cells into the
single pooling cell, symmetrically and one for one, which is the coefficient pattern $(+1,-2,+1)$
that Step~8 factors. Two features of the construction are load-bearing and neither is a
normalisation. The informed mark must lie strictly outside the reach of the uninformed mark plus
noise, $2\bar z > 0+\bar z$, which is what makes both end cells fully revealing and pins the support
endpoints at $0$ and $1$ at every $\kappa$. And the pre-order engagement share must be exactly
$\tfrac12$, which is what puts the pooling cell's posterior at the midpoint of the chord rather than
elsewhere in it; the word ``symmetric'' in Assumption~\ref{asm:Atau} carries that second requirement.

This example has $\bar\pi = 1$ and so cannot be swept toward zero, which part~(c) requires. A family
in the same class with $\bar\pi$ free is obtained directly: put atoms at $\{0,\bar\pi/2,\bar\pi\}$
with $A_1(\kappa) = A_0(\kappa) = \alpha-c\kappa$ and $A_{1/2}(\kappa) = 1-2\alpha+2c\kappa$, so that
$\Akap = -c$ and the derivative pattern holds by inspection; $\alpha = 0.4$ and $c = 0.3$ keep all
three weights in $[0.1,0.8]$ for $\kappa \in [0,1]$. This is a genuine information structure and not
a list of numbers. For a binary state with prior
$\varpi := A_{1/2}\tfrac{\bar\pi}{2}+A_1\bar\pi$, which Step~15 shows is $\kappa$-free and which here
equals $\bar\pi/2$, the likelihoods
\begin{equation}
\label{eq:L3-likelihood}
\Prb(\text{signal } i \mid a=1) = \frac{A_i\pi_i}{\varpi},
\qquad
\Prb(\text{signal } i \mid a=0) = \frac{A_i(1-\pi_i)}{1-\varpi}
\end{equation}
over the three signals $i$ with $\pi_i \in \{0,\bar\pi/2,\bar\pi\}$ are nonnegative and sum to
$\varpi/\varpi = 1$ and $(1-\varpi)/(1-\varpi) = 1$ respectively, and Bayes returns
$\Prb(a=1 \mid i) = A_i\pi_i/\bigl(A_i\pi_i+A_i(1-\pi_i)\bigr) = \pi_i$, which are the posited
posteriors.

\emph{Step 17 (a structure that does not satisfy it, and it is the one-round predecessor's own
no-disclosure benchmark).} Keep
the single trading date and the ternary noise, but let the engaging plan's mark be $q_0 = +\bar z$
against a non-engaging mark of $0$, with $\Prb(a=1) = \varpi \in (0,1)$. Then
$X_0 \in \{-\bar z,0,\bar z,2\bar z\}$ and the cells are
\begin{equation}
\label{eq:L3-exampleB}
\begin{array}{c|l|c}
X_0 & \text{contributing } (q_0,z_0)\text{, unnormalised mass} & \text{posterior}\\
\hline
-\bar z & (0,-\bar z):\ (1-\varpi)\tfrac{\kappa}{2}
  & 0\\
0 & (\bar z,-\bar z):\ \varpi\tfrac{\kappa}{2};\quad (0,0):\ (1-\varpi)(1-\kappa)
  & \pi_-(\kappa) = \dfrac{\varpi\kappa/2}{\varpi\kappa/2+(1-\varpi)(1-\kappa)}\\
\bar z & (\bar z,0):\ \varpi(1-\kappa);\quad (0,\bar z):\ (1-\varpi)\tfrac{\kappa}{2}
  & \pi_+(\kappa) = \dfrac{\varpi(1-\kappa)}{\varpi(1-\kappa)+(1-\varpi)\kappa/2}\\
2\bar z & (\bar z,\bar z):\ \varpi\tfrac{\kappa}{2}
  & 1
\end{array}
\end{equation}
The support condition fails twice over in \eqref{eq:L3-exampleB}. The support has four points rather
than three. And the two interior points move with $\kappa$: the likelihood ratio at $X_0 = 0$ is
$(\kappa/2)/(1-\kappa)$, strictly increasing on $(0,1)$, so $\pi_-$ is strictly increasing, while at
$X_0 = \bar z$ it is $(1-\kappa)/(\kappa/2)$, strictly decreasing, so $\pi_+$ is strictly decreasing.
The consequence for this lemma can be written down exactly. Differentiating
$\Eop_\kappa[h] = \sum_x \Prb(X_0=x)\,h(\pi(x))$ gives
\begin{equation}
\label{eq:L3-twoterm}
\partial_\kappa \Eop_\kappa[h]
 \;=\; \underbrace{\sum_x \bigl[\partial_\kappa\Prb(X_0=x)\bigr]h(\pi(x))}_{\text{the term the
   representation keeps}}
 \;+\; \underbrace{\sum_x \Prb(X_0=x)\,h'(\pi(x))\,\partial_\kappa\pi(x)}_{\text{the term it has no
   room for}},
\end{equation}
and the second sum in \eqref{eq:L3-twoterm} is generically nonzero by the monotonicity just shown.
Even the first sum is not proportional to $C_h(\bar\pi)$: by Step~9 it equals
$\sum_i A_i'(h-\ell_h)(\pi_i)$ over four atoms, a weighted sum of chord gaps at two distinct interior
points, and no single scalar multiple of the second difference at the midpoint reproduces it. This is
not a contrived counterexample. It is the structure the one-round predecessor's no-disclosure
benchmark solves: its no-disclosure
order-flow enumeration has four cells whose two middle posteriors are ratios in which the noise
probabilities $1-\kappa$ and $\kappa/2$ appear in numerator and denominator alike and therefore move
with $\kappa$, and only the two chord ends survive the limits $\kappa \downarrow 0$ and
$\kappa \uparrow 1$. The predecessor's own route to the chord is accordingly not the three-point
representation of \eqref{eq:atau} but the chord-gap identity of Step~9, in which the affine part
contributes a $\kappa$-free constant because the unnormalised engagement moment is pinned and the
interior motion is the motion of the gap between the kernel and its chord. Step~9 transfers to that
structure; Step~8's clean proportionality does not.

\emph{Step 18 (whether the two-round pooled cell is in the satisfying class, and the weakest
sufficient conditions).} This is not settled here and is not claimed either way. Three things can be
stated precisely.

First, Step~16 does not settle it, because the example produces $\bar\pi = 1$ and does so by force
rather than by accident. Within the one-round primitives, non-engaging plans carry marks that are
weakly negative or zero, Hold paths being constant and Exit paths weakly decreasing by clause (TR-ii)
of Assumption~\ref{asm:TR}, while Voice plans carry positive increments; so the largest order-flow
realisation is attainable only by a Voice plan and the top posterior atom is $1$. Any one-round
ternary-noise market with a non-degenerate pooled law therefore has $\bar\pi = 1$, and part~(c)'s
limit $\bar\pi \downarrow 0$ is empty in that class.

Second, the two-round structure is the only place in this model where an endpoint below one could
come from. On the flagged cell the posterior is identically one by Definition~\ref{def:prices}, and
the two-round timing of Section~\ref{sec:timing} removes from the pooled cell exactly the histories
that would carry the revealing top atom into the flagged cell. Whether it does so is a separate
question from whether it could, and the answer at the implemented calibration is that it does not:
Remark~\ref{rem:AtauRecord} records $\bar\pi = 1$ at every non-degenerate node there, because
unflagged Voice types still generate fully revealing pooled order flows. The conjecture that the
two-round timing leaves the pooled cell with a top atom strictly below one is therefore false at that
calibration, and nothing below rests on it.

Third, what would have to be shown can be named, and it is narrower than the assumption looks. By
Step~15 it suffices to establish the support condition alone: that the pooled cell's
engagement-posterior law, at fixed policies, is supported on exactly three points
$0 < \bar\pi/2 < \bar\pi$ with none of them varying with $\kappa$. By Step~16 that decomposes into
two checkable requirements, which are the weakest sufficient conditions this lemma can name:
\begin{enumerate}[label=(S\arabic*),leftmargin=3.2em,itemsep=3pt]
\item every pooled order-flow cell is either fully revealing of $a=0$, or fully revealing of $a=1$ up
  to the pooled cell's own ceiling $\bar\pi$, or a single cell whose posterior is $\bar\pi/2$; and
\item that middle cell's likelihood ratio is free of $\kappa$, which in Step~16 came from the two
  contributing types reaching the cell through noise events of equal probability.
\end{enumerate}
Step~17 fails (S1) and (S2) simultaneously. Whether a two-round plan menu can be built to satisfy
(S1) and (S2) across a whole window of $H$ trading dates, where the pooled history is the vector
$(X_0,\dots,X_d)$ and the pooled cell is itself carved out by the stake-path event
$\{B_j(s,H-T)<\tau\}$, is not determined here, and it is the sense in which the question is open.
At the implemented calibration both conditions are refuted together, and by measurement rather than
by argument: Remark~\ref{rem:AtauRecord} records a support carrying $23$ to $767$ distinct posterior
values with no mass at $\bar\pi/2$ at any node, which is (S1), and interior atoms that move with
$\kappa$ at a two-sided Hausdorff distance reaching $0.4608$ between adjacent-$\kappa$ support sets,
which is (S2). That is applicability evidence at one calibration and it weakens no result; the
question is one about the assumption's domain, and a different menu, a different horizon or a different
calibration could still satisfy (S1) and (S2). What the openness costs is stated exactly: it is
load-bearing for this lemma, for leg~3 of Lemma~\ref{lem:L4} and for Part~(B) of
Theorem~\ref{thm:T1} jointly.
\end{proof}

\begin{proof}[Proof of Lemma~\ref{lem:L4} (L4)]
In this proof, $j(\cdot)$ is the plan the frozen cutoff vector selects at each signal,
$\mathcal N := \mathcal C_F(\tau',T)\setminus\mathcal C_F(\tau,T)$ the newly flagged set,
$\nu := \Prb(\mathcal N)$ for its mass, and $\rho_P := 1-\Omega(\tau,T)$ for the pooled mass at the
looser threshold. The three legs are proved in order; nestedness is derived at Step~4 as part of
leg~1 rather than assumed.

\emph{Step 1 (only the clock objects move with the threshold).} At fixed policies the selection map
$j(\cdot)$, the engagement labels and the path family $B_j(s,\cdot)$ are the same objects at $\tau$
and at $\tau'$. Definition~\ref{def:plans} writes the stake path with no threshold argument, whereas
the crossing date, the filing date, the stake at filing and the disclosure indicator all carry one.
In the comparison below, therefore, the only objects that move when the threshold moves are those
four; the path itself is common to the two environments.

\emph{Step 2 (the path is a function of the signal and the date alone).} By the no-feedback timing of
Assumption~\ref{asm:nofeedback} the stake path is a function of the plan, the signal and the date
alone. Composing with Step~1's fixed selection map, $d \mapsto B_{j(s)}(s,d)$ is a function of the
signal and the date: it does not depend on the noise draw, on the realised pooled order flow, on the
pooled prices, or on noise-trading intensity.

\emph{Step 3 (product form of the disclosure indicator, at each threshold).} For every plan, every
signal and each of the two thresholds,
\begin{equation}
\label{eq:L4-product}
D_j(s;\tau,T) \;=\; \ind\{a_j = 1\}\cdot\ind\{B_j(s,H-T) \ge \tau\}.
\end{equation}
If $a_j = 0$ the left side of \eqref{eq:L4-product} is zero by Definition~\ref{def:plans}, whose
first conjunct fails, and the right side is zero through its first factor. If $a_j = 1$ the left side
reduces to $\ind\{c_j<\infty,\ f_j \le H\}$, and the conjunct $c_j < \infty$ is redundant: with
$f_j = c_j+T$ and $T$ finite, $f_j \le H$ implies $c_j \le H-T < \infty$, while the reverse inclusion
is immediate. Lemma~\ref{lem:D1}\,(b), the clock equivalence, then converts $\{f_j \le H\}$ into
$\{B_j(s,H-T) \ge \tau\}$. Two hypotheses put that citation inside its domain.
Assumption~\ref{asm:A4} is what makes the crossing date the first passage and pins the filing to land
exactly at $c_j+T$, which is what Lemma~\ref{lem:D1}\,(b) is an equivalence about; and the
restriction $b_0 < \tau' < \tau$ imposes the core domain at \emph{both} thresholds, without which
Lemma~\ref{lem:D1}\,(b) would be cited off its domain at the tighter one. Combining with Step~1, the
realised indicator $D(s;\tau,T) = \ind\{a_{j(s)}=1\}\cdot\ind\{B_{j(s)}(s,H-T)\ge\tau\}$ is a
function of the signal alone, measurable by Lemma~\ref{lem:D1}\,(a), and free of the noise draw and
of $\kappa$.

\emph{Step 4 (leg~1, the inclusion).} Fix any control-node history with $D_j(s;\tau,T) = 1$. By
\eqref{eq:L4-product} this means $a_j = 1$ and $B_j(s,H-T) \ge \tau$. Since $\tau' < \tau$, the
inequality between real numbers $B_j(s,H-T) \ge \tau > \tau'$ gives $B_j(s,H-T) \ge \tau'$, and
applying \eqref{eq:L4-product} at $\tau'$ to the same plan and signal gives $D_j(s;\tau',T) = 1$. The
\emph{same} number $B_j(s,H-T)$ appears in both comparisons, and this is exactly where fixed policies
and Step~2 are consumed: without path fixity the left-hand side of the comparison at $\tau'$ would be
a different path's stake and the implication would not go through. Since the history was arbitrary,
$\mathcal C_F(\tau,T) \subseteq \mathcal C_F(\tau',T)$, and taking complements in the exclusive and
exhaustive partition of Definition~\ref{def:prices} gives
$\mathcal C_P(\tau',T) \subseteq \mathcal C_P(\tau,T)$: tightening the threshold can only shrink the
pooled cell. Nestedness is a conclusion of this step, not an input to it.

\emph{Step 5 (leg~1, every newly flagged history is generated by a Voice plan).} Every history in
$\mathcal N$ satisfies $D(\tau') = 1$ by the definition of that set, and clause (TR-iii) of
Assumption~\ref{asm:TR} records $D = 1 \Rightarrow a = 1$, which holds because the disclosure
indicator carries $a_j = 1$ as a conjunct. Hence the engagement indicator is identically one on
$\mathcal N$, so every newly flagged history is generated by a Voice plan and, whenever $\nu > 0$,
$\Prb(a=1 \mid \mathcal N) = 1$. This is derived, not assumed; it is consumed at Step~9 and its force
is analysed at Step~11.

\emph{Step 6 (leg~1, the weight comparison).} By Assumption~\ref{asm:A1} the joint law of the
primitives is itself a primitive and does not depend on the threshold, which enters only through the
event $\{D=1\}$. Monotonicity of a probability measure applied to Step~4's inclusion gives
$\Omega(\tau',T) = \Prb\bigl(\mathcal C_F(\tau',T)\bigr) \ge \Prb\bigl(\mathcal C_F(\tau,T)\bigr)
= \Omega(\tau,T)$, completing leg~1.

\emph{Step 7 (the newly flagged set, read on the pooled side).} By Steps~4 and~6 together with
exhaustiveness, $\mathcal N = \mathcal C_P(\tau,T)\setminus\mathcal C_P(\tau',T)$; and since
$\mathcal C_F(\tau,T)$ and $\mathcal N$ are disjoint with union $\mathcal C_F(\tau',T)$, finite
additivity gives $\nu = \Omega(\tau')-\Omega(\tau) \ge 0$, the inequality by Step~6. Read on the
pooled side this is the disjoint union
$\mathcal C_P(\tau,T) = \mathcal C_P(\tau',T)\ \uplus\ \mathcal N$.

\emph{Step 8 (both shares are defined).} The hypothesis $\Omega(\tau',T) < 1$ gives
$\Prb\bigl(\mathcal C_P(\tau',T)\bigr) = 1-\Omega(\tau') > 0$, and by Step~6
$\Omega(\tau) \le \Omega(\tau') < 1$, so $\rho_P = 1-\Omega(\tau) > 0$ as well. Both pooled
engagement shares $\bar\pi_{\mathrm{pr}}(\tau) = \Prb\bigl(a=1 \mid D(\tau)=0\bigr)$ and
$\bar\pi_{\mathrm{pr}}(\tau') = \Prb\bigl(a=1 \mid D(\tau')=0\bigr)$ are therefore defined. Were
$\Omega(\tau') = 1$ instead, the pooled cell at the tighter threshold would be null and
$\bar\pi_{\mathrm{pr}}(\tau')$ undefined rather than imputed, which is the treatment
Lemma~\ref{lem:L1} gives a null cell; the comparison would then have no right-hand side rather than a
wrong one.

\emph{Step 9 (the conditional-probability arithmetic).} Apply finite additivity to the event
$\{a=1\}\cap\mathcal C_P(\tau,T)$ along Step~7's disjoint union. The first resulting term is
$\bar\pi_{\mathrm{pr}}(\tau')(\rho_P-\nu)$, using
$\Prb\bigl(\mathcal C_P(\tau',T)\bigr) = \rho_P-\nu$ from Steps~7 and~8; the second is $\nu\cdot 1$
by Step~5, which is the point at which ``every newly flagged history is Voice'' is consumed; and the
left-hand term is $\bar\pi_{\mathrm{pr}}(\tau)\rho_P$. Dividing by $\rho_P > 0$,
\begin{equation}
\label{eq:L4-convex}
\bar\pi_{\mathrm{pr}}(\tau) \;=\; \Bigl(1-\tfrac{\nu}{\rho_P}\Bigr)\,\bar\pi_{\mathrm{pr}}(\tau')
 \;+\; \tfrac{\nu}{\rho_P}\cdot 1 .
\end{equation}

\emph{Step 10 (leg~2, with its exact identity).} By Step~7, $\mathcal N \subseteq \mathcal
C_P(\tau,T)$, so $0 \le \nu \le \rho_P$ and the weight $\nu/\rho_P$ lies in $[0,1]$, the denominator
being strictly positive by Step~8. Equation \eqref{eq:L4-convex} therefore exhibits the share at the
looser threshold as a convex combination of the share at the tighter one and $1$. Rearranging it
gives the exact identity for the difference between the two shares,
\begin{equation}
\label{eq:L4-share}
\bar\pi_{\mathrm{pr}}(\tau) - \bar\pi_{\mathrm{pr}}(\tau')
 \;=\; \frac{\nu}{\rho_P}\Bigl(1-\bar\pi_{\mathrm{pr}}(\tau')\Bigr) \;\ge\; 0,
\end{equation}
the inequality because $\nu \ge 0$, $\rho_P > 0$ and $\bar\pi_{\mathrm{pr}}(\tau') \le 1$, the last
holding because it is a conditional probability. Hence
$\bar\pi_{\mathrm{pr}}(\tau') \le \bar\pi_{\mathrm{pr}}(\tau)$, which is leg~2. Solving
\eqref{eq:L4-convex} for the tighter share when $\rho_P > \nu$ gives the equivalent closed form
$\bar\pi_{\mathrm{pr}}(\tau') = \bigl(\rho_P\bar\pi_{\mathrm{pr}}(\tau)-\nu\bigr)/(\rho_P-\nu)$,
from which the same difference reads $\nu\bigl(1-\bar\pi_{\mathrm{pr}}(\tau)\bigr)/(\rho_P-\nu)$.

\emph{Step 11 (why engagement share one delivers leg~2 unconditionally).} Were Step~5 weakened, so
that the newly flagged set carried some share $\Prb(a=1\mid\mathcal N) \in [0,1]$ rather than exactly
one, the arithmetic of Step~9 would be unchanged except that the last term of \eqref{eq:L4-convex}
would carry that share, and \eqref{eq:L4-share} would become
$\tfrac{\nu}{\rho_P}\bigl(\Prb(a=1\mid\mathcal N)-\bar\pi_{\mathrm{pr}}(\tau')\bigr)$, whose sign is
the sign of the bracket. The conclusion would then be conditional on the newly flagged histories
being at least as engagement-intensive as the pool they leave, which is an assumption about
\emph{which} histories move. Step~5 supplies the value one, the maximum a conditional probability can
take, so the bracket is nonnegative for every admissible value of $\bar\pi_{\mathrm{pr}}(\tau')$ with
no further restriction. Two corollaries: the inequality is an equality exactly when $\nu = 0$ or
$\bar\pi_{\mathrm{pr}}(\tau') = 1$; and no strict inequality is available without assuming
$\nu > 0$, which is not assumed and is not claimed. Finally, legs~1 and~2 hold at every $\kappa$
simultaneously rather than on average, since by Step~3 the indicator is a function of the signal
alone, by Step~1 so is the engagement label, and by Assumption~\ref{asm:A1} the marginal law of the
signal does not depend on $\kappa$; so $\Omega$, $\nu$ and both shares are invariant to $\kappa$ at
fixed policies, and the statement's common $\kappa$ costs nothing.

\emph{Step 12 (leg~3: the pooled sensitivity in product form at each threshold).}
Lemma~\ref{lem:L3} is consumed here by its statement and not through its proof, with the one
exception noted at Step~13. Its part~(a) gives
$\partial_\kappa\Eop_\kappa[h] = \Akap C_h(\bar\pi)$ at a policy at which
Assumption~\ref{asm:Atau} holds, and reading that through Definition~\ref{def:premium}, which sets
$M_P = \Delta_m\Eop[h \mid D=0]$ and $\mathcal S_P = \lvert\partial_\kappa M_P\rvert$, gives
$\partial_\kappa M_P = \Delta_m\Akap C_h(\bar\pi)$ and hence
$\mathcal S_P = \Delta_m\lvert\Akap\rvert\lvert C_h(\bar\pi)\rvert$ at that policy. Two clauses of
Assumption~\ref{asm:Abr} are what make this available at \emph{both} compared thresholds and in the
form the comparison needs. Clause (br-i) carries the
representation \eqref{eq:atau} for the pooled class under $\tau$ and under $\tau'$, with chord
endpoints $\bar\pi(\tau)$ and $\bar\pi(\tau')$ and coefficients $\Akap(\tau)$ and $\Akap(\tau')$;
clause (br-ii) localises all $\kappa$-dependence of $M_P$ in the weights, so that the total
$\kappa$-derivative that Definition~\ref{def:premium} defines $\mathcal S_P$ to be carries no
composition remainder. The gap that (br-ii) closes is real and not bookkeeping: if the support points
or the kernel moved with $\kappa$, the total derivative would carry the extra term
$\bigl[\tfrac12 A_{1/2}h'(\bar\pi/2)+A_1h'(\bar\pi)\bigr]\partial_\kappa\bar\pi$ together with a term
in $\partial_\kappa h$, and Lemma~\ref{lem:L3} bounds neither. Under the two clauses,
\begin{equation}
\label{eq:L4-SPproduct}
\mathcal S_P(\tau,T) = \Delta_m\,\lvert\Akap(\tau)\rvert\,\lvert C_h(\bar\pi(\tau))\rvert,
\qquad
\mathcal S_P(\tau',T) = \Delta_m\,\lvert\Akap(\tau')\rvert\,\lvert C_h(\bar\pi(\tau'))\rvert .
\end{equation}

\emph{Step 13 (leg~3: transfer from the share to the chord endpoint).} Step~10 is a statement about
the pooled prior engagement share, while \eqref{eq:L4-SPproduct} is written in the chord endpoint.
Clause (br-iv) of Assumption~\ref{asm:Abr} is precisely the assumed link: the endpoint is the upper
support point of the pooled posterior law and is a weakly increasing function of the share, the
\emph{same} function at both thresholds. Applying that function to \eqref{eq:L4-share} gives
$\bar\pi(\tau') \le \bar\pi(\tau)$. Without (br-iv), leg~2 says nothing about the endpoint, since
leg~2 moves a share and the chord is drawn to a support point. The identity branch, in which the
endpoint is read as the mean of the pooled law, is excluded as degenerate by (br-iv) itself and, at
the one point where this proof reads a step of Lemma~\ref{lem:L3}'s proof rather than its statement,
by Step~15 there: under that branch the pooled law is a point mass at the endpoint
with $\Akap = 0$, so \eqref{eq:L4-SPproduct} returns
$\mathcal S_P(\tau) = \mathcal S_P(\tau') = 0$ and leg~3 would hold only because both sides vanish.

\emph{Step 14 (leg~3: the chord comparison).} Assumption~\ref{asm:Atau} maintains
$\lvert C_h\rvert$ weakly increasing in the endpoint, and this is used as the maintained orientation
rather than derived. That property orders two values of \emph{one} functional at a single policy,
while the comparison here reads the chord functional at $\tau$ and at $\tau'$; clause (br-v) of
Assumption~\ref{asm:Abr} is what makes those one functional, by requiring the chord functional and
the kernel it is built from to be the same functions of the posterior at both compared thresholds.
The two clauses do different jobs and both are needed: (br-i) fixes the endpoints and the
coefficients, (br-ii) freezes the support points and the kernel along $\kappa$ only, and (br-iv)
speaks about the endpoint map rather than about the kernel. With them and Step~13,
\begin{equation}
\label{eq:L4-chord}
\lvert C_h(\bar\pi(\tau'))\rvert \;\le\; \lvert C_h(\bar\pi(\tau))\rvert .
\end{equation}
Only the magnitude half of the maintained orientation is read here. The sign half $C_h \le 0$ is
consumed at no step of this leg: \eqref{eq:L4-chord} and \eqref{eq:L4-SPproduct} carry absolute
values throughout, and Step~16 multiplies nonnegative numbers.

\emph{Step 15 (leg~3: the coefficient comparison).} Clause (br-iii) of Assumption~\ref{asm:Abr} gives
$\lvert\Akap(\tau')\rvert \le \lvert\Akap(\tau)\rvert$. This clause cannot be dispensed with. The
coefficient is a property of the pooled information structure, and by Step~7 the pooled class at
$\tau'$ is a strictly different collection of histories from the pooled class at $\tau$ whenever
$\nu > 0$, so nothing in the model ties the two coefficients together; Definition~\ref{def:premium}
bounds the coefficient on $[0,1]$ but says nothing about how it moves with the policy.

\emph{Step 16 (leg~3).} All four factors in \eqref{eq:L4-SPproduct} are nonnegative reals and
$\Delta_m > 0$ by clause (TR-i) of Assumption~\ref{asm:TR}. Multiplying the two weak inequalities
\eqref{eq:L4-chord} and Step~15 and reading \eqref{eq:L4-SPproduct} at each threshold,
\begin{equation}
\label{eq:L4-leg3}
\mathcal S_P(\tau',T)
 = \Delta_m\lvert\Akap(\tau')\rvert\,\lvert C_h(\bar\pi(\tau'))\rvert
 \;\le\; \Delta_m\lvert\Akap(\tau)\rvert\,\lvert C_h(\bar\pi(\tau))\rvert
 = \mathcal S_P(\tau,T),
\end{equation}
which is leg~3. It is the only leg that rests on Assumption~\ref{asm:Abr}.

\emph{Step 17 (the vanishing-chord case).} The maintained orientation is weak, so the boundary case
belongs to the statement rather than being assumed away. If $C_h(\bar\pi(\tau)) = 0$ then
\eqref{eq:L4-chord} gives
$0 \le \lvert C_h(\bar\pi(\tau'))\rvert \le \lvert C_h(\bar\pi(\tau))\rvert = 0$, so the chord
vanishes at the tighter threshold too and \eqref{eq:L4-SPproduct} returns
$\mathcal S_P(\tau') = \mathcal S_P(\tau) = 0$. The conclusion then holds with equality, which is the
equality clause of leg~3.

Three limits on the reach of this lemma are part of it. Every statement above is at fixed policies:
if the blockholder re-optimises when the threshold moves (a signal that selects an aggressive
Voice plan at $\tau$ switching to Hold at $\tau'$ once the earlier certain flag removes the
pooled-round camouflage), then the engagement label at that signal changes, Step~4's inclusion
fails, and the flagged weight can fall with a tighter threshold. That is the general-equilibrium
cutoff response, which is the subject of Proposition~\ref{prop:C1}, and this lemma must not be quoted
as an equilibrium statement. No strict inequality is claimed at any leg, since $\nu > 0$ is nowhere
assumed. And nothing here concerns the window margin: the window is held fixed throughout, and a
threshold change bundled with a longer window breaks Step~4 outright, since a history crossing
weakly earlier at $\tau'$ may then file after the horizon.
\end{proof}


%==============================================================================
\section{Proof of Proposition~\ref{prop:P1} (existence)}
\label{sec:proofs-existence}
%==============================================================================

\begin{proof}[Proof of Proposition~\ref{prop:P1} (P1: existence)]
The argument builds the equilibrium object rather than characterising it. Parts~A and~B fix a
conjectured cutoff vector and solve the price system it induces, separating a finite pooled layer
from a flagged layer that is never a finite nonempty indexed family; Part~C supplies beliefs at every
history that carries a requirement, on path and off; Part~D discharges sequential optimality of the
flagged component at every flagged pair; Part~E builds the outer best-response map, applies Brouwer's
fixed-point theorem and assembles the six requirements of Definition~\ref{def:cpbe}; Part~F is the
addendum under Assumption~\ref{asm:A8} (A8). The parameter vector $\vartheta$ is fixed throughout and
suppressed where it does not vary.

\smallskip\noindent\emph{Notation for this proof.} Write $\sigma_F$ for a generic value of
the flagged tuple $\Sfl$ of Definition~\ref{def:prices}, and
$[\underline s,\overline s]$ for the common signal bracket underlying the polytope $\Theta$ of
Definition~\ref{def:theta}. Write $\Phi_s$ and $\varphi_s$ for the c.d.f.\ and density of the signal
$s$, reserving $\Phi$ for the standard normal c.d.f.\ of \eqref{eq:entry} and $\phi$ for its density.
At a control-node information set $\mathcal I$ put
\[
\hat v(\mathcal I) = \Eop[v \mid \mathcal I],
\qquad
\pi(\mathcal I) = \Prb(a = 1 \mid \mathcal I),
\qquad
\bar m(\mathcal I) = m_0 + \pi(\mathcal I)\,\Delta_m .
\]
Further symbols ($\mathcal P_{\mathcal I}$, $\mathfrak R_{\mathcal I}$, $A$, $\mathcal G_F$, $\iota_F$, $j_k$,
$j^\star$, $\mathfrak S(k)$, $\mathfrak w$, $\mathfrak m$, $\mathcal Q_j(s)$, $G_j$, $E_j$, $\mu_n$, $L_j$, $w_n$,
$t_n$, $Z_n$, $\Lambda_k$, $\Lambda_u$, $\pi_n$, $(\hat v_\circ,\pi_\circ)$, $P_\circ$,
$d_i(s,k)$, $c_i(k)$, $\operatorname{supp}(z_d)$, $s_F$) are declared where they are first used. The
selection set $\mathfrak S(k)$ is written in fraktur so that it does not collide with the liquidity
sensitivities $\mathcal S$ and $\mathcal S_P$ of Definition~\ref{def:premium}. The engagement cost
$C_j(s)$ and the execution outlay $\mathcal C_j^{\mathrm{trade}}$ of Definition~\ref{def:U} are
always written with their subscripts, so that neither collides with the chord $C_h$ of
\eqref{eq:chord}, the composition ratios $C_\tau,C_T$, or the cells $\mathcal C_F,\mathcal C_P$.

\medskip
\noindent\emph{Part A: the game at a fixed conjecture.}

\smallskip
\noindent\textbf{Step 1 (the conjecture induces a measurable plan-selection map).}
Fix a conjectured cutoff vector $k = (k_1 \le \dots \le k_{J-1})$ in
\[
\Theta \;=\; \bigl\{k \in [\underline s,\overline s]^{J-1} :
\underline s \le k_1 \le \dots \le k_{J-1} \le \overline s\bigr\},
\]
the ordered box built on the common bracket. It is a set of the kind Definition~\ref{def:theta}
requires: nonempty, compact and convex, being the intersection of a cube with the $J-2$ half-spaces
$\{k_i \le k_{i+1}\}$, and nonemptiness is not decoration, since Brouwer's theorem is vacuous without
it. That definition states the properties $\Theta$ must have and names no particular set, so the
choice is made here, once, and this ordered box is the $\Theta$ used throughout. Nothing below
weakens into an assumption the membership $\Tmap(k) \in \Theta$ derived in Step~13: on the ordered
box, the ordering and bracket inequalities established there imply it. Define the map induced by the
conjecture,
\begin{equation}
\label{eq:P1-jk}
j_k(s) \;=\; 1 + \#\bigl\{i \in \{1,\dots,J-1\} : k_i \le s\bigr\}.
\end{equation}
Each set $\{s : k_i \le s\}$ is a half-line, so $j_k$ is a weakly increasing step function of $s$
with values in $\mathcal J$, and it is Borel. This is the object requirement~(i) of
Definition~\ref{def:cpbe} calls a weakly ordered cutoff vector mapping the signal into a plan. The
second clause of Assumption~\ref{asm:A3} (A3), hypothesis (P-3), that the preferred plan is weakly
increasing in $s$, is what makes a representation of this shape the right one for a best response;
Step~13 returns to that point and derives the representation rather than assuming it.

\smallskip
\noindent\textbf{Step 2 (every date-$0$ object is a deterministic Borel function of the pair
$(j,s)$).}
By hypothesis (P-8), the no-feedback timing of Assumption~\ref{asm:nofeedback}, the executed path
carries no feedback from realised order flow or from realised prices. Hence for each $j \in
\mathcal J$ the objects $B_j(s,d)$ for $d = 0,\dots,H$, the pooled marks $q_{jd}(s) =
\Gamma\bigl(B_j(s,d)-B_j(s,d-1)\bigr)$, the crossing date $c_j(s;\tau)$, the filing date $f_j(s) =
c_j(s)+T$, the stake at filing $B_j^F(s) = B_j(s,f_j(s))$ and the flagged order $Q_j^F(s) =
b_j^*(s)-B_j^F(s)$ are functions of $(j,s)$ and of the policy pair $(\tau,T)$ alone.

Measurability in the signal is where care is needed, and it is consumed from the primitive table
rather than inherited. That $s \mapsto B_j(s,d)$ is Borel \emph{for every plan, Exit included}, is
clause (TR-ii) of Assumption~\ref{asm:TR} (TR), carried by hypothesis (P-13) and consumed here
\emph{directly}: the monotonicity $\partial_s B_j \ge 0$ of clause (TR-iii) is imposed on Voice
alone, Exit paths are restricted in $d$ and unrestricted in $s$, and the conclusion of
Lemma~\ref{lem:D1} is measurability of $D$ and of the cell map, not of the objects listed above. The
coarsening $\Gamma$ has finite image by Assumption~\ref{asm:A2p} (A2$'$), hypothesis (P-2), hence is Borel.
The crossing date $c_j(\cdot;\tau) = \inf\{d : B_j(\cdot,d) \ge \tau\}$ is the pointwise minimum,
over the finite calendar supplied by the same assumption, of the indices of the Borel sets
$\{B_j(\cdot,d) \ge \tau\}$, so it is Borel with values in $\{0,\dots,H\} \cup \{+\infty\}$. The
filing objects are Borel through a finite decomposition indexed by the filing date itself, over its
whole admissible range,
\begin{equation}
\label{eq:P1-BF}
\widetilde B_j^F(s) \;=\; \sum_{\ell=T}^{H} \ind\{f_j(s) = \ell\}\; B_j(s,\ell),
\qquad
\widetilde Q_j^F(s) \;=\; b_j^*(s) - \widetilde B_j^F(s),
\end{equation}
each a finite sum of products of Borel functions, hence Borel; indexing the sum by the filing date is
what makes the surviving factor the stake at $f_j(s)$, which is the object defined above. On
$\{D_j = 1\}$ one has $f_j(s) \in \{T,\dots,H\}$ by hypothesis (P-7), so
$\widetilde B_j^F(s) = B_j(s,f_j(s)) = B_j^F(s)$ and $\widetilde Q_j^F(s) = Q_j^F(s)$ there. Off
$\{D_j = 1\}$ every indicator vanishes and the tilde objects are a conventional extension by
$\widetilde B_j^F := 0$; Definition~\ref{def:plans} defines $B_j^F$ and $Q_j^F$ on the flagged set and
no step reads them elsewhere, so the extension is never consumed. Suppressing the tildes on the
flagged set, $B_j^F$ and $Q_j^F$ are Borel there. By hypothesis (P-7), Lemma~\ref{lem:D1} with its own
hypotheses travelling, the disclosure indicator is $D_j(s;\tau,T) = \ind\{a_j = 1\}\,\ind\{B_j(s,H-T)
\ge \tau\}$ (the clock equivalence of part~(b) of that lemma), and is Borel in $s$. Composing with
\eqref{eq:P1-jk}, all of these become Borel functions of the signal alone at the fixed conjecture
$k$. The definitions of $c_j$, $f_j$, $B_j^F$, $Q_j^F$ and $b_j^*$ used here are clause (TR-iii)'s,
again by hypothesis (P-13).

\smallskip
\noindent\textbf{Step 3 (the pooled public-history family is finite; the flagged family is empty or
uncountable, and in neither case a finite nonempty indexed family).}
Definition~\ref{def:prices} sets $\mathcal H_d^P = (X_0,\dots,X_d;\ \text{flag landed by } d)$ with
$X_d = q_{jd}+z_d$. By hypothesis (P-2) the image of $\Gamma$ is finite and $z_d$ takes values in
$\{-\bar z,0,+\bar z\}$, so each $X_d$ takes values in a finite set; the calendar
$\{0,\dots,H\}$ is finite; and the flag coordinate is a single bit. The collection of pooled public
histories is therefore finite.

The flagged layer is different. By Step~2 the flagged tuple $\sigma_F$ is Borel with \emph{codomain}
$[0,\bar b]^2 \times \{1\}$, a continuum: Definition~\ref{def:plans} puts $B_j(s,d) \in [0,\bar b]$
with $s$ Gaussian and imposes monotonicity only, clause (TR-ii) records that stake levels are
continuum-valued, and the finiteness of Assumption~\ref{asm:A2p} covers the plan menu, the image of
$\Gamma$, the noise support and the calendar, not the stake level. What the construction below runs
on is not that codomain but the \emph{image} of the flagged-pair map
$(j,s) \mapsto \bigl(B_j^F(s),Q_j^F(s),1\bigr)$ on $\{(j,s) : D_j = 1\}$, and that image obeys a
dichotomy. If no plan flags, which the existence half admits because Assumption~\ref{asm:A8} is not
among its hypotheses (a threshold above $\bar b$ realises the case), then the flagged-pair set and
its image are empty and Steps~6, 10 and~12 are vacuous. If some plan flags at one signal, that plan
is Voice by Assumption~\ref{asm:A4}; with the disclosure indicator of Step~2 and
$\partial_s B_j \ge 0$ on Voice from clause (TR-iii), its flagged signal set is a nonempty upper set
of the line, hence contains a half-line and is uncountable, and Assumption~\ref{asm:A7J} maps it
injectively, so the image is uncountable as well. Either way the flagged layer cannot be handled as a
finite nonempty indexed family, and that is what forces the two layers to be treated differently in
Steps~5 and~6. The finiteness of the pooled family rests on Definition~\ref{def:prices}'s own
content, in which the flag enters $\mathcal H_d^P$ as a bit and the filing content $B^F$ does not;
were post-filing pooled histories to carry $B^F$, the pooled family would join the continuum and the
selection argument of Step~6 would have to be run there as well.

\medskip
\noindent\emph{Part B: inner prices.}

\smallskip
\noindent\textbf{Step 4 (every control-node pricing fixed point reduces to one scalar equation, and
depends on the information set only through the pair $(\hat v,\pi)$).}
Fix a control-node information set $\mathcal I$. By hypothesis (P-13), clause (TR-iv) of
Assumption~\ref{asm:TR}, the terminal payoff is \eqref{eq:Y}, prices obey $P(\mathcal I) =
\Eop[Y \mid \mathcal I]$ and entry obeys \eqref{eq:entry}; writing $\mathsf B$ for the entry
indicator, \eqref{eq:entry} is $\mathsf B = \ind\{\xi \ge P(\mathcal I)+K+\bar m(\mathcal I)-\bar
S\}$. Given $\mathcal I$, the quantities $P(\mathcal I)$ and $\bar m(\mathcal I)$ are
$\mathcal I$-measurable, so $\mathsf B$ is a function of $\xi$ alone. By hypothesis (P-1),
Assumption~\ref{asm:A1} (A1), $\xi$ is independent of $(v,\varepsilon)$ and of every $z_d$, hence of
$(v,s,z_{0:H})$ jointly and therefore of $(v,a,\mathcal I)$; conditionally on $\mathcal I$ it is
independent of $(v,a)$. Writing $p = \Prb(\mathsf B = 1 \mid \mathcal I)$ and taking conditional
expectations of \eqref{eq:Y} term by term,
\begin{equation}
\label{eq:P1-EY}
\Eop[Y \mid \mathcal I]
\;=\; (1-p)\bigl(\hat v + \pi\Delta_V\bigr) + p\,(P+m_0) + \Delta_m\,p\,\pi
\;=\; (1-p)\bigl(\hat v + \pi\Delta_V\bigr) + p\bigl(P + \bar m\bigr).
\end{equation}
With $\xi \sim N(0,\sigma_\xi^2)$, which is clause (TR-i)'s distributional form carried by hypothesis
(P-13) (what hypothesis (P-1) supplies is the independence just used and the strict positivity of the
variance, not the Gaussian law itself), $p = 1-\Phi\bigl((P+K+\bar m-\bar S)/\sigma_\xi\bigr)$, which
is \eqref{eq:entry} verbatim and lies in $(0,1)$ for every finite $P$. Define the inner pricing map
\begin{equation}
\label{eq:P1-inner}
\mathcal P_{\mathcal I}(P) \;=\; \bigl(1-p(P)\bigr)\bigl(\hat v + \pi\Delta_V\bigr)
+ p(P)\bigl(P+\bar m\bigr),
\qquad
p(P) = 1-\Phi\!\Bigl(\tfrac{P+K+\bar m-\bar S}{\sigma_\xi}\Bigr).
\end{equation}
The requirement $P(\mathcal I) = \Eop[Y \mid \mathcal I]$ is then the scalar equation
$\mathcal P_{\mathcal I}(P) = P$, and the map depends on $\mathcal I$ only through the two scalars
$\bigl(\hat v(\mathcal I),\pi(\mathcal I)\bigr)$. That reduction is what Steps~5 to~8 run on, and it
is the object the trailing paragraph of the proposition calls the scalar reduction.

\smallskip
\noindent\textbf{Step 5 (the pooled layer, in two parts, because only one of them is a fixed point).}
By Step~3 there are finitely many pooled public histories. Step~4's map was derived at a control
node, which is where $\mathsf B$ is a function of $\xi$ alone given the conditioning, so the two
pooled layers are treated separately.

\emph{(a) The pooled control-node cell, $D = 0$ at date $H$.} Here $\mathcal I = \mathcal I_H$ is a
control node and Step~4 applies as derived. Step~7 supplies a unique fixed point of
$\mathcal P_{\mathcal I}$ from hypothesis (P-11), the restriction $m_0 \ge 0$ of \eqref{eq:m0}, and
Step~8 supplies joint continuity in the belief pair $(\hat v,\pi)$. This is a genuine fixed point: the
price appears on both sides of \eqref{eq:P1-inner} through the entry indicator. Neither clause is
taken from Assumption~\ref{asm:A5} (A5), which the proposition does not assume; Steps~7 and~8 do not
depend on the present step, so citing them forward is not circular.

\emph{(b) Intermediate pooled dates $d < H$.} A history $\mathcal H_d^P$ with $d < H$ is not a
control node. Display \eqref{eq:Y} writes the takeover branch as $\mathsf B(P+m_0+a\Delta_m)$ with
$P$ unqualified. Under the reading Step~4 itself adopts, that $P$ is the control-node price
$P(\mathcal I_H)$, so by the tower property
$P_d^P = \Eop\bigl[Y \mid \mathcal H_d^P\bigr]$ is a plain conditional expectation of
already-solved control-node values, with no self-reference and no fixed point, which is the
reading the paragraph after Definition~\ref{def:prices} records. Under the alternative reading, in
which the $P$ inside $Y$ is the price at whichever information set is conditioning, part~(a)'s
fixed-point argument applies at these dates too. The conclusion this step is used for survives on
either reading, and that is why nothing downstream turns on the adjudication.

That conclusion is: a finite family requires no selection argument, and the pooled price family
$k \mapsto \bigl(P_d^P(\mathcal H_d^P;k)\bigr)_{\mathcal H_d^P}$ is a finite vector, indexed by the
finite pooled alphabet of Step~3, each entry an inner root at the belief summaries carried at its own
history: at the control-node cell by (a) with Steps~7 and~8, and at $d < H$ by (b), a conditional
expectation of already-solved control-node values, which integrates over the continuous signal law
and, on flagged branches, over the flagged-tuple law, and is \emph{not} a finite sum over terminal
states. \emph{No continuity of this family in the conjecture $k$ is claimed here.} Continuity in the
belief summaries is the content of Steps~7 and~8; continuity of the composition through the
conditioning is not derived anywhere below (Step~15) and enters only through hypothesis (P-5), which
is the split the paragraph after Assumption~\ref{asm:A5} draws in the same words. Histories carrying
zero probability under $k$, and histories carrying zero probability under every profile, are handled
in Step~9(b) and Step~9(c) respectively.

\smallskip
\noindent\textbf{Step 6 (the flagged layer is a measurably selected family, built rather than
assumed).}
Work on the image of the flagged-pair map of Step~3. If that image is empty, parts (a) to (d) below
are vacuous and no flagged price is constructed or needed. Otherwise Step~3 makes it uncountable, its
elements written $\sigma_F \in [0,\bar b]^2 \times \{1\}$, and the construction applies. A pointwise
statement (at each $\sigma_F$ there is a unique root) does not by itself
yield a \emph{function} of $\sigma_F$ that the model can integrate against, which is what
$M_F = \Delta_m\Eop[h \mid D=1]$ of Definition~\ref{def:premium} and the timing split of
Lemma~\ref{lem:D1} both require. The family is constructed in four moves.

\emph{(a) The engagement posterior is degenerate on the flagged cell.} Assumption~\ref{asm:A4} (A4),
hypothesis (P-4), together with clause (TR-iii) gives $D = 1 \Rightarrow a = 1$, and hypothesis (P-7)
makes $\{D = 1\}$ an event of the control-node history, so $\Prb(a = 1 \mid \sigma_F, D=1) = 1$. This
is the row of Definition~\ref{def:prices} that sets $\pi = 1$ on $\mathcal C_F$, and hence
$\bar m = m_0+\Delta_m = m_1$ on the whole flagged cell, a constant.

\emph{(b) The flagged pricing map therefore depends on one scalar.} By Step~4 and (a), the flagged
map depends on $\sigma_F$ only through $\hat v(\sigma_F;k) = \Eop[v \mid \sigma_F, D=1]$. Write
$\mathcal G_F$ for the map sending a belief $\hat v$ to the root of
$\mathcal P_{\mathcal I}(\cdot)-\mathrm{id}$ at $(\hat v,\pi=1)$. Parts (ii) and (iii) of Step~7 make
$\mathcal G_F$ single-valued, by existence and uniqueness of the root under hypothesis (P-11), and
Step~8 makes it continuous in the belief, indeed $1$-Lipschitz in the $\hat v$ direction at fixed
$\pi$, since $\partial P/\partial\hat v \in (0,1]$ there and $\pi \equiv 1$ on the flagged cell by
(a). This is the one place at which continuity in the belief is load-bearing, because (d) below
composes a Borel map with a continuous one.

\emph{(c) The belief is a Borel function of the tuple.} Two routes are available and both are taken.
First, $\hat v(\cdot;k)$ is a conditional expectation with respect to $\sigma(\sigma_F)$ and is
therefore $\sigma(\sigma_F)$-measurable by construction. Second, under hypothesis (P-6),
Assumption~\ref{asm:A7J} (A7-J), joint tuple injectivity, the map $(j,s) \mapsto \sigma_F$ is
injective on the flagged-pair set and, by Step~2, Borel; both $\mathcal J \times \mathbb R$ and
$[0,\bar b]^2\times\{1\}$ are Borel subsets of Polish spaces, so the Lusin--Souslin theorem supplies
a Borel inverse $\iota_F$ on the image, and
$\hat v(\sigma_F;k) = \mu_v + \beta\bigl(\iota_F(\sigma_F)_s - \mu_v\bigr)$
with $\beta$ the Gaussian projection coefficient of Definition~\ref{def:primitives}. Injectivity
together with the Borel regularity of clause (TR-ii) already delivers the measurable inverse; no
separate selection assumption is introduced, which is the sense in which the measurable-selection
content of Assumption~\ref{asm:A5} is derived here from Assumption~\ref{asm:A7J} (A7-J) and clause (TR-ii)
rather than assumed.

\emph{(d) The family.} Therefore $P^F(\sigma_F;k) = \mathcal G_F\bigl(\hat v(\sigma_F;k)\bigr)$ is
the composition of a Borel map with a continuous map, hence Borel. This is the measurably selected
flagged price family, and it is pinned rather than chosen: uniqueness of the root at each $\sigma_F$
leaves no freedom, so no selection principle is invoked and no two runs of the argument can produce
different families. That is a statement about the \emph{price} family given the belief; the sense in
which the flagged \emph{belief} is pinned is Step~10's, and is weaker, being a version rather than a
forcing.

\smallskip
\noindent\textbf{Step 7 (under $m_0 \ge 0$ the inner root exists and is unique, by derivation).}
Write $A = \hat v + \pi\Delta_V$ for the no-takeover branch value at $\mathcal I$ (a symbol used
only in this step and the next, and never subscripted, so that it does not collide with the weights
$A_0,A_{1/2},A_1$ of \eqref{eq:atau}) and put
\begin{equation}
\label{eq:P1-rho}
\mathfrak R_{\mathcal I}(P) \;=\; \mathcal P_{\mathcal I}(P) - P \;=\; \bigl(1-p(P)\bigr)(A-P) + p(P)\,\bar m,
\end{equation}
continuous in $P$ because $\Phi$ is. By hypothesis (P-11), the restriction $m_0 \ge 0$ of
\eqref{eq:m0}, together with $\Delta_m > 0$ and $\pi \in [0,1]$ from clause (TR-i), we have
$\bar m \ge 0$.

\emph{(i) No root below $A$.} For $P < A$ both terms of \eqref{eq:P1-rho} are nonnegative and the
first is strictly positive, since $p(P) < 1$ by Step~4; hence $\mathfrak R_{\mathcal I}(P) > 0$.

\emph{(ii) A root exists.} We have $\mathfrak R_{\mathcal I}(A) = p(A)\bar m \ge 0$. If $\bar m = 0$ then $P = A$ is a
root. If $\bar m > 0$ then $\mathfrak R_{\mathcal I}(A) > 0$, while as $P \to +\infty$ we have $p(P) \to 0$ and
$(A-P) \to -\infty$, so $\mathfrak R_{\mathcal I}(P) \to -\infty$. An explicit bracket: for
$P \ge \bar S-K-\bar m+\sigma_\xi$ one has $p(P) \le 1-\Phi(1) < 0.159$, whence
$\mathfrak R_{\mathcal I}(P) \le 0.159\,\bar m - 0.841\,(P-A) \le 0$ once additionally $P \ge A+0.19\,\bar m$. The
intermediate value theorem on $\bigl[A,\ \max\{\bar S-K-\bar m+\sigma_\xi,\ A+0.19\bar m\}\bigr]$
delivers a root.

\emph{(iii) The root is unique.} The residual is differentiable, with
$\mathfrak R_{\mathcal I}'(P) = p'(P)\bigl(P+\bar m-A\bigr) + p(P) - 1$ and
$p'(P) = -\phi\bigl((P+K+\bar m-\bar S)/\sigma_\xi\bigr)/\sigma_\xi < 0$. Take any $P \ge A$. Then
$P+\bar m-A \ge \bar m \ge 0$ by hypothesis (P-11), so the first term is at most zero, and
$p(P)-1 < 0$ because $p < 1$ by Step~4. Hence
\[
\mathfrak R_{\mathcal I}'(P) \;<\; 0 \qquad \text{for every } P \ge A,
\]
and not merely at roots. By (i) every root lies in $[A,\infty)$, and a strictly decreasing function
has at most one zero on an interval, so there is at most one root; (ii) supplies existence, so the
root is unique. The global form of the sign statement, rather than its restriction to roots, is what
the implicit-function argument of Step~8 consumes.

On the maintained sign $m_0 \ge 0$ the existence and uniqueness content of Assumption~\ref{asm:A5} is
therefore a theorem rather than an assumption, which is the first of the three parts of the
proposition's trailing paragraph on that assumption, and Step~8 adds its continuity content in the
belief. What is left over is continuity of the \emph{composition} in the conjecture $k$, which runs
through the conditioning pair $(\hat v,\pi)$ rather than through the pricing map; Step~15 takes that
up and says exactly where it is assumed.

\smallskip
\noindent\textbf{Step 8 (the inner root is monotone and non-expansive in the belief, and continuous
in the belief pair).}
Continuity in the belief comes from the same scalar reduction. On \eqref{eq:P1-rho} the residual is
continuously differentiable jointly in $(P,\hat v)$, and by Step~7(iii) $\partial_P\mathfrak R_{\mathcal I} < 0$
strictly at every root, so the implicit-function argument applies to $\mathfrak R_{\mathcal I}(P;\hat v) = 0$ and
yields
\begin{equation}
\label{eq:P1-dP}
\frac{\partial P}{\partial \hat v}
\;=\; \frac{1-p}{\,1-p+\lvert p'(P)\rvert\,(P+\bar m-A)\,}
\;\in\; (0,1],
\end{equation}
the denominator being at least $1-p > 0$ by hypothesis (P-11) and Step~7(i). So the root map is
single-valued, strictly increasing and $1$-Lipschitz in $\hat v$ at fixed $\pi$.

The pooled layer moves both summaries, so the root is needed as a function of the pair
$(\hat v,\pi)$ and not of $\hat v$ alone. Write the residual with both arguments,
\[
\mathfrak R_{\mathcal I}(P;\hat v,\pi) \;=\; \bigl(1-p(P;\pi)\bigr)\bigl(A(\hat v,\pi)-P\bigr) + p(P;\pi)\,\bar m(\pi),
\]
\[
A = \hat v + \pi\Delta_V,
\qquad
\bar m = m_0 + \pi\Delta_m,
\qquad
p(P;\pi) = 1-\Phi\!\Bigl(\tfrac{P+K+\bar m(\pi)-\bar S}{\sigma_\xi}\Bigr).
\]
Both $A$ and $\bar m$ are affine in $(\hat v,\pi)$ by clause (TR-i) and \eqref{eq:m0}, and $\Phi$ is
smooth, so $\mathfrak R_{\mathcal I}$ is $C^1$ jointly in $(P,\hat v,\pi)$ on $\mathbb R \times \mathbb R \times
[0,1]$; by Step~7(iii) in its global form, $\partial_P\mathfrak R_{\mathcal I} < 0$ at every $P \ge A$ and hence at
every root. The implicit-function theorem therefore gives a locally $C^1$ root
$P = P(\hat v,\pi)$ around every belief pair, and the global existence and uniqueness of Step~7(ii)
and~7(iii) make those local root functions agree wherever their domains overlap. Hence the unique
inner root is a single-valued continuous function of $(\hat v,\pi)$ jointly on
$\mathbb R \times [0,1]$. The bound \eqref{eq:P1-dP} is the non-expansiveness of that root in the
$\hat v$ direction at fixed $\pi$, which is what Step~6(b) consumes on the flagged cell, where
$\pi \equiv 1$; the joint statement is what Steps~5(a) and~9(b) consume, where $\pi$ moves.

\medskip
\noindent\emph{Part C: beliefs, on path and off.}

\smallskip
\noindent\textbf{Step 9 (pooled beliefs as limits of one full-support perturbation family, on the
reachable history set, at every $\kappa \in [0,1]$).}
Requirement~(vi) of Definition~\ref{def:cpbe} asks for off-path beliefs as limits of full-support
perturbations. The family is fixed once here and used at every $k \in \Theta$, not only at the fixed
point of Step~16, which is what the deviation payoffs defining the outer map require, since those
payoffs read off-path pooled histories.

\emph{(a) The alphabet and the reachable set.} Definition~\ref{def:primitives} sets
$\Prb(z_d = 0) = 1-\kappa$ and $\Prb(z_d = \pm\bar z) = \kappa/2$ on the whole domain
$\kappa \in [0,1]$, so the \emph{realised} support is
\[
\operatorname{supp}(z_d) = \{0\} \ \text{at } \kappa = 0,
\qquad
\{-\bar z,+\bar z\} \ \text{at } \kappa = 1,
\qquad
\{-\bar z,0,+\bar z\} \ \text{at every } \kappa \in (0,1).
\]
There is no conflict with hypothesis (P-2): the set $\{-\bar z,0,+\bar z\}$ is the noise
\emph{alphabet}, finite at every $\kappa$ and the only thing Step~3's finiteness argument needs,
while $\operatorname{supp}(z_d)$ is the subset of that alphabet carrying positive probability at the
maintained $\kappa$, and it is the object over which a zero-probability argument must be quantified.
Call a pooled history $\mathcal H_d^P$ \emph{reachable} if it carries strictly positive probability
under some plan and a positive-probability set of signals: there are $j \in \mathcal J$ and a Borel
set $S$ with $\Prb(s \in S) > 0$ such that $X_{d'}-q_{jd'}(s) \in \operatorname{supp}(z_{d'})$ for
every $d' \le d$ and every $s \in S$, with the flag coordinate agreeing with $\ind\{f_j(s) \le d\}$.
Reachability so defined is a property of the menu and the noise law alone: it does not depend on $k$,
and it does not depend on the perturbation stage. The whole-history form is what the argument needs,
since a history each of whose marks is individually attainable under \emph{some} plan need not be
attainable under any \emph{one} plan.

\emph{(b) The limit exists at every reachable history, and it is a limit of the joint posterior.}
Index the perturbation by $n$: at stage $n$ every signal type plays every plan $j \in \mathcal J$
with weight at least $1/n$, the remaining mass following $j_k$ of \eqref{eq:P1-jk}. For integers
$n \ge J$ write
\begin{equation}
\label{eq:P1-wn}
w_n(j \mid s) \;=\; (1-t_n)\,\ind\{j = j_k(s)\} + \frac{t_n}{J},
\qquad
t_n := \frac{J}{n} \in (0,1],
\qquad
t_n \downarrow 0,
\end{equation}
for the stage-$n$ mixing weight: the weights are nonnegative, sum to one, and give every plan at
least $t_n/J = 1/n$. (For $n < J$ the display is not a probability distribution, since $t_n > 1$
makes $1-t_n$ negative, and a weight of at least $1/n$ on every plan is infeasible there; the limit
is unaffected by dropping finitely many initial indices.) The perturbation mass is written $t_n$
rather than with a Greek letter because $\varepsilon$ is the signal noise of
Definition~\ref{def:primitives}. With
\[
L_j\bigl(\mathcal H_d^P \mid s\bigr)
= \prod_{d' \le d} \Prb\bigl(z_{d'} = X_{d'}-q_{jd'}(s)\bigr)
\cdot \ind\bigl\{\text{flag coordinate} = \ind\{f_j(s) \le d\}\bigr\}
\]
for the likelihood of the history of (a), the stage-$n$ joint density over $(j,s)$ at that history is
\begin{equation}
\label{eq:P1-mun}
\mu_n(j,s) \;=\; \frac{w_n(j \mid s)\,L_j(\mathcal H_d^P \mid s)\,\varphi_s(s)}
{\sum_{j' \in \mathcal J} \int w_n(j' \mid s')\,L_{j'}(\mathcal H_d^P \mid s')\,\varphi_s(s')\,
\mathrm{d}s'}.
\end{equation}
By hypothesis (P-2) the plan menu is finite and by Step~3 the pooled history alphabet is finite, so
numerator and denominator are finite sums of terms polynomial in $1/n$ with coefficients not
depending on $n$. At a reachable history the denominator is strictly positive for every $n$: the
witnessing pair $(j,S)$ contributes at least $(1/n)\int_S L_j(\mathcal H_d^P \mid s)\varphi_s(s)\,
\mathrm{d}s$, which is strictly positive because each factor is positive on $S$ by the definition of
reachability and, the mark $q_{jd'}(\cdot)$ taking finitely many values, $L_j$ takes finitely many
positive values on $S$ and is bounded below by their minimum; every other term is nonnegative. A
ratio of polynomials in $1/n$ whose denominator is nonzero for all large $n$ converges as
$n \to \infty$, pointwise in $(j,s)$; write $\mu_\infty$ for the limit. Which law $\mu_\infty$ is
depends on the same denominator split the envelope below needs, and the two cases are \emph{not} the
same law; both are displayed there. This is where the finiteness clauses pay: with a continuum of
pooled histories the limit would need a separate argument.

The joint form is what the step must deliver rather than a flourish. Step~4 shows that the pricing
map reads the information set through $(\hat v,\pi)$; while $\pi$ is a functional of the posterior
over plans alone, $\hat v(\mathcal I) = \Eop[v \mid \mathcal I]$ is a functional of the posterior over
\emph{signals}. Integrating the signal out first would deliver $\pi$ and leave $\hat v$ undefined.
Passing from $\mu_n$ to the belief, $\hat v_n = \sum_j \int \bigl(\mu_v+\beta(s-\mu_v)\bigr)
\mu_n(j,s)\,\mathrm{d}s$, and the envelope needs a case split on the denominator. Write
\begin{equation}
\label{eq:P1-Zn}
Z_n \;=\; (1-t_n)\,\Lambda_k + \frac{t_n}{J}\,\Lambda_u
\end{equation}
for the denominator of \eqref{eq:P1-mun} with $w_n$ expanded, where
\[
\Lambda_k = \int L_{j_k(s')}\bigl(\mathcal H_d^P \mid s'\bigr)\varphi_s(s')\,\mathrm{d}s',
\qquad
\Lambda_u = \sum_{j'} \int L_{j'}\bigl(\mathcal H_d^P \mid s'\bigr)\varphi_s(s')\,\mathrm{d}s'
\]
are the unperturbed and the plan-uniform likelihood aggregates.

If $\Lambda_k > 0$, so that the history is on path under $k$, then $Z_n \ge \Lambda_k/2$ for all
large $n$, and $2\lvert \mu_v+\beta(s-\mu_v)\rvert\varphi_s L_j/\Lambda_k$ is an integrable envelope:
the likelihood is a product of noise probabilities and an indicator, so $L_j \le 1$, the envelope's
signal factor grows only linearly in $\lvert s\rvert$, and $\lvert s\rvert\varphi_s$ is integrable
under the Gaussian law of clause (TR-i). Dominated convergence therefore applies, and since
$w_n(j \mid s) \to \ind\{j = j_k(s)\}$ pointwise and $Z_n \to \Lambda_k$,
\[
\mu_\infty(j,s) \;=\; \frac{\ind\{j = j_k(s)\}\,L_j(\mathcal H_d^P \mid s)\,\varphi_s(s)}{\Lambda_k},
\]
the ordinary Bayes posterior generated by the conjectured map $j_k$: the uniform tremble washes out,
and the limit is \emph{not} the plan-uniform law. If $\Lambda_k = 0$, the $k$-null case, then the
$(1-t_n)$ term vanishes $\Phi_s$-almost everywhere in numerator and denominator alike, so
\[
\mu_n(j,s) \;=\; \mu_\infty(j,s) \;=\; \frac{L_j(\mathcal H_d^P \mid s)\,\varphi_s(s)}{\Lambda_u}
\]
exactly and free of $n$, as an identity between densities almost everywhere rather than pointwise at
every signal, and there is nothing to pass to the limit: the limit is the plan-uniform posterior
restricted to the history. Without the split the bare claim would be false as stated, since
$\mu_n \le \varphi_s L_j/Z_n$ with $Z_n \downarrow 0$ is not a uniform envelope. In both cases the
displayed density is a probability density at the history.

Both belief summaries must be carried, because the pricing map of Step~4 reads the information set
through the pair. Alongside $\hat v_n$ define the stage-$n$ engagement posterior
\[
\pi_n \;:=\; \sum_{j \in \mathcal J} a_j \int \mu_n(j,s)\,\mathrm{d}s \;\in\; [0,1],
\]
which is the engagement posterior $\pi(\mathcal I)$ of Definition~\ref{def:prices} evaluated at
$\mu_n$ rather than a new object. If $\Lambda_k > 0$, then $a_j \in \{0,1\}$ makes
$a_j\mu_n(j,s) \le \mu_n(j,s)$, so the same denominator bound gives the integrable envelope
$2\varphi_s L_j/\Lambda_k$, again with $L_j \le 1$ and $\varphi_s$ integrable under clause (TR-i),
and dominated convergence yields $\pi_n \to \pi_\infty$; if $\Lambda_k = 0$, then $\mu_n$ is free of
$n$ by the display above, so $\pi_n$ is constant and the convergence is immediate. The same two cases
give $\hat v_n \to \hat v_\infty$. Hence $(\hat v_n,\pi_n) \to (\hat v_\infty,\pi_\infty)$ and, by the
joint continuity of the unique inner root in $(\hat v,\pi)$ established in Step~8, the prices converge
with the beliefs, $P_n \to P_\infty$; the bound \eqref{eq:P1-dP} remains available for comparisons at
fixed $\pi$. The limiting belief and the limiting price therefore exist at every reachable pooled
history, on path and off, and on path the belief agrees with the Bayes posterior. This is
load-bearing where it is least obvious: Step~13 evaluates the payoff to \emph{every} plan in the
menu, including plans carrying zero probability under $k$ on a collapse face, whose pooled-execution
bracket reads prices at $k$-null histories, so those prices must be defined, not merely constrained.

\emph{(c) Unreachable histories carry no requirement, and a convention keeps the payoff defined at
every signal.} An unreachable history has probability zero under \emph{every} plan profile,
perturbed or not. It is null under nature rather than off the players' path, so requirement~(vi) of
Definition~\ref{def:cpbe} asks nothing of it and requirement~(iv) prices nothing there. Almost
everything downstream is already clear of such histories: Step~11's pooled-execution bracket
integrates $P_d^P(\mathcal H_d^P)$ against the law of $z_{0:H}$ under the plan actually played, which
for $\Phi_s$-almost every $s$ puts mass only on reachable histories, and the same holds for every
deviation in the round-2 action set of hypothesis (P-9), since those share the same pooled path.

The exceptional signals must nonetheless be named. Reachability in (a) requires a
positive-probability signal set, while the joint mark-and-flag vector
$\bigl((q_{jd'}(s))_{d' \le d},\,\ind\{f_j(s) \le d\}\bigr)$ takes finitely many values by hypothesis
(P-2), so its level sets partition the line into finitely many Borel cells, and a particular cell
may be $\Phi_s$-null while being nonempty. At such a signal, plan $j$'s own realised pooled histories
are unreachable and Step~11's bracket would read a price that (b) has not defined. We therefore fix
once and for all a $k$-independent \emph{reference belief} $(\hat v_\circ,\pi_\circ) \in \mathbb R
\times [0,1]$, for definiteness $(\mu_v,1)$, the prior mean of $v$ of
Definition~\ref{def:primitives} paired with certain engagement, any other pair fixed independently of
the conjecture serving equally well, and assign to every unreachable pooled history the inner root
$P_\circ$ at that belief, the unique solution of
$\mathcal P_{(\hat v_\circ,\pi_\circ)}(P_\circ) = P_\circ$, which exists and is unique by Step~7 under
hypothesis (P-11). Independence of the conjecture is what the convention has to deliver, for two
reasons that meet here: clause (i) of the proposition asks for one price system fixed once and used
at every $k \in \Theta$, and, as the rider below records, the reference belief can move a component
of the outer map through a $\Phi_s$-null set of signals. The payoff is then defined at every triple
$(j,s,k)$, and no requirement of Definition~\ref{def:cpbe} is touched, since requirements (iii), (iv)
and (vi) constrain beliefs and prices only at histories carrying positive probability under some
profile.

The form of the convention matters, and the unconditional mean $\Eop[Y]$ will not serve. That
quantity is an unconditional average of realised control-node values and is in general not a root of
$\mathcal P_{\mathcal I}$ at any belief, so the constructed object would carry prices that are not
inner fixed points at nodes where the blockholder does trade, while the proposition's conclusion says
``prices at their inner fixed points'' without qualification. The precedent of
$P_{-1}^P = \Eop[Y]$ in Definition~\ref{def:prices} does not transfer: that is the pre-trading node, a
different object. With the reference-belief root, every price in the constructed object is an inner
fixed point at the belief carried there. The choice is not innocuous in one narrow respect, and we do
not pretend otherwise: it can change the payoff on a $\Phi_s$-null set of signals, and hence can move
a component of the outer map, which is an infimum over a pointwise condition. What follows is that
the proposition is an existence statement about the object built from this fixed convention; a
different admissible convention yields another equilibrium, not a failure of this one.

\emph{(d) The boundary values of $\kappa$.} At $\kappa \in \{0,1\}$ the realised support degenerates
as displayed in (a): to $\{0\}$ and to $\{\pm\bar z\}$ respectively. A pooled history needing a mark
outside the surviving support is null under \emph{every} plan profile and at every perturbation
stage, so it is off nature's path rather than off the players', it carries no requirement under
Definition~\ref{def:cpbe}\,(vi), and no step above or below reads it. Parts (a) to (c) are quantified
over $\operatorname{supp}(z_d)$ and so cover both endpoints in one sentence; the claim therefore
holds on the full domain $\kappa \in [0,1]$, with no cut to $\kappa \in [0,1)$ anywhere. This is the
extension route of clause~(ii) of the proposition rather than a restriction. What the endpoints
change is \emph{which} histories are reachable, never whether the reachable ones carry a limit; and
the correction runs in the direction of honesty rather than retreat, since the earlier claim that a
limiting belief exists at \emph{every} pooled history is false at $\kappa = 1$ and is withdrawn.
Nothing here asserts that the flagged cell carries positive probability at $\kappa = 1$, and nothing
below does either.

\smallskip
\noindent\textbf{Step 10 (the flagged belief is the point mass at the generating pair, and it is a
version).}
By hypothesis (P-6), Assumption~\ref{asm:A7J} (A7-J), the map $(j,s) \mapsto \sigma_F$ is injective on the
flagged-pair set, so each flagged tuple in its image is generated by exactly one pair, and the Borel
inverse $\iota_F$ of Step~6(c) returns it.

The version must be stated explicitly, because the naive argument is unavailable. The signal is
Gaussian by clause (TR-i), so the pair $(j,s)$ carries probability zero under every stage-$n$
perturbation: what carries positive weight is the plan conditional on the type. A conditional law
given $\sigma_F$ is therefore defined only up to $\Phi_s$-null sets, and the sentence ``that pair has
strictly positive weight, so the perturbed posterior at $\sigma_F$ places probability one on
$\iota_F(\sigma_F)$'' would apply a positive-probability argument to a null event. What is true, and
is all that is needed, is this: because $\iota_F$ is a genuine pointwise Borel map,
$\delta_{\iota_F(\sigma_F)}$ is a \emph{version} of the regular conditional law given $\sigma_F$ at
every stage $n$ (the defining disintegration identity holds tuple by tuple, since the conditioning
$\sigma$-field separates the generating pairs), and it is invariant in $n$, hence its own limit.
This proof selects that version at every tuple in the image; any $\Phi_s$-almost-everywhere-equal
version satisfies requirements (iii) and (vi) of Definition~\ref{def:cpbe} equally well, and nothing
below distinguishes them. So the belief is available and forced up to null sets, not unique.
Consequently the flagged belief is $\hat v(\sigma_F) = \mu_v + \beta\bigl(\iota_F(\sigma_F)_s -
\mu_v\bigr)$ at every image tuple, on path and off, and Step~6's family is simultaneously the on-path
Bayes family and the off-path limit family.

What ``off path'' covers here must be said precisely. It covers flagged tuples generated by pairs
$(j,s)$ that the conjecture $k$ does not select: these are exactly the pairs
Assumption~\ref{asm:A7J} (A7-J) reaches and its on-path form, Assumption~\ref{asm:A7p} (A7$'$), does not. It does
\emph{not} by itself cover tuples outside the \emph{image} of $(j,s) \mapsto \sigma_F$, the tuples a
round-2 deviation to an off-menu order would produce, and no step assigns those a belief. The step
stands here on hypothesis (P-9), the definitional round-2 action-set hypothesis: because the round-2
action set \emph{is} the plan-generated set, no such tuple arises, and the image exhausts the flagged
tuples that can be reached. Off-path beliefs at flagged nodes therefore carry no free parameter,
which is a consequence of Assumption~\ref{asm:A7J} (A7-J) worth recording, since it removes the usual
arbitrariness at exactly the nodes the disclosure mechanism runs through. By hypothesis (P-1) the
pair $(v,\xi)$ remains conditionally independent of the pooled residual given the signal, so nothing
further is needed to price the node.

\medskip
\noindent\emph{Part D: sequential optimality.}

\smallskip
\noindent\textbf{Step 11 (there are exactly two decision points, and the flagged continuation is
deterministic given $(j,s)$).}
Section~\ref{sec:timing} places the plan choice at date $0$ and, when $D = 1$, the flagged order in
round~2, with no re-optimisation in between by hypothesis (P-8). There is thus no pooled decision
node after date $0$: requirement~(ii) of Definition~\ref{def:cpbe}, read on the pooled component, is
satisfied by the timing itself rather than by an argument, and the only genuine sequential-optimality
requirement is the round-2 order. Two features of the decomposition below come from the
flag-termination clause of Assumption~\ref{asm:flagterm}, the second half of hypothesis (P-8), and
not from the no-feedback clause: that the pooled execution runs over $d \le f_j$ and stops there, and
that $Q_j^F = b_j^*(s)-B_j^F(s)$ is the blockholder's \emph{whole} residual position.

On the flagged branch, by Step~2 the objects $B_j^F(s)$ and $Q_j^F(s)$ are deterministic in $(j,s)$;
by Step~6 the flagged price is a function of $\sigma_F$ alone and hence deterministic in $(j,s)$; and
by Definition~\ref{def:prices} the control-node price on that branch is $P^F$. Write $G_j(s;k)$ for
the flagged \emph{trading terms} and $E_j(s;k)$ for the pooled-execution expectation,
\begin{align*}
G_j(s;k) &\;=\; b_j^*(s)\,\Eop\bigl[Y \mid s,j,D=1\bigr] - P^F(\sigma_F)\,Q_j^F(s),\\
E_j(s;k) &\;=\; \Eop_{z}\Bigl[\textstyle\sum_{d \le f_j} P_d^P\bigl(\mathcal H_d^P\bigr)
\bigl(B_j(s,d)-B_j(s,d-1)\bigr)\Bigr],
\end{align*}
the second expectation being over the noise alone. The objective of hypothesis (P-12),
Definition~\ref{def:U} with the display \eqref{eq:Uj}, then splits as
\begin{equation}
\label{eq:P1-Udecomp}
U_j(s;k) \;=\; \underbrace{G_j(s;k) - a_j C_j(s)}_{\text{flagged continuation: deterministic in }(j,s)}
\;-\; \underbrace{E_j(s;k)}_{\text{pooled execution: fixed by the pooled path}} .
\end{equation}
On the flagged branch $a_j = 1$ by
Assumption~\ref{asm:A4}, so the engagement term is $C_j(s)$; the flag factor is carried to match
\eqref{eq:Uj} and changes nothing in Steps~12 and~13, since those run on the flagged set. By
Step~9(c), for $\Phi_s$-\emph{almost every} signal every pooled history the second bracket weighs is
reachable, so the prices it reads are the ones Steps~5 and~9 supply; at the exceptional signals it
reads the reference-belief root of Step~9(c) instead, which is itself an inner root, and which is
what keeps the payoff defined at every signal for Step~13's pointwise argmax. The noise enters the
first bracket nowhere: $\Eop[Y \mid s,j,D=1]$ depends on $(v,\xi)$ and on $P^F$, and $\xi$ is
independent of the noise by hypothesis (P-1) while $P^F$ is noise-free by Step~6.

\smallskip
\noindent\textbf{Step 12 (the flagged component is sequentially optimal at \emph{every} flagged pair,
selected or not: price invariance, cancellation, and the continuation-cost clause).}
This is the step that discharges requirement~(ii) of Definition~\ref{def:cpbe} on the flagged
component, and it does so without appealing to date-$0$ optimality, which is precisely what lets
it cover flagged pairs that the cutoff vector does not select. Such pairs are genuine round-2
information sets: a date-$0$ deviation to a non-selected plan that flags creates one, and
requirement~(ii) binds there exactly as it binds on the selected plan.

Fix a flagged pair $(j,s)$, with no assumption that $j = j_k(s)$, and let $j'$ range over the class
generating the round-2 action set of hypothesis (P-9),
\[
\mathcal Q_j(s) \;:=\; \bigl\{Q_{j'}^F(s) : j' \in \mathcal J \text{ shares } j\text{'s pooled path up
to } f_j(s) \text{ and } a_{j'} = a_j\bigr\},
\]
which is what the flagged-round action set \emph{is} under that hypothesis. The shared pooled path
forces $c_{j'}(s) = c_j(s)$, since the crossing date is a first-hitting index of a path the two
plans share and $c_j \le f_j - T \le f_j$; hence $f_{j'}(s) = f_j(s)$ and $B_{j'}^F(s) = B_j^F(s)$ by
Step~2 and hypothesis (P-7), and $a_{j'} = a_j = 1$ on the flagged set by Assumption~\ref{asm:A4}.
The deviation's flagged tuple is therefore $\sigma_F(j',s) = \bigl(B_j^F(s),Q_{j'}^F(s),1\bigr)$,
in the image of the flagged-pair map where Step~10 applies, and it differs from
$\sigma_F(j,s)$ in the $Q^F$ coordinate alone. Two remarks on the class before the argument. The
sharing relation is an equivalence relation on the flagged set at each signal, since the mutual
agreement just derived makes it symmetric; and it is a genuine pair only on menus carrying two or
more Voice plans that share a pooled path, being the singleton $\{j\}$ on any single-Voice menu,
where Exit and Hold never cross $\tau$ under $b_0 < \tau$.

\emph{(a) The flagged price does not move across the class.} By Step~10 the belief at
$\sigma_F(j',s)$ is the point mass at $(j',s)$, so
$\hat v\bigl(\sigma_F(j',s)\bigr) = \mu_v + \beta(s-\mu_v)$, a function of the signal alone and the
same for every member of the class, and $\pi = 1$ on the flagged cell by Step~6(a), so
$\bar m = m_1$. By Step~4 the inner pricing map depends on the information set only through
$(\hat v,\pi)$. Hence
\begin{equation}
\label{eq:P1-PFinv}
P^F\bigl(\sigma_F(j',s)\bigr) \;=\; \mathcal G_F\bigl(\hat v(s)\bigr) \;=:\; P^F(s)
\qquad \text{for every } j' \text{ in the class,}
\end{equation}
so the round-2 order carries no price impact across the menu. It is Assumption~\ref{asm:A7J} (A7-J) that
pins the belief at the same signal for every class member, and it is uniqueness of the inner root
from Step~7 that makes $P^F(s)$ a number rather than a selection. The argument at this display is
where the joint form of the injectivity hypothesis is indispensable and its on-path form,
Assumption~\ref{asm:A7p} (A7$'$), is not enough: the tuples in question are generated by pairs the cutoff
vector need not select.

\emph{(b) At a flagged node the blockholder values a share at exactly $P^F(s)$.} The blockholder
knows $(j,s)$ and the realised pooled history; given $(j,s)$ the latter is a function of
$z_{0:f_j}$ by Step~2, and the noise is independent of $(v,\varepsilon,\xi)$ by hypothesis (P-1), so
it carries no information about the terminal payoff:
$\Eop\bigl[Y \mid s,j',\mathcal H_{f^-}^P,D=1\bigr] = \Eop\bigl[Y \mid s,j',D=1\bigr]$. By the same
hypothesis $\mathsf B$ is a function of $\xi$ alone given the $\sigma_F$-measurable $P^F$, so with
$p = p\bigl(P^F(s)\bigr)$ from Step~4,
\[
\Eop\bigl[Y \mid s,j',D=1\bigr]
\;=\; (1-p)\bigl(\Eop[v \mid s] + \Delta_V\bigr) + p\bigl(P^F(s)+m_1\bigr).
\]
Assumption~\ref{asm:A7J} (A7-J) makes the market's flagged posterior the blockholder's own: by (a) we
have $\hat v\bigl(\sigma_F(j',s)\bigr) = \Eop[v \mid s]$, so the right-hand side above is
$\mathcal P_{\mathcal I}\bigl(P^F(s)\bigr)$ evaluated at the flagged information set, which equals
$P^F(s)$ because $P^F(s)$ solves the inner fixed point by Steps~4 and~6. Hence
\begin{equation}
\label{eq:P1-EYflag}
\Eop\bigl[Y \mid s,j',\mathcal H_{f^-}^P,D=1\bigr] \;=\; P^F(s).
\end{equation}
No informational rent survives into round~2: full separation is what Assumption~\ref{asm:A7J} (A7-J) buys,
and \eqref{eq:P1-EYflag} is what it costs.

\emph{(c) The flagged-order terms cancel.} The flag terminates the pooled round by
Assumption~\ref{asm:flagterm}, so at round~2 the pooled execution is complete, sunk, and common to
every class member: it is the bracket $E_j(s;k)$ of \eqref{eq:P1-Udecomp}, which the shared path
makes identical across the class for every noise draw. Write $V(j')$ for the round-2 continuation
value of submitting $Q_{j'}^F(s)$: the terminal position valued at the control node, less what the
flagged order costs, less the engagement cost the deviator bears. Then
\begin{equation}
\label{eq:P1-Vcancel}
\begin{aligned}
V(j')
&\;=\; \bigl(B_j^F(s)+Q_{j'}^F(s)\bigr)\,\Eop[Y \mid \cdot\,]
\;-\; P^F(s)\,Q_{j'}^F(s) \;-\; (\text{engagement cost})\\[2pt]
&\;=\; B_j^F(s)\,P^F(s) \;-\; (\text{engagement cost}),
\end{aligned}
\end{equation}
using \eqref{eq:P1-EYflag} for the valuation and \eqref{eq:P1-PFinv} for the price at the deviation
tuple. Every appearance of $Q_{j'}^F$ has cancelled, and with it every appearance of $b_{j'}^*$: the
class member's identity survives only through the engagement cost. In the notation of
\eqref{eq:P1-Udecomp} the same computation reads $G_{j'}(s;k) = B_j^F(s)P^F(s)$ for every class
member, so on flagged plans $U_{j'}(s;k) = B_j^F(s)P^F(s) - C_{j'}(s) - E_j(s;k)$.

\emph{(d) The continuation-cost clause closes the step, and the conclusion is convention-free.}
Display \eqref{eq:Uj} does not date the engagement cost, and Remark~\ref{rem:Ctiming} records the two
available readings. Under the \emph{plan-completion} reading, submitting $Q_{j'}^F(s)$ is completing
plan $j'$, so the deviator bears $C_{j'}(s)$; under the \emph{sunk} reading, the filing has landed
and $D = 1 \Rightarrow a = 1$ is public by Assumption~\ref{asm:A4}, so the engagement cannot be
unmade and the deviator bears $C_j(s)$ whatever order is submitted. Under the sunk reading
\eqref{eq:P1-Vcancel} is constant on the class outright. Under the plan-completion reading,
hypothesis (P-10) (continuation-cost equivalence on this same set, $C_{j'}(s) = C_j(s)$ for every
class member) makes it constant again. Either way
\[
V(j') \;=\; B_j^F(s)\,P^F(s) - C_j(s)
\qquad \text{for every } j' \text{ in the class,}
\]
so every element of $\mathcal Q_j(s)$, the specified $Q_j^F(s)$ included, attains the maximum. The
flagged component is sequentially optimal at every flagged pair $(j,s)$, on the selected plan and off
it, which is requirement~(ii) of Definition~\ref{def:cpbe} in full on its flagged half.

Where hypothesis (P-10) bites is worth stating exactly, since it is vacuous over a large part of its
range. Under the sunk reading it is not consumed at all. Under the plan-completion reading, at a
\emph{selected} $j$ date-$0$ optimality would already do the work: by (c) the within-class comparison
$U_j \ge U_{j'}$ \emph{is} $C_j \le C_{j'}$, which defeats the deviation. But at a flagged pair the
cutoff vector does not select there is no date-$0$ optimality to appeal to, and without the clause the
deviator strictly prefers the class member with the smallest engagement cost. So the clause's bite is
precisely the flagged nodes off the equilibrium plan under the plan-completion reading, and what it
buys is that the conclusion holds under both readings without this proof adjudicating a silence in
the model's own accounting. It is stated as an equality rather than one-sidedly for two reasons: the
sharing relation is symmetric, as shown above, so a clause imposed at every pair of a class in one
direction is imposed in both; and the clause cannot be indexed by the equilibrium's selection, since
Brouwer's theorem does not say which fixed point it returns and the requirement lands at pairs no
cutoff vector selects, where ``the selected plan'' names nothing.

Two things this step does \emph{not} establish should be recorded here rather than left to be
inferred. First, a class on which the \emph{trading} terms differ cannot be built. A witness fixing
$G_{j'}(s;k) - G_j(s;k) = \delta > 0$ across a class is inconsistent with requirements (iv) and (vi)
of Definition~\ref{def:cpbe} in force: by (a) and (b), at pinned beliefs and inner-fixed-point prices
$G_{j'} = B_j^F(s)P^F(s)$ for \emph{every} class member, so $\delta = 0$ necessarily. Such a witness
fixes as a primitive what equilibrium determines, and is available only where the flagged price is
not the fixed point of \eqref{eq:P1-inner} or the flagged belief is not the one
Assumption~\ref{asm:A7J} (A7-J) pins. The live wedge is in the engagement cost, which is where equilibrium
leaves room for one, and that is what hypothesis (P-10) pays for. Second, the converse direction is
about hypothesis (P-9) and not about the argument above. Without it, if round~2 offered the full
interval $[0,\bar b-B_j^F(s)]$, an order outside the plan-generated set would produce a flagged
tuple outside the image of $(j,s) \mapsto \sigma_F$, where Step~10 pins nothing and no step assigns a
belief; the deviation's price, and with it the whole comparison, would be undefined until a belief
were supplied, and the supplied belief would decide the answer. By (c) the on-image flagged payoff is
$B_j^F(s)P^F(s)-C_j(s)$ with the order cancelled out, so an off-image deviation is profitable or not
entirely according to how the assigned off-image belief compares with $\Eop[v \mid s]$, and Step~9's
plan-only perturbation constrains that choice at no stage $n$. Sequential optimality of the flagged
component is therefore not a free consequence of complete contingent plans, and it does not follow
from Assumptions~\ref{asm:A1} through \ref{asm:A7J} (A1 through A7-J) on their own: hypothesis (P-9) is \emph{a} sufficient
condition that delivers it, and it is a restriction on the round-2 action set rather than on the
menu. Strengthening the injectivity hypothesis does not change this, since the menu witnessing the
failure may be taken with $b^*$ strictly increasing on the whole line, so that
Assumption~\ref{asm:A7J} (A7-J) holds on it and requirement~(ii) still fails. Whether hypothesis (P-9) is
the \emph{weakest} such condition is not claimed and is open: the textbook route to sequential
rationality at an unreached node is not a restriction on the action set but a choice of off-path
beliefs, and Definition~\ref{def:cpbe}\,(vi) leaves round-2 orders outside the menu image reached at
no stage $n$, so their limit beliefs are unconstrained by the perturbation used here. Whether some
admissible choice deters every off-menu deviation is a genuine question: by (c) an off-image
deviation to $Q'$ at belief $\hat v'$ earns
$B_j^F\,\Eop[Y \mid \cdot\,] + Q'\bigl(\Eop[Y \mid \cdot\,]-P'\bigr)$ less the cost, so the assigned
belief moves the deviation's gain and the incumbent's own valuation at once, in opposite directions,
and no step of this proof addresses it.

\medskip
\noindent\emph{Part E: the outer map and Brouwer.}

\smallskip
\noindent\textbf{Step 13 (a well-defined weakly ordered best-response map; the ordering and self-map
content of Assumption~\ref{asm:A6} (A6) is a consequence, not an assumption).}
Fix $k \in \Theta$. Steps~5, 6, 9 and~10 determine the pooled and flagged price families and the
belief system, and Step~11 then determines $U_j(s;k)$ for every plan and \emph{every} signal, by the
convention of Step~9(c); each value is finite by hypothesis (P-2), which supplies local boundedness
in $(s,\vartheta)$ pointwise and integrability in expectation. By the second clause of
Assumption~\ref{asm:A3}, hypothesis (P-3), the preferred plan is weakly increasing in the signal, so
the set $\mathfrak S(k)$ of \emph{weakly increasing selections} from
$s \mapsto \arg\max_{j \in \mathcal J} U_j(s;k)$ is nonempty. The selection is named rather than
merely asserted to exist:
\begin{equation}
\label{eq:P1-jstar}
j^\star(\cdot\,;k) \;:=\; \text{the largest element of } \mathfrak S(k),
\qquad
j^\star(s;k) = \max\bigl\{\mathfrak w(s) : \mathfrak w \in \mathfrak S(k)\bigr\}.
\end{equation}
It exists and is itself a weakly increasing selection. \emph{Closure under pointwise maximum:} for
$\mathfrak w_1,\mathfrak w_2 \in \mathfrak S(k)$ the map $s \mapsto \max\{\mathfrak w_1(s),
\mathfrak w_2(s)\}$ takes at each signal one of the two values, both in the argmax there, so it is a
selection, and a pointwise maximum of two weakly increasing maps is weakly increasing.
\emph{Attainment:} the menu is finite by hypothesis (P-2), so the pointwise supremum is a maximum;
fix $s$ and let $\mathfrak w \in \mathfrak S(k)$ attain it there. Then $j^\star$ is a selection, since
$j^\star(s;k) = \mathfrak w(s)$ lies in the argmax at $s$, and it is weakly increasing, since for
$s < s'$ we have $j^\star(s';k) \ge \mathfrak w(s') \ge \mathfrak w(s) = j^\star(s;k)$. So
$j^\star(\cdot\,;k)$ is canonical, single-valued and monotone.

The largest \emph{maximiser} would not do, and the counterexample sits inside
Assumption~\ref{asm:A3}: let $U_2-U_1 \le 0$ everywhere with equality at exactly one signal $s_0$, a
tangential touch, so that there are no crossings and the constant selection $j \equiv 1$ is weakly
increasing, whence both clauses of that assumption hold. The largest maximiser is then $1,\dots,1,
2,1,\dots,1$: it is not weakly increasing, the set $\{s : j^\star \ge 2\} = \{s_0\}$ is not an upper
set, and no cutoff vector represents it, which would break the assembly in Step~17(i). The largest
weakly increasing selection is $j \equiv 1$ there, as it should be. Define now
\begin{equation}
\label{eq:P1-Tmap}
\Tmap_i(k;\vartheta) \;=\; \inf\bigl\{s \in [\underline s,\overline s] :
j^\star(s;k) \ge i+1\bigr\},
\qquad i = 1,\dots,J-1,
\qquad \inf\emptyset := \overline s .
\end{equation}
Since $\{s : j^\star \ge i+2\} \subseteq \{s : j^\star \ge i+1\}$, the infima satisfy
$\Tmap_1(k) \le \Tmap_2(k) \le \dots \le \Tmap_{J-1}(k)$, and every component lies in
$[\underline s,\overline s]$ by construction. The convention $\inf\emptyset := \overline s$ in
\eqref{eq:P1-Tmap} is a \emph{totalisation} and not a code for ``never'': it makes each $\Tmap_i$
defined at every $k$, and under hypothesis (P-5) it is never triggered.

What the bracket clause of hypothesis (P-5) is assumed to deliver can now be written out.
Assumption~\ref{asm:A6} asserts that all best-response cutoffs lie in the common compact ordered
polytope $\Theta$; read at the level Step~14 consumes it, that is the statement that for every
$k \in \Theta$ and every $i \in \{1,\dots,J-1\}$ the set $\{s \in \mathbb R : j^\star(s;k) \ge i+1\}$
is a \emph{nonempty} upper set whose infimum lies in $[\underline s,\overline s]$. Under that clause,
and because $j^\star(\cdot\,;k)$ is weakly increasing, so that each such set is indeed an upper set,
the infima genuinely represent $j^\star$ on the whole signal line: $j^\star(\cdot\,;k)$ agrees with
the induced map $j_{\Tmap(k)}$ of \eqref{eq:P1-jk} at every signal except possibly the finitely many
$\Tmap_i(k)$ at which the upper set fails to contain its own infimum, where the two differ by the
boundary convention alone, and requirement~(i) of Definition~\ref{def:cpbe} pins no convention there.
Weak increase on its own does not deliver the representation; it is the bracket clause that does, and
Step~14 says what that clause excludes.

Hence $\Tmap(k;\vartheta) \in \Theta$ for every $k \in \Theta$, including on the collapse faces where
consecutive components coincide and the corresponding plan carries zero probability, which the
weak inequalities of Definition~\ref{def:cpbe} admit. The ``weakly ordered'' and ``maps $\Theta$ into
itself'' halves of Assumption~\ref{asm:A6} are thus derived here from the monotone-preferred-plan
clause of Assumption~\ref{asm:A3}. What remains genuinely assumed in hypothesis (P-5) is the bracket
$[\underline s,\overline s]$ and the continuity of $\Tmap$; Steps~14 and~15 take those in turn. The
tie-break and the corner convention named in \eqref{eq:P1-jstar} and \eqref{eq:P1-Tmap} are the named
selection the proposition's ``A6 is read'' paragraph requires before
Assumption~\ref{asm:A6} can be applied at all, since a correspondence cannot be called continuous.

\smallskip
\noindent\textbf{Step 14 (the bracket: derivable in the four-action specialisation, assumed at the
level of generality the model is stated at).}
In the four-action version of this model, the bracket is proved rather than assumed. There the
blockholder's payoff to each action is affine in the posterior mean with intercepts bounded uniformly
over conjectures (prices lie in a bounded interval and the entry probability lies in $[0,1]$),
and with totally ordered slopes: zero for Exit, one for Hold and Quiet Voice, and strictly more than
one for Public Voice, while the engagement cost is continuous, strictly positive and strictly
decreasing with full range on the half-line. Each adjacent indifference condition then equates two
affine functions whose slopes differ, except for the Hold and Quiet Voice pair, where the slopes tie
and the comparison reduces to a strictly decreasing cost schedule meeting a bounded constant. Every
indifference signal is therefore finite and bounded uniformly in the conjecture, and the union over
the finitely many adjacent pairs gives one bracket serving all of them.

That argument uses the affine payoff form with ordered slopes, and nothing in this model's
hypotheses imposes it on a general finite menu. Assumption~\ref{asm:A3} imposes single crossing and a
monotone preferred plan only; Definition~\ref{def:U} fixes the accounting of the payoff without
restricting its shape in the posterior mean, and names only plan-locality and integrability as the
properties ever used, neither of which is a shape restriction. At this level of generality the common
bracket is therefore \emph{assumed}, and it is the first of the two things hypothesis (P-5) is doing.

Said in the form Step~13 consumes it: what the bracket clause supplies is that \emph{every}
adjacent-pair indifference signal \emph{exists} and lies in $[\underline s,\overline s]$. That is
exactly what the four-action argument establishes in its own specialisation, and exactly what
excludes both an adjacent pair with no crossing anywhere (in the canonical dominated-top case, where
the top plan is dominated and therefore optimal nowhere) and an adjacent pair whose crossing sits
outside the bracket. The exclusion is a genuine narrowing of the domain rather than a formality. A
menu with a dominated top plan can satisfy every enumerated hypothesis other than (P-5) and fail this
clause, and on such a menu the proposition asserts nothing. What goes wrong there is the
representation itself: under the coding \eqref{eq:P1-jk} the finite cutoff vector cannot represent
the best-response plan map on the Gaussian tails at such a menu, so the components returned by
\eqref{eq:P1-Tmap} would not represent a best response.

\smallskip
\noindent\textbf{Step 15 (continuity: where hypothesis (P-5) assumes rather than derives, and what it
assumes).}
Nothing in Steps~4 to~10 derives continuity of $k \mapsto U_j(s;k)$ at fixed $(j,s)$, and this step
does not assert it. One link of the chain is derived: Step~7(iii) and Step~8 make the inner root
exist, be unique and be continuous in its belief summaries $(\hat v,\pi)$, and by Step~4 those
summaries are the only channel through which the information set enters the price. The other link is
not. The $k$-dependence runs through the \emph{conditioning}: the pair $(\hat v,\pi)$ consists of
ratios of integrals over signal intervals with endpoints $k$, continuous in $k$ only while the
conditioning event's probability stays bounded away from zero, and at a history whose probability
vanishes as $k$ moves it is the construction of Step~9, not a continuity argument, that supplies the
value. The two branches of Step~9(b) are different laws (the $j_k$-concentrated Bayes posterior where
$\Lambda_k > 0$, the plan-uniform posterior where $\Lambda_k = 0$), and no step here shows that they
agree as $\Lambda_k$ falls to zero. The paragraph after Assumption~\ref{asm:A5} draws exactly this
distinction and records the composition $k \mapsto (\hat v,\pi) \mapsto P$ as the underived one, and
Remark~\ref{rem:A6record} records it measured to jump, on the boundaries of the reaching sets of the
pooled histories, at the implemented calibration. The proposition therefore consumes hypothesis (P-5)
as an assumption at precisely this point: for the named largest-weakly-increasing-selection
tie-break, the totalisation of \eqref{eq:P1-Tmap} and the reference-belief convention of Step~9(c),
$\Tmap$ is a continuous self-map of $\Theta$, which is what Step~16 applies. Neither of the two
separate continuities is derived above, so the observation that their conjunction is strictly weaker
than the joint statement the crossing-point argument consumes stands as a logical point and names no
available route.

Conditions (i) and (ii) below are candidate primitive sufficient conditions for that assumption, the
weakest pair we can name, and they are not consequences of the steps above.

\begin{enumerate}[label=(\roman*),leftmargin=2.4em,itemsep=3pt]
\item \emph{Joint continuity.} The map $(s,k) \mapsto U_j(s;k)$ is continuous on
  $[\underline s,\overline s] \times \Theta$ for each plan. This is stated as a condition, not
  inferred from the two separate continuities. It is plausible from the structure (finitely many
  plans, the inner root $1$-Lipschitz in the belief, and $(\hat v,\pi)$ ratios of integrals over
  signal intervals with endpoints $k$), and what it needs in the signal direction is continuity of
  $s \mapsto \bigl(B_j(s,\cdot),b_j^*(s),C_j(s)\bigr)$. Clause (TR-ii) imposes monotonicity on the
  stake path and Borel regularity, and nothing more; a plan acquiring a block discontinuously at a
  signal trigger is permitted and makes the payoff jump, at which point the best-response cutoff is a
  jump point rather than a crossing point and moves discontinuously with $k$.
\item \emph{Robust threshold identification.} For each adjacent pair write
  $d_i(s,k) := U_{i+1}(s;k)-U_i(s;k)$, a function of the signal and of the cutoff vector which always
  carries its adjacent-pair subscript, so that it is not read as the calendar index $d$ of
  Definition~\ref{def:plans}. The condition is that for every $k \in \Theta$ there is a \emph{unique}
  $c_i(k) \in [\underline s,\overline s]$ with
  \[
  d_i(s,k) < 0 \ \text{ for } s < c_i(k),
  \qquad
  d_i(s,k) > 0 \ \text{ for } s > c_i(k),
  \]
  and that across pairs these thresholds are weakly ordered,
  $c_1(k) \le c_2(k) \le \dots \le c_{J-1}(k)$ for every $k \in \Theta$, equivalently that no plan is
  absent from the argmax at every signal. (The threshold $c_i(k)$ carries an adjacent-pair subscript
  and the conjecture as its argument, so it does not collide with the crossing date $c_j(s;\tau)$ of
  Definition~\ref{def:plans}, which carries a signal argument and the threshold policy.) This asks
  strictly more than Assumption~\ref{asm:A3} together with a tie set of empty interior. That
  assumption says the difference \emph{crosses} zero at most once, which admits two failures: an
  interval of ties, and an isolated \emph{tangency}, at which $d_i \le 0$ everywhere with a single
  zero, so that there is no crossing, no sign change and no threshold at all. The counterexample in
  Step~13 that rules out the largest maximiser is that tangency, and it sits inside
  Assumption~\ref{asm:A3}.
\end{enumerate}

Under (i) and (ii), each $\Tmap_i$ equals $c_i$ and each $c_i$ is continuous, by the topological
argument that follows; the implicit-function theorem is the wrong tool here, since it would need the
payoff differentiable in $(s,k)$ and no hypothesis supplies that.

\emph{The identification.} Fix $k$ and $s$ and write $\mathfrak m := \#\{i : s > c_i(k)\}$ for the
number of thresholds strictly below the signal. By the weak ordering the counted indices form a
prefix $\{1,\dots,\mathfrak m\}$, so the increments of the payoff sequence
$U_1(s;k),\dots,U_J(s;k)$ read $d_1,\dots,d_{\mathfrak m} > 0$ and
$d_{\mathfrak m+1},\dots,d_{J-1} < 0$, each sign by the threshold clause: the sequence strictly
increases up to index $1+\mathfrak m$ and strictly decreases after it. Hence, off the finitely many
threshold points, the argmax is single-valued and equal to the counting map
$1+\#\{i : s > c_i(k)\}$, which is \eqref{eq:P1-jk} evaluated at the cutoff vector
$\bigl(c_1(k),\dots,c_{J-1}(k)\bigr)$; at a threshold point $s = c_i(k)$ the continuity in the signal
supplied by (i) forces $d_i(s,k) = 0$, so the argmax there can only widen across the plans tied at
that point. The counting map is weakly increasing in the signal, so it is itself a weakly increasing
selection from the argmax, and every selection agrees with it off the threshold points. Therefore
$j^\star(\cdot\,;k)$ is that counting map, and the vector $\bigl(c_1(k),\dots,c_{J-1}(k)\bigr)$
represents it. Reading the infimum \eqref{eq:P1-Tmap} off that representation, the set
$\{s : j^\star(s;k) \ge i+1\}$ is $\{s : s > c_i(k)\}$ up to the threshold points themselves, and
below $c_i(k)$ the argmax never reaches plan $i+1$, so
$\Tmap_i(k) = \inf\{s : s > c_i(k)\} = c_i(k)$ for every $i$.

\emph{Continuity.} Let $k_n \to k$ in $\Theta$ and take any convergent subsequence
$c_i(k_{n'}) \to c$, which exists because $[\underline s,\overline s]$ is compact. For every $s < c$
one eventually has $s < c_i(k_{n'})$, so $d_i(s,k_{n'}) < 0$, and the joint continuity of (i) gives
$d_i(s,k) \le 0$; symmetrically $d_i(s,k) \ge 0$ for every $s > c$. Uniqueness of the sign threshold
at $k$ forces $c = c_i(k)$. Every convergent subsequence therefore has the same limit, so
$c_i(k_n) \to c_i(k)$, which by the identification just argued is continuity of each component of
$\Tmap$.

These two conditions are sufficient and are not claimed weakest. The proposition continues to assume
their conclusion through hypothesis (P-5), and (i) with (ii) is the weakest pair we can name that
would replace it. Note also that (i) is not independent of hypothesis (P-6): a stake path flat on a
signal interval destroys injectivity there, so continuity cannot be bought by weakening
monotonicity.

Two configurations in which the named sufficient conditions fail are worth recording, because they
show hypothesis (P-5) doing more work than this step makes it look like it is doing. An engagement
cost constant on a signal interval, with the conjecture such that $d_i(\cdot\,,k)$ vanishes
identically there, leaves the pair with no strict sign threshold at all, so condition (ii) fails;
every cutoff in that interval represents a best response, and the selection jumps as the plateau
opens and closes with $k$. Worse, plateaus are structural on exactly the menus
hypothesis (P-10) exists for. By (c) of Step~12,
$U_{j'}(s;k) = B_j^F(s)P^F(s)-C_j(s)-E_j(s;k)$ is the \emph{same function of the signal} for every
member of a deviation class (the shared pooled path gives a common stake at filing and a common
execution bracket, and hypothesis (P-10) gives a common cost), so on any menu carrying two
\emph{adjacent} class members the difference $U_{i+1}-U_i$ vanishes identically on the whole flagged
region, and condition (ii) fails identically rather than on a knife-edge. Neither observation breaks
the proof: hypothesis (P-5) asserts continuity of $\Tmap$ outright, and the selection
\eqref{eq:P1-jstar} is single-valued across a plateau. Nor is any alternative route claimed to
survive such a configuration, since the correspondence-valued argument of Step~18 draws no
conclusion. But it should be said plainly that on multi-Voice shared-path menus that hypothesis is
being assumed at a configuration where its own named sufficient condition provably fails.
Remark~\ref{rem:A6record} carries the separate, measured record that the continuity clause fails at
the implemented calibration for the $\Theta$ this argument constructs.

\smallskip
\noindent\textbf{Step 16 (Brouwer).}
By Step~1 and Definition~\ref{def:theta}, $\Theta$ is nonempty, compact and convex. By Step~13,
$\Tmap(\cdot\,;\vartheta)$ is a single-valued self-map of $\Theta$ under the largest weakly
increasing selection \eqref{eq:P1-jstar} and the corner convention of \eqref{eq:P1-Tmap}, and this is
the reading of Assumption~\ref{asm:A6} the step uses: that hypothesis asserts that $\Tmap$, so
selected, is a continuous self-map of $\Theta$. Brouwer's fixed-point theorem then supplies
$k^\star \in \Theta$ with $k^\star = \Tmap(k^\star;\vartheta)$. The fixed point may lie on a collapse
face, in which case the corresponding plan carries zero probability; the weak inequalities of
Definition~\ref{def:cpbe}\,(i) admit this.

\smallskip
\noindent\textbf{Step 17 (assembling the six requirements of Definition~\ref{def:cpbe}).}
Take $k^\star$ from Step~16 and check the definition item by item.

\begin{enumerate}[label=(\roman*),leftmargin=2.4em,itemsep=3pt]
\item \emph{Weakly ordered cutoff vector.} $k^\star \in \Theta$ by Step~16, and the equilibrium plan
  map is $j^\star(\cdot\,;k^\star)$ of \eqref{eq:P1-jstar}, the largest weakly increasing selection
  from the pointwise argmax, weakly increasing by construction rather than by an appeal to
  Assumption~\ref{asm:A3} after the fact, and represented by $k^\star$ because a weakly increasing
  menu-valued map has upper sets for its upper level sets. It and the induced map $j_{k^\star}$ agree
  off the cutoff points and can differ \emph{at} them, since each $\Tmap_i$ is an infimum that need
  not be attained; requirement~(i) asks for a weakly ordered vector mapping the signal into a plan and
  does not pin the tie convention at the cutoffs, so taking the map to be $j^\star$, optimal at
  \emph{every} signal by construction, is admissible and is what (ii) needs. Consistency with the
  conjecture the prices are built on is immediate: Steps~5, 6 and~9 price against $k^\star$, whose
  induced map is $j_{k^\star}$, and the two maps agree off the finitely many cutoff points, hence
  $\Phi_s$-almost surely, so every conditional probability, posterior and price is the same under
  either. The disagreement is invisible to (iii), (iv) and (v), which are statements about
  probabilities.
\item \emph{Sequentially optimal pooled and flagged components.} Pooled: there is no decision node
  after date $0$, by Step~11. Flagged: Step~12, under hypotheses (P-9) and (P-10), at every flagged
  pair, selected or not. Date-$0$ plan optimality: $k^\star$ is a fixed point of $\Tmap$ and the plan
  map is $j^\star(\cdot\,;k^\star)$, so $j^\star(s;k^\star) \in \arg\max_j U_j(s;k^\star)$ at every
  signal by the definition of the selection. No appeal to indifference at the cutoff points is
  needed, and none is available: indifference there would require continuity of
  $s \mapsto U_j(s;k)$, which is condition (i) of Step~15, explicitly not a hypothesis and explicitly
  not derived, and a plan acquiring a block at a signal trigger, which the primitive table permits,
  makes the payoff jump.
\item \emph{Bayes-consistent on-path beliefs.} Step~9 for pooled histories of positive probability
  under $k^\star$; Step~10 for flagged tuples, where Assumption~\ref{asm:A7J} (A7-J) supplies the point mass
  at $\iota_F(\sigma_F)$ as the selected version.
\item \emph{Competitive pooled and flagged prices at their fixed points.} Step~5 for the finite
  pooled family and Step~6 for the measurable flagged family, both solving
  $P(\mathcal I) = \Eop[Y \mid \mathcal I]$ by Step~4, at the beliefs of (iii) where the history
  carries positive probability under $k^\star$, and at Step~9(b)'s limit belief at the reachable
  histories that do not, since those are the histories the deviation payoffs defining $\Tmap$ read.
  At unreachable histories the requirement asks nothing and Step~9(c)'s convention supplies a value
  that is itself an inner root.
\item \emph{Bidder-entry rule.} The entry probability \eqref{eq:entry} is the one implied by the same
  pair $(P,\pi)$ at each control-node information set, by the derivation in Step~4.
\item \emph{Off-path beliefs as limits of full-support perturbations.} Steps~9 and~10: one family,
  fixed once in \eqref{eq:P1-wn} and used to define the price system at every $k \in \Theta$, at
  every pooled history reachable with positive probability under some plan profile; and, at flagged
  nodes, the version of Step~10 at every tuple in the image of the flagged-pair map, with no tuple
  outside that image arising under hypothesis (P-9).
\end{enumerate}

All six requirements hold, so the assembled object is a cutoff perfect Bayesian equilibrium, and the
construction has been carried out at an arbitrary $\kappa \in [0,1]$, endpoints included, by
Step~9(d). This proves the existence half of the proposition together with its clauses (i) to (v).

\smallskip
\noindent\textbf{Step 18 (a possible route, outside the proposition and not established here).}
Define instead the best-response correspondence
\[
\mathfrak T(k) \;=\; \bigl\{k' \in \Theta : k' \text{ represents some optimal weakly increasing plan
selection at } k\bigr\}.
\]
It is nonempty by Assumption~\ref{asm:A3}, and its values are compact, being closed subsets of the
compact $\Theta$. A Kakutani argument would additionally need a lemma that, under the cutoff encoding
used here, $\mathfrak T$ has convex values and a closed graph, and no such lemma is proved above. The
plateau picture, in which the admissible values of a component form an interval at an indifference
plateau and the ordering constraints cut the product of those intervals by half-spaces, does not by
itself show that every admissible cutoff vector is an independent component-wise choice, and it does
not cover plans absent from the argmax or ties between non-adjacent plans. The maximum theorem, which
concerns choice from a fixed set at a parameter, does not by itself deliver a closed graph for a
\emph{selection-valued} problem whose limits must remain represented under the conventions of
Step~13, and those conventions represent the best response only under the bracket clause of
hypothesis (P-5). No Kakutani conclusion is therefore drawn here. Definition~\ref{def:cpbe} fixes the
Brouwer route for this proposition, so this step is a remark outside the claim, and no step above or
below reads it.

\medskip
\noindent\emph{Part F: Assumption~\ref{asm:A8} and both cells on path.}

\smallskip
\noindent\textbf{Step 19 (at an equilibrium at which Assumption~\ref{asm:A8} holds, both cells carry
positive probability).}
At $k^\star$, hypothesis (P-7) makes $\mathcal C_F = \{D=1\}$ and $\mathcal C_P = \{D=0\}$ exclusive
and exhaustive, so $\Prb(\mathcal C_F) = \Omega(\kappa,\tau,T)$ and
$\Prb(\mathcal C_P) = 1-\Omega(\kappa,\tau,T)$, with $\Omega$ evaluated under the equilibrium plan
map $j^\star(\cdot\,;k^\star)$ of Step~17(i). The value of $\Omega$ is unaffected by the choice
between that map and the induced $j_{k^\star}$, since the two can differ only at the finitely many
cutoff points, a $\Phi_s$-null set. Assumption~\ref{asm:A8} asserts $0 < \Omega < 1$, so both
probabilities are strictly positive: both cells are reached with positive probability under the
equilibrium, that is, both are on path. This is also the condition under which the cell averages
$M_F$ and $M_P$ of Definition~\ref{def:premium} are defined, and clause (TR-iii)'s implication
$D = 1 \Rightarrow a = 1$ makes the flagged cell an engagement cell throughout.

\smallskip
\noindent\textbf{Step 20 (what Assumption~\ref{asm:A8} does and does not do, and the single-threshold
restatement).}
Assumption~\ref{asm:A8} is a restriction on an \emph{equilibrium object}: $\Omega$ is computed at
$k^\star$, not from primitives. No step above rules out $\Omega(k^\star) \in \{0,1\}$, so this
proposition does not produce an equilibrium satisfying that assumption; it states that if the
constructed equilibrium satisfies it then both cells are on path, which is why the addendum is
written ``at any such equilibrium at which Assumption~\ref{asm:A8} holds''. Read literally, Step~19
is close to a restatement of the assumption, and its content is the consistency check that the
partition of Lemma~\ref{lem:D1} is non-degenerate at the fixed point.

The restatement that gives the assumption something to bite on is the following. Suppose in addition
(a) the upper-set engagement-flag hypothesis: the flags $a_j$ equal $1$ exactly on an upper set of
the ordered menu, so that there is $j_a$ with $a_j = 1$ for $j \ge j_a$ and $a_j = 0$ for $j < j_a$;
(b) $\partial_s B_j \ge 0$ on Voice plans, which is clause (TR-iii) and hence already carried by
hypothesis (P-13); and (c) H-ord, Voice stake monotonicity across plans: for Voice plans $j' > j$,
$B_{j'}(s,d) \ge B_j(s,d)$ at every $(s,d)$. Of these only (b) is supplied by the standing
hypotheses; (a) and (c) are the two additions the proposition's final paragraph names, and this is
the one step that consumes them. Definition~\ref{def:plans} orders the menu by aggressiveness without
tying that order either to the engagement flags or to the stake path, which is why neither follows.
Then the flagged set
\[
\bigl\{s : a_{j^\star(s;k^\star)} = 1 \ \text{ and } \ B_{j^\star(s;k^\star)}(s,H-T) \ge \tau\bigr\}
\]
(the equivalence $f_j \le H \iff B_j(s,H-T) \ge \tau$ being part~(b) of Lemma~\ref{lem:D1}) is
an upper interval of signals. The first condition defines an upper set because
$j^\star(\cdot\,;k^\star)$ is weakly increasing, by construction in \eqref{eq:P1-jstar}, together with
(a); and within that set the map $s \mapsto B_{j^\star(s;k^\star)}(s,H-T)$ is weakly increasing,
because it increases in the signal at a fixed plan by (b) and increases across plans by (c). Write
$s_F(k^\star)$ for the infimum of that upper interval, and adopt, \emph{for this step alone}, the
extended-real conventions $\inf\emptyset := +\infty$ and $\inf\mathbb R := -\infty$. These are not
the totalisation $\inf\emptyset := \overline s$ of \eqref{eq:P1-Tmap}, which totalises a component of
the outer map inside the bracket and is a different convention for a different object. Then
\begin{equation}
\label{eq:P1-Omega}
\Omega \;=\; 1 - \Phi_s\bigl(s_F(k^\star)\bigr) \;\in\; [0,1],
\end{equation}
with endpoint values $\Omega = 0$ at $s_F = +\infty$, where nothing flags, as when the threshold
exceeds $\bar b$, and $\Omega = 1$ at $s_F = -\infty$, where everything flags. Hence
\[
0 < \Omega < 1
\qquad \Longleftrightarrow \qquad
s_F(k^\star) \in \mathbb R,
\]
which is Assumption~\ref{asm:A8} restated as a single signal threshold. Without (a) and (c) the
flagged set need not be an interval and $\Omega$ need not be a single-threshold object at all.

\medskip
\noindent\textbf{Scope of the conclusion.}
Six disclaimers travel with the construction above and are stated here rather than left to be
inferred.

First, \emph{uniqueness is not claimed}, in any form: not uniqueness of $k^\star$, not local
uniqueness, not uniqueness of the induced price system, and not uniqueness within a collapse face.
Brouwer's theorem is an existence theorem, and no step bounds $\lVert D_k\Tmap\rVert$:
Assumption~\ref{asm:AGE} (AGE) does that, on a region, and only Proposition~\ref{prop:C1} consumes it.

Second, the addendum is conditional at the equilibrium and nowhere else.
Assumption~\ref{asm:A8} is imposed at the fixed point, Step~20 shows that no step produces a fixed
point with $\Omega \in (0,1)$, and the existence half of the proposition consumes that assumption
nowhere. That an equilibrium satisfying it exists at any given parameter vector is not claimed.

Third, the boundary values of $\kappa$ are handled by extension and not by restriction. At
$\kappa \in \{0,1\}$ a pooled history needing a mark outside the surviving noise support is null
under every plan profile, so it is off nature's path rather than off the players', and
Definition~\ref{def:cpbe}\,(vi) asks nothing of it; the earlier claim that a limiting belief exists at
\emph{every} pooled history is false at $\kappa = 1$ and is withdrawn. No positivity of the flagged
cell at $\kappa = 1$ is asserted anywhere above, and none follows.

Fourth, the equilibrium is an object built from one fixed convention. The reference-belief root of
Step~9(c) can change the payoff on a $\Phi_s$-null set of signals and so can move a component of
$\Tmap$; a different admissible reference belief yields another equilibrium rather than a failure of
this one, and the proposition does not claim convention-independence. In the same spirit, the date at
which the engagement cost is incurred is not settled here: Remark~\ref{rem:Ctiming} records it as open
between two readings, and hypothesis (P-10) is what makes the conclusion the same under both. The
reading of the takeover-branch price in \eqref{eq:Y} is settled, by the model's own convention that
the genuine fixed point sits at the control nodes while earlier pooled prices are tower expectations
of already-solved control-node values; Step~5 records both readings and shows its conclusion survives
either.

Fifth, the flagged order is optimal but not \emph{uniquely} optimal. Step~12 proves the opposite of
uniqueness: on each deviation class the flagged price is invariant and the order cancels, so every
element of $\mathcal Q_j(s)$ delivers the same continuation and the blockholder is exactly
indifferent over the whole action set. The specified $Q_j^F(s)$ is \emph{a} maximiser, which is all
requirement~(ii) of Definition~\ref{def:cpbe} asks; the model does not pin the flagged order by
incentives. This is the round-2 face of the full separation Assumption~\ref{asm:A7J} (A7-J) buys: once the
filing reveals the signal and the price is competitive, no informational rent survives to make one
order strictly better than another.

Sixth, two of the hypotheses are sufficient rather than necessary and are not verified menu by menu.
Assumption~\ref{asm:A7J} (A7-J) is exhibited as satisfiable (the pro-rata single-Voice menu with terminal
target strictly increasing on the whole line satisfies it, and it is the paper's pinned instance),
but nothing above establishes it on a general finite menu, and the discussion after that assumption
records its failure boundary. Hypothesis (P-9) is \emph{a} sufficient condition delivering
requirement~(ii) at flagged nodes and is not claimed to be the weakest, the off-path-belief route
permitted by Definition~\ref{def:cpbe}\,(vi) being untouched by the plan-only perturbation used here;
hypothesis (P-10) inherits the same three disclaimers, being vacuous on any single-Voice menu and a
genuine restriction only on menus whose Voice plans share a pooled path, where no step above verifies
it. Finally, Assumption~\ref{asm:A6} is not claimed to be derivable at this level of generality: its
ordering and self-map content is derived in Step~13, while the bracket of Step~14 and the continuity
of Step~15 are assumed. The bracket clause is a named narrowing of the domain and not a technical
convenience: menus on which some adjacent pair has no indifference signal in the common bracket,
those with a dominated top plan among them, are excluded, and on such a menu the proposition asserts
nothing. Remark~\ref{rem:A6record} carries the measured record that the continuity clause fails at
the implemented calibration for the $\Theta$ this argument constructs.

\medskip
With those six scope statements attached, the proof of the proposition is complete.
\end{proof}

%==============================================================================
\section{Proofs of Theorem~\ref{thm:T1} and Proposition~\ref{prop:C1}}
\label{sec:proofs-ge}
%==============================================================================

\begin{proof}[Proof of Theorem~\ref{thm:T1} (T1)]
Throughout, policies are frozen in the sense of hypothesis~(T-1): the plan menu $\mathcal J$, the
execution policies $B_j(\cdot,\cdot)$, $b_j^*(\cdot)$ and $Q_j^F(\cdot)$, and the cutoff vector $k$
are held fixed in $\kappa$ and held fixed across the two rules compared at each margin. Nothing
below lets $k$ re-solve $k = \Tmap(k;\vartheta)$ at the second rule; that is the whole difference
between this theorem and Proposition~\ref{prop:C1}.

Two pieces of hypothesis traffic are worth naming before the argument starts, because several steps
draw on them. First, hypothesis~(T-5) imports Lemma~\ref{lem:L2} \emph{with its own hypotheses
travelling}, and those hypotheses are Assumption~\ref{asm:A1}, Assumption~\ref{asm:A2p} together
with the primitive table restrictions of Assumption~\ref{asm:TR}, Assumption~\ref{asm:A4},
Assumption~\ref{asm:A5}, Assumption~\ref{asm:A7p} in its on-path injective form,
Lemma~\ref{lem:D1}, the no-feedback timing of Assumption~\ref{asm:nofeedback}, $\Omega > 0$, and an
explicit bidder-entry rule. Where a step below consumes one of them it is cited by name, and it is
consumed as a travelling hypothesis of~(T-5) rather than as a free-standing assumption of this
theorem. The A7 form matters and is therefore named every time: the form
Lemma~\ref{lem:L2} consumes is Assumption~\ref{asm:A7p} (A7$'$), not the weak identification
wording of Assumption~\ref{asm:A7}, which permits two pairs $(j,s)$ with different pooled paths and
is that lemma's first failure case. Second, hypothesis~(T-8), $\kappa$-differentiability of
$M_P$, is supplied by no standing hypothesis of Section~\ref{sec:hypotheses} and is carried here
as a hypothesis of the theorem; Step~4 is where it is used, and the integrability clause of
Assumption~\ref{asm:A2p} is \emph{not} cited for it, since boundedness is not differentiability.

\medskip
\noindent\emph{Part~(A): the factorisation.}

\emph{Step 1 (the two-cell identity).} By hypothesis~(T-2), Assumption~\ref{asm:A8} holds at the
policy under consideration, so $0 < \Omega < 1$ and Lemma~\ref{lem:L1}, imported by
hypothesis~(T-4), applies in its non-degenerate branch: identity~\eqref{eq:L1} holds at
$(\kappa,\tau,T)$, with $M_F$ and $M_P$ the cell averages of Definition~\ref{def:premium}.

\emph{Step 2 (the flagged cell does not move with $\kappa$).} By Lemma~\ref{lem:L2}, imported with
its hypotheses by~(T-5), and at the fixed cutoff and execution policies of~(T-1), the flagged
posterior, the flagged price, the entry probability and $M_F$ are invariant to $\kappa$. Hence
$\partial_\kappa M_F = 0$. This is the step at which Lemma~\ref{lem:L2}'s stack is consumed, and in
particular Assumption~\ref{asm:A7p}: on a menu that violates it the flagged tuple no longer
identifies the informed component of the selected plan, $M_F$ retains direct $\kappa$-dependence,
and everything after this step is void.

\emph{Step 3 (the flagged weight does not move with $\kappa$ either; it is derived, not assumed).}
Hypothesis~(T-7) is PE-$\Omega$, the statement that $\Omega(\cdot,\tau,T)$ does not move with
$\kappa$ at fixed policies. It is derivable at fixed policies, and deriving it is what makes the
partial-equilibrium restriction visible rather than buried. By the clock equivalence of
Lemma~\ref{lem:D1}\,(b), imported by~(T-6), a Voice plan $j$ files by the horizon exactly when
$B_j(s,H-T) \ge \tau$, so with $j(s)$ the plan the frozen cutoff vector assigns to the signal,
\[
D \;=\; \ind\bigl\{a_{j(s)} = 1,\ B_{j(s)}(s,H-T) \ge \tau\bigr\}.
\]
The no-feedback timing of Assumption~\ref{asm:nofeedback}, carried as hypothesis~(T-13), makes
$B_j(s,d)$ a function of $(j,s,d)$ alone: no realised order flow and no realised price enters the
executed path. With the policies frozen by~(T-1), $D$ is therefore a function of $s$ alone. The law
of $s$ carries no $\kappa$: Assumption~\ref{asm:A1} makes $v$ and $\varepsilon$ independent with
strictly positive variances, and the distributional forms of clause~(TR-i) of
Assumption~\ref{asm:TR}, both travelling with Lemma~\ref{lem:L2} under~(T-5), give
$s = v + \varepsilon \sim N(\mu_v,\sigma_v^2+\sigma_\varepsilon^2)$, in which $\kappa$ does not
appear. Indeed $\kappa$ enters the primitives of Definition~\ref{def:primitives} in exactly one
place, the ternary noise-mark law of $z_d$, which reaches observed pooled order flow
$X_d = q_{jd} + z_d$ and reaches nothing that $D$ depends on. Hence $\Omega = \Prb(D=1)$ is
literally the same number at every $\kappa$. The constancy is the form Step~6 needs; the derivative
form $\partial_\kappa\Omega = 0$ is its corollary, and it is the form Step~4 needs.

What this derivation consumes is the freezing of the cutoff vector, not any property of the
disclosure rule. In equilibrium $k$ solves $k = \Tmap(k;\vartheta)$ and moves with $\kappa$,
$\Omega$ moves with it, and the term discarded in Step~4 returns at full size. That term is the
object Proposition~\ref{prop:C1} bounds.

\emph{Step 4 (differentiating in $\kappa$).} Two of the three factors on the right of \eqref{eq:L1}
are constant in $\kappa$: $M_F$ by Step~2 and $\Omega$ by Step~3. Neither step uses the present one,
so nothing here is circular. With those two factors constant, $\kappa \mapsto \Dact$ is an affine
image of $\kappa \mapsto M_P$, which hypothesis~(T-8) makes differentiable; hence
$\partial_\kappa\Dact$ exists and equals $(1-\Omega)\,\partial_\kappa M_P$. Written as one term per
factor that could carry $\kappa$, the same computation reads
\begin{equation}
\label{eq:T1-three}
\partial_\kappa\Dact
\;=\; \Omega\,\partial_\kappa M_F
\;+\; (1-\Omega)\,\partial_\kappa M_P
\;+\; (\partial_\kappa\Omega)\,(M_F - M_P),
\end{equation}
the three terms being the flagged cell's own motion, the pooled cell's own motion, and the
reallocation of mass between the cells. Step~2 kills the first and Step~3 the third. Note what the
third deletion does \emph{not} assume: $M_F - M_P$ is neither small nor signed here, and the term
goes because its coefficient is zero.

\emph{Step 5 (Part~(A)'s factorisation).} By~(T-2), $1-\Omega \in (0,1)$, so
$|1-\Omega| = 1-\Omega$. Taking absolute values in the conclusion of Step~4, and writing
$\mathcal S = |\partial_\kappa\Dact|$ and $\mathcal S_P = |\partial_\kappa M_P|$ as in
Definition~\ref{def:premium},
\begin{equation}
\label{eq:T1-factor}
\mathcal S(\kappa,\tau,T) \;=\; \bigl(1-\Omega(\tau,T)\bigr)\,\mathcal S_P(\kappa,\tau,T),
\end{equation}
exactly, which is the first clause of Part~(A). Step~3 licenses writing the argument of $\Omega$ as
$(\tau,T)$ rather than $(\kappa,\tau,T)$, and that convention is kept from here on. Two consequences
are used later. First, $\mathcal S > 0$ exactly when $\mathcal S_P > 0$, which hypothesis~(T-3)
supplies. Second, when $\mathcal S_P > 0$ the map $x \mapsto |x|$ is differentiable at
$\partial_\kappa M_P \neq 0$, so $\mathcal S_P$ inherits whatever differentiability
$\partial_\kappa M_P$ has in a policy coordinate. That transfer creates no differentiability; it
only moves it, which is why Part~(C)'s local form carries its own smoothness hypothesis.

\emph{Step 6 (the aggregate form, with no differentiability used).} Fix any grid
$\kappa_0 < \kappa_1 < \dots < \kappa_n$ inside the maintained parameter set, with~(T-2) holding at
each node. Applying Step~1 at $\kappa_i$ and at $\kappa_{i+1}$, and using Step~2 and Step~3 in their
\emph{constancy} form ($M_F$ and $\Omega$ are one and the same number at the two nodes, which a
vanishing derivative at a single $\kappa$ would not deliver) gives
$\Dact(\kappa_{i+1}) - \Dact(\kappa_i) = (1-\Omega)\bigl[M_P(\kappa_{i+1}) - M_P(\kappa_i)\bigr]$.
Summing absolute values,
\begin{equation}
\label{eq:T1-tv}
\sum_{i=0}^{n-1}\bigl\lvert \Dact(\kappa_{i+1}) - \Dact(\kappa_i)\bigr\rvert
\;=\; (1-\Omega)\sum_{i=0}^{n-1}\bigl\lvert M_P(\kappa_{i+1}) - M_P(\kappa_i)\bigr\rvert,
\end{equation}
which is the total-variation clause of Part~(A). Hypothesis~(T-8) is not used: \eqref{eq:T1-tv} is
finite differences throughout. The same argument runs for any aggregator of the increment vector
that is positively homogeneous of degree one (total variation, mean absolute slope, the supremum
of the absolute increments), because $1-\Omega$ is one nonnegative scalar common to every
increment. Every ratio statement in Parts~(B) and~(C) therefore holds verbatim with $\mathcal S$ and
$\mathcal S_P$ replaced by their aggregate counterparts. This matters because a measurement
convention that reports $\kappa$-sensitivity as a total variation over a grid rather than as a
pointwise derivative is then measuring the same functional the theorem is about. Part~(A) is
Steps~5 and~6.

\medskip
\noindent\emph{Part~(B): the threshold margin.}

Throughout Part~(B) the window $T$ is common to the two rules and the comparison is
$b_0 < \tau' < \tau$, so $\tau'$ is the tighter threshold. The maintained configuration
$b_0 < \tau$ is clause~(TR-i) of Assumption~\ref{asm:TR}, and it is imposed at both thresholds.

\emph{Step 7 (the chord form of the pooled sensitivity).} By hypothesis~(T-9),
Assumption~\ref{asm:Atau} holds at each compared policy, with $\bar\pi$ read as the upper support
point of the pooled engagement posterior in \eqref{eq:atau} and not as the pooled engagement share.
Clause~(br-ii) of Assumption~\ref{asm:Abr}, imported by hypothesis~(T-11), places all
$\kappa$-dependence of $M_P$ in the weights of \eqref{eq:atau} and delivers
$\partial_\kappa M_P = \Delta_m\,\Akap\,C_h(\bar\pi)$ exactly, with no
composition-through-$\kappa$ remainder; this is the $M_P$-level form of part~(a) of
Lemma~\ref{lem:L3}, imported by hypothesis~(T-10), which gives
$\partial_\kappa\Eop_\kappa[h] = \Akap\,C_h(\bar\pi)$. Taking absolute values, and using
$\Delta_m > 0$ from clause~(TR-i) of Assumption~\ref{asm:TR} as it travels with~(T-5),
\begin{equation}
\label{eq:T1-chord}
\mathcal S_P \;=\; \Delta_m\,\lvert\Akap\rvert\,\lvert C_h(\bar\pi)\rvert
\;=\; \frac{\Delta_m}{4}\,\lvert\Akap\rvert\,\bar\pi^2\,\lvert h''(\zeta)\rvert
\qquad\text{for some }\zeta\in(0,\bar\pi),
\end{equation}
the second equality by part~(b) of Lemma~\ref{lem:L3}, which is an identity rather than an
approximation. The reading used below is that $\mathcal S_P$ is a \emph{product} of a
weight-derivative magnitude and a chord magnitude. That is why leg~3 of Lemma~\ref{lem:L4} needs
both a coefficient comparison and a chord comparison, and neither alone.

\emph{Step 8 (the ratio identity).} Apply \eqref{eq:T1-factor} at $(\tau',T)$ and at $(\tau,T)$,
which~(T-2) permits at both. By~(T-3), $\mathcal S_P(\kappa,\tau,T) > 0$, so
$\mathcal S(\kappa,\tau,T) > 0$ by \eqref{eq:T1-factor} and the quotient is defined:
\begin{equation}
\label{eq:T1-ratio-tau}
\frac{\mathcal S(\kappa,\tau',T)}{\mathcal S(\kappa,\tau,T)}
\;=\; \frac{\bigl(1-\Omega(\tau',T)\bigr)\mathcal S_P(\kappa,\tau',T)}
            {\bigl(1-\Omega(\tau,T)\bigr)\mathcal S_P(\kappa,\tau,T)}
\;=\; W_\tau\,C_\tau,
\qquad
W_\tau = \frac{1-\Omega(\tau',T)}{1-\Omega(\tau,T)},\quad
C_\tau = \frac{\mathcal S_P(\kappa,\tau',T)}{\mathcal S_P(\kappa,\tau,T)},
\end{equation}
with $W_\tau$ and $C_\tau$ the weight-effect and composition-effect ratios of
Definition~\ref{def:premium} at this margin. The identity is exact, being
\eqref{eq:T1-factor} twice and one division, and it consumes~(T-1) to guarantee that the two
sides are one model with one primitive changed rather than two different equilibria.

\emph{Step 9 (the weight ratio lies in $[0,1]$).} By leg~1 of Lemma~\ref{lem:L4}, imported by
hypothesis~(T-12) and proved outright there from the clock equivalence of Lemma~\ref{lem:D1}\,(b)
and $b_0 < \tau' < \tau$, we have $\Omega(\tau',T) \ge \Omega(\tau,T)$. Subtracting from one
reverses the inequality, and both sides are strictly positive by~(T-2), so dividing preserves the
direction: $0 \le W_\tau \le 1$, the lower bound because $1-\Omega > 0$ at both policies.

\emph{Step 10 (the composition ratio lies in $[0,1]$).} By leg~3 of Lemma~\ref{lem:L4},
$\mathcal S_P(\kappa,\tau',T) \le \mathcal S_P(\kappa,\tau,T)$, and that leg holds only under
Assumption~\ref{asm:Abr}. With~(T-3) supplying a strictly positive denominator, and with
$\mathcal S_P$ an absolute value by Definition~\ref{def:premium}, $0 \le C_\tau \le 1$. The chain
leg~3 executes is cited rather than re-derived, but it is worth setting out so that each clause can
be seen carrying its own inch of the argument.
\begin{enumerate}[label=(\roman*),leftmargin=2.4em,itemsep=3pt]
\item A tighter threshold weakly lowers the pooled engagement \emph{share}
  $\bar\pi_{\mathrm{pr}} = \Prb(a=1\mid D=0)$. This is leg~2 of Lemma~\ref{lem:L4}, proved outright
  by conditional-probability arithmetic on the nested flagged sets of leg~1.
\item A weakly lower share gives a weakly lower chord endpoint $\bar\pi$. This is clause~(br-iv) of
  Assumption~\ref{asm:Abr}, and it is the clause that keeps the two objects apart: leg~2 moves a
  share, the chord moves an upper support point, and only~(br-iv) links them. Without it leg~2 says
  nothing about $\bar\pi$.
\item A weakly lower $\bar\pi$ gives a weakly smaller $\lvert C_h(\bar\pi)\rvert$. Two clauses do
  different jobs here. The maintained magnitude monotonicity of $\lvert C_h\rvert$ in $\bar\pi$,
  carried by Assumption~\ref{asm:Atau}, orders two values of \emph{one} functional; and
  clause~(br-v) of Assumption~\ref{asm:Abr} is what makes the values at $\tau$ and at $\tau'$ two
  values of one functional rather than of two, since at fixed policies a change in $\tau$ moves
  which histories are pooled, which moves the pooled price, which enters the kernel
  $h = \pi\,p$ through the entry probability \eqref{eq:entry}. The sign half $C_h \le 0$ of
  Assumption~\ref{asm:Atau} is not used at this leg.
\item $\lvert\Akap(\tau')\rvert \le \lvert\Akap(\tau)\rvert$, which is clause~(br-iii) of
  Assumption~\ref{asm:Abr}.
\item Multiplying the two nonnegative factors of the product form \eqref{eq:T1-chord}, using
  items~(iii) and~(iv), gives leg~3.
\end{enumerate}
Items~(iii)--(v) are unavailable unless the representation \eqref{eq:atau} holds at \emph{both}
compared policies, which is clause~(br-i); and item~(v) is unavailable without clause~(br-ii),
because a composition-through-$\kappa$ remainder in $\partial_\kappa M_P$ would destroy the product
form. The threshold leg therefore rests on Assumption~\ref{asm:Atau} at both policies and on all
five clauses~(br-i)--(br-v) of Assumption~\ref{asm:Abr}.

\emph{Step 11 (Part~(B)'s conclusion).} By Steps~9 and~10 both ratios lie in $[0,1]$, so their
product does: $W_\tau C_\tau \le 1\cdot 1 = 1$. Substituting into \eqref{eq:T1-ratio-tau} and
multiplying through by the strictly positive number $\mathcal S(\kappa,\tau,T)$ gives
\eqref{eq:T1B} and, with it,
$\mathcal S(\kappa,\tau',T) \le \mathcal S(\kappa,\tau,T)$. Threshold tightening attenuates at
fixed policies. The structural point is that no dominance condition appears anywhere in Steps~9--11:
at this margin the weight effect and the composition effect push the same way, so the product is
bounded by one without the two being compared in size. The inequality is weak, and it is an
equality when the move from $\tau$ to $\tau'$ reclassifies no history (when no pair $(j,s)$ has
$\tau' \le B_j(s,H-T) < \tau$), since then $\Omega$, $\bar\pi_{\mathrm{pr}}$, $\bar\pi$ and
$\mathcal S_P$ are all unchanged and both ratios equal one. No strict attenuation is claimed.

\medskip
\noindent\emph{Part~(C): the window margin.}

Throughout Part~(C) the threshold $\tau$ is common to the two rules and the comparison is
$T' < T$, so $T'$ is the tighter window.

\emph{Step 12 (the weight ratio is at most one, and this is proved).} Let $j$ be a Voice plan and
$s$ a signal with $D_j(s;\tau,T) = 1$. By the clock equivalence of Lemma~\ref{lem:D1}\,(b),
$B_j(s,H-T) \ge \tau$. Since $T' < T$ gives $H-T' > H-T$, and since $\partial_d B_j \ge 0$ for Voice
plans by clause~(TR-ii) of Assumption~\ref{asm:TR} as it travels with~(T-5), monotonicity of the
stake path in the date gives $B_j(s,H-T') \ge B_j(s,H-T) \ge \tau$; and the engagement flag $a_j=1$
is untouched by the window. Applying the clock equivalence in the other direction at $T'$ yields
$D_j(s;\tau,T') = 1$. By Assumption~\ref{asm:A4} only Voice plans cross the threshold in the core
model, so no other history has to be checked. Hence $\mathcal C_F(\tau,T) \subseteq
\mathcal C_F(\tau,T')$, so $\Omega(\tau,T') \ge \Omega(\tau,T)$, and, exactly as in Step~9 with
\mbox{(T-2)} supplying a strictly positive denominator and~(T-1) holding the plans fixed across the
two windows,
\[
0 \;\le\; W_T \;=\; \frac{1-\Omega(\tau,T')}{1-\Omega(\tau,T)} \;\le\; 1 .
\]
This inequality is a consequence of Lemma~\ref{lem:D1} and the monotone Voice stake path rather
than a hypothesis of the theorem, and it is the bridge from ``shorter window'' to ``higher flagged
weight'' that the legal clock supplies.

\emph{Step 13 (the composition ratio carries no sign).} No hypothesis of this theorem signs $C_T$.
Three reasons block a window analogue of Lemma~\ref{lem:L4}'s third leg, and any one of them
suffices.
\begin{enumerate}[label=(\roman*),leftmargin=2.4em,itemsep=3pt]
\item Assumption~\ref{asm:Abr} is quantified over the threshold pair. Clause~(br-i) asserts the
  representation under $\tau$ and under $\tau'$; clause~(br-iii) compares $\lvert\Akap(\tau')\rvert$
  with $\lvert\Akap(\tau)\rvert$; clause~(br-iv) asserts one endpoint function at $\tau$ and at
  $\tau'$; clause~(br-v) asserts one chord functional at $\tau$ and at $\tau'$. None of them says
  anything about two window environments, so Step~10's chain has no window instance to run on, and
  Lemma~\ref{lem:L4} has no window leg to cite.
\item The window has a channel the threshold does not. At fixed policies, changing $T$ changes the
  filing date $f_j = c_j + T$ and hence the objects clause~(TR-iii) of Assumption~\ref{asm:TR}
  already signs: $T' < T$ implies $B^F(T') \le B^F(T)$ and $Q^F(T') \ge Q^F(T)$. Trade moves from
  the pooled round into the flagged round on histories flagged under \emph{both} windows. A
  threshold change at a fixed window moves no unit of trade across the filing date in this way; it
  relabels which histories file at all.
\item The pooled cell's own information changes. The pooled public history of
  Definition~\ref{def:prices} is $\mathcal H_d^P = (X_0,\dots,X_d;\ \text{flag landed by }d)$, so a
  pooled history carries the event ``no flag has landed by $d$''. Under a tighter window that event
  excludes more crossing dates: it excludes $c \le d - T'$ rather than only $c \le d - T$, so
  the pooled cell's updating changes even holding the set of pooled histories fixed. Clause~(br-ii)
  requires the support points and the kernel not to move with $\kappa$; here the conditioning event
  itself moves with the policy, which is a composition change of the kind~(br-ii) excludes at the
  $\kappa$ margin and says nothing about at the $T$ margin.
\end{enumerate}

\emph{Step 14 (the exact finite equivalence).} Apply \eqref{eq:T1-factor} at $(\tau,T')$ and at
$(\tau,T)$, which~(T-2) permits at both. By~(T-3) and \eqref{eq:T1-factor},
$\mathcal S(\kappa,\tau,T) > 0$, so
\begin{equation}
\label{eq:T1-ratio-T}
\frac{\mathcal S(\kappa,\tau,T')}{\mathcal S(\kappa,\tau,T)}
\;=\; \frac{\bigl(1-\Omega(\tau,T')\bigr)\mathcal S_P(\kappa,\tau,T')}
            {\bigl(1-\Omega(\tau,T)\bigr)\mathcal S_P(\kappa,\tau,T)}
\;=\; W_T\,C_T,
\qquad
C_T = \frac{\mathcal S_P(\kappa,\tau,T')}{\mathcal S_P(\kappa,\tau,T)} .
\end{equation}
Multiplying an inequality by the strictly positive number $\mathcal S(\kappa,\tau,T)$ preserves it
in both directions, so
\begin{equation}
\label{eq:T1-iff}
\mathcal S(\kappa,\tau,T') \;\le\; \mathcal S(\kappa,\tau,T)
\qquad\Longleftrightarrow\qquad
W_T\,C_T \;\le\; 1 ,
\end{equation}
which is the equivalence of Part~(C). Both directions hold, and neither side is a sign claim: the
identity \eqref{eq:T1-ratio-T} is the falsifiable content at this margin, and the equivalence is
empty of economics until $C_T$ is measured. Since $W_T > 0$, \eqref{eq:T1-iff} rearranges to
$C_T \le 1/W_T = \bigl(1-\Omega(\tau,T)\bigr)/\bigl(1-\Omega(\tau,T')\bigr)$: attenuation holds
exactly when the composition effect does not exceed the reciprocal of the weight effect. That is
what ``the weight effect dominates the composition effect'' means here, and it means nothing more.
In particular it does not mean $\lvert 1-W_T\rvert \ge \lvert 1-C_T\rvert$, and it does not mean
$C_T \le 1$.

\emph{Step 15 (the local form).} Adopt the smooth window interpolation of hypothesis~(T-14), which
is consumed here and nowhere else: there is an open interval of window values containing $[T',T]$
and continuously differentiable extensions $r \mapsto \Omega(r)$ and
$r \mapsto \partial_\kappa M_P(r)$ of the model's integer-valued objects, agreeing with them at
$r = -T$ and $r = -T'$, with $\Omega(r) \in (0,1)$ and $\mathcal S_P(r) > 0$ throughout and with
$r \mapsto \Omega(r)$ weakly increasing. Write $r = r_T = -T$ as in Definition~\ref{def:theta}, so
higher $r$ is tighter, and set $r_0 = -T < r_1 = -T'$. On the interpolating interval
$\mathcal S_P(r) > 0$, so $x \mapsto |x|$ is differentiable at $\partial_\kappa M_P(r) \neq 0$ and
$\mathcal S_P$ is continuously differentiable in $r$ by Step~5's second consequence. Differentiating
\eqref{eq:T1-factor} in $r$,
\[
\partial_{r_T}\mathcal S \;=\; -\,\partial_{r_T}\Omega\,\mathcal S_P \;+\; (1-\Omega)\,\partial_{r_T}\mathcal S_P .
\]
The monotonicity clause of~(T-14) gives $\partial_{r_T}\Omega \ge 0$, so the first term is the attenuating
weight effect and the second is the unsigned composition effect. That sign comes from~(T-14) and not
from Step~12: Step~12 compares two \emph{integer} windows and delivers the inequality at the
endpoints only, and nothing outside~(T-14) forbids an interpolant that dips in between. Dividing by
the strictly positive $\mathcal S = (1-\Omega)\mathcal S_P$ and writing, for this proof only,
\begin{equation}
\label{eq:T1-Lambda}
\Lambda(r) \;:=\; \frac{\partial_{r}\mathcal S_P(r)}{\mathcal S_P(r)} \;-\; \frac{\partial_{r}\Omega(r)}{1-\Omega(r)},
\end{equation}
we get $\partial_{r_T}\mathcal S/\mathcal S = \Lambda(r)$, so
$\operatorname{sgn}(\partial_{r_T}\mathcal S) = \operatorname{sgn}(\Lambda)$ and
\begin{equation}
\label{eq:T1-local}
\partial_{r_T}\mathcal S \;\le\; 0
\qquad\Longleftrightarrow\qquad
\frac{\partial_{r_T}\mathcal S_P}{\mathcal S_P} \;\le\; \frac{\partial_{r_T}\Omega}{1-\Omega} ,
\end{equation}
an equivalence at each $r$, with no sign supplied for either side. Display \eqref{eq:T1-local} is
pure algebra (a division by a strictly positive number) and holds whether or not the
interpolant is monotone; only the \emph{reading} of the first term as attenuating uses
$\partial_{r_T}\Omega \ge 0$.

\emph{Step 16 (in what sense \eqref{eq:T1C} is the same criterion as \eqref{eq:T1-iff}).} The
statement of Part~(C) says that \eqref{eq:T1C} holds on average along the tightening path,
integrated over $[-T,-T']$, and exactly in the infinitesimal limit, and that it is false read
pointwise. Each of those three clauses is a separate assertion, and each is proved here. Let
$\Lambda$ be as in \eqref{eq:T1-Lambda}, continuous on $[r_0,r_1]$ by~(T-14).

\emph{(i) The finite product criterion is the sign of the integrated local difference.} By~(T-14),~(T-3)
and \eqref{eq:T1-factor}, $\mathcal S(r) > 0$ and is continuously differentiable on $[r_0,r_1]$, so
$\log\mathcal S$ is continuously differentiable there and the fundamental theorem of calculus gives
$\log\bigl(\mathcal S(r_1)/\mathcal S(r_0)\bigr) = \int_{r_0}^{r_1}\partial_r\log\mathcal S(r)\,dr
= \int_{r_0}^{r_1}\Lambda(r)\,dr$, the second equality by Step~15. By \eqref{eq:T1-ratio-T},
$\mathcal S(r_1)/\mathcal S(r_0) = W_T C_T$, so
\begin{equation}
\label{eq:T1-integral}
W_T\,C_T \;=\; \exp\!\left(\int_{r_0}^{r_1}\Lambda(r)\,dr\right),
\qquad\text{hence}\qquad
W_T\,C_T \le 1 \iff \int_{r_0}^{r_1}\Lambda(r)\,dr \le 0 .
\end{equation}
The exponential is strictly increasing, so the equivalence holds in both directions. Written out,
the integrand is the left side of \eqref{eq:T1C} minus its right side: the product criterion is
the local criterion integrated along the tightening path. This is the ``on average'' reading, and
it is exact at finite scale.

\emph{(ii) Pointwise local implies the finite product.} If $\Lambda(r) \le 0$ at every
$r \in [r_0,r_1]$, the integral of a nonpositive continuous function over that interval is
nonpositive, so $W_T C_T \le 1$ by \eqref{eq:T1-integral}. If in addition $\Lambda < 0$ on a
subinterval of positive length, the integral is strictly negative and $W_T C_T < 1$.

\emph{(iii) The finite product implies the local criterion at one point, and no more, so read
pointwise, \eqref{eq:T1C} is false.} If $W_T C_T \le 1$ then $\int_{r_0}^{r_1}\Lambda \le 0$ by
\eqref{eq:T1-integral}, and since $\Lambda$ is continuous the mean value theorem for integrals
supplies some $r^\star \in (r_0,r_1)$ with
$\Lambda(r^\star) = (r_1-r_0)^{-1}\int_{r_0}^{r_1}\Lambda \le 0$. The local criterion therefore
holds \emph{somewhere} in the interval. It does not hold everywhere in general, and the
counterexample shape is a configuration \eqref{eq:T1C} is not entitled to exclude: a $\Lambda$ that
is positive on $[r_0,\tfrac12(r_0+r_1)]$ and sufficiently negative on $[\tfrac12(r_0+r_1),r_1]$
integrates to a nonpositive number while violating the local criterion on the first half. The
finite comparison from $T$ to $T'$ then reports attenuation even though an intermediate window
tightening amplifies. The implication in this direction is ``at some point'', never ``at every
point''.

\emph{(iv) The two forms coincide exactly in the infinitesimal limit.} Fix $r_0$ and let
$r_1 \downarrow r_0$. By \eqref{eq:T1-ratio-T}, $W_T C_T = \mathcal S(r_1)/\mathcal S(r_0)$ as a
function of $r_1$, and it equals $1$ at $r_1 = r_0$, so
\[
\lim_{r_1\downarrow r_0}\frac{W_TC_T-1}{r_1-r_0}
\;=\; \frac{d}{dr_1}\left.\frac{\mathcal S(r_1)}{\mathcal S(r_0)}\right|_{r_1=r_0}
\;=\; \frac{\partial_r\mathcal S(r_0)}{\mathcal S(r_0)} \;=\; \Lambda(r_0).
\]
Hence if $\Lambda(r_0) < 0$ then $W_TC_T < 1$ for every $r_1 > r_0$ close enough to $r_0$; and if
$W_TC_T \le 1$ for every such $r_1$, the limit is at most zero, so $\Lambda(r_0) \le 0$. The
boundary case $\Lambda(r_0) = 0$ is undetermined at first order, the finite sign there being decided
by higher-order terms, and it is named rather than resolved.

\emph{Step 17 (no unconditional window sign).} Nothing in hypotheses~(T-1)--(T-15) signs $C_T$, by
Step~13, and nothing in them signs $\Lambda$, by Step~15. So nothing in them signs
$\partial_{r_T}\mathcal S$ or $W_TC_T - 1$. Both branches are consistent with every hypothesis
maintained here: $W_TC_T \le 1$ and $W_TC_T > 1$ each require only a value of $C_T$ that the
hypotheses leave free. Part~(C) is an equivalence and an identity, and the sign is an empirical
question about the composition ratio. This is what the statement's closing clause records, and it is
why the model does not deliver window-margin attenuation as a theorem.

Parts~(A), (B) and~(C) are Steps~5--6, Step~11, and Steps~12--17 respectively.
\end{proof}

\medskip
\noindent\emph{Scope of the argument above.}
Hypothesis~(T-15), threshold-side smoothness, is consumed by no step of the proof above. That is the
content of the statement's ``confirmed non-load-bearing'': Part~(B)'s conclusion \eqref{eq:T1B} is a
finite-difference comparison throughout, and Parts~(A) and~(C) never differentiate in $r_\tau$. A
smooth local reading of Part~(B) is available under~(T-15) and would say that
$\partial_{r_\tau}\mathcal S \le 0$ with slack on both sides of zero, which is the same statement as
``no dominance condition needed''; it is not claimed here, and if~(T-15) fails nothing above moves.

Two limits on the proof are worth restating where the argument is, rather than only beside the
statement. First, Part~(B)'s Steps~7 and~10 consume Assumption~\ref{asm:Atau} at both compared
policies and clauses~(br-i)--(br-v) of Assumption~\ref{asm:Abr}; Remark~\ref{rem:AtauRecord} records
that the support condition of $\Atau$ fails at every non-degenerate node of the implemented
calibration, and clause~(br-i) carries that representation at both thresholds, so at that
calibration Part~(B)'s antecedent is not satisfied and Part~(B) says nothing about the implemented
pooled cell. The implication is unaffected by that fact: the measurement bears on whether the
antecedent holds at one calibration, not on the proof. Second, the window-margin
evidence recorded at Remark~\ref{rem:T1record} is node evidence at one calibration; it is not a
sign, and Part~(C) claims none.

\begin{proof}[Proof of Proposition~\ref{prop:C1} (C1)]
Four conventions are fixed first, because the statement's hypotheses are stated against them.

\emph{The margin.} By hypothesis~(C-7) the strictness coordinate is the threshold one,
$r = r_\tau = -\tau$ of Definition~\ref{def:theta}, and $\vartheta = (\kappa,r)$.
Definition~\ref{def:primitives} restricts only the window to integers and places no discreteness on
$\tau$, so $r$ is a coordinate on an interval and the derivatives below have a domain to live on.
The window coordinate is excluded: $T$ takes values in $\{1,\dots,H\}$ there, so $\partial_{r_T}$,
$\partial_{\kappa r_T}\Dact$ and $g_{r_T}^{PE}$ have no meaning without an interpolation of the
model in $T$, and nothing local is claimed at that margin.

\emph{The norm.} By hypothesis~(C-1) one norm $\lVert\cdot\rVert$ is fixed on the cutoff space
$\mathbb R^{J-1}$, and every magnitude bar in the $\bar k_x$ and $\bar k_{\kappa r}$ rows of
Definition~\ref{def:theta}, in $L_{\mathcal R}$, and in \eqref{eq:BGE} is read as the member of the
family that norm induces: $\lVert\cdot\rVert$ itself for the vector-valued objects
$\partial_\kappa\Tmap$, $\partial_r\Tmap$ and $\Tmap_{\kappa r}$; the induced operator norm
$\sup_{\lVert w\rVert\le1}\lVert Aw\rVert$ for the matrices $D_k\Tmap$, $\Tmap_{\kappa k}$ and
$\Tmap_{rk}$; $\sup_{\lVert w_1\rVert,\lVert w_2\rVert\le1}\lVert\Tmap_{kk}[w_1,w_2]\rVert$ for the
vector-valued bilinear form $\Tmap_{kk}$; the \emph{dual} norm
$\lVert\phi\rVert_* = \sup_{\lVert w\rVert\le1}\lvert\phi(w)\rvert$ for the covectors $\Delta_k$,
$\Delta_{\kappa k}$ and $\Delta_{kr}$; and
$\sup_{\lVert w_1\rVert,\lVert w_2\rVert\le1}\lvert\Delta_{kk}[w_1,w_2]\rvert$ for the scalar
bilinear form $\Delta_{kk}$. Two consequences are used. The operator norm must be induced, hence
submultiplicative, because Step~1 needs $\lVert A^n\rVert \le \lVert A\rVert^n$; and the pairing
must be the matching one, because pairing a covector with a vector through a magnitude that is not
the dual norm can make \eqref{eq:BGE} \emph{understate} the remainder it is meant to bound, which
voids the implication silently. At $J-1 = 1$ the whole family collapses to the absolute value and
the convention is vacuous.

\emph{The equilibrium objects.} By hypothesis~(C-3) there is, at every $\vartheta$ in the region
$\mathcal R_r$, a cutoff vector $k(\vartheta)$ in the interior of $\Theta$ with
$k(\vartheta) = \Tmap(k(\vartheta);\vartheta)$, and $\vartheta\mapsto k(\vartheta)$ is one branch,
with no switch of selected equilibrium inside $\mathcal R_r$. Interiority is a genuine restriction:
Definition~\ref{def:cpbe}\,(i) has weak inequalities and so permits collapsed action regions, and a
collapsed region puts $k$ on the boundary of $\Theta$, where the fixed-point equation may hold only
as a one-sided condition. Write, for this proof,
\[
\mathcal D(\kappa,r) \;:=\; \Dact\bigl(k(\kappa,r),\kappa,r\bigr),
\qquad
\mathcal S^{GE} \;:=\; \Bigl\lvert \frac{d\mathcal D}{d\kappa} \Bigr\rvert ,
\]
so that $\mathcal D$ is the premium \eqref{eq:Dact} evaluated along the equilibrium branch, not
the disclosure indicator $D$, and $\mathcal S^{GE}$ is the equilibrium liquidity sensitivity of
hypothesis~(C-5). It is a different object from the fixed-policy $\mathcal S = |\partial_\kappa\Dact|$
of Definition~\ref{def:premium}, which holds $k$ frozen; the two agree only where the cutoff
response contributes nothing, and the entire content of this proposition sits in the difference.
Every partial derivative of $\Dact$ and of $\Tmap$ written below is a partial of the
\emph{fixed-policy} map $(k,\kappa,r)\mapsto\Dact(k,\kappa,r)$, respectively of
$(k,\kappa,r)\mapsto\Tmap(k;\kappa,r)$, evaluated at the equilibrium point
$(k(\vartheta),\vartheta)$.

\emph{The region and its regularity.} Hypothesis~(C-1) puts Assumption~\ref{asm:AGE} in force on
$\mathcal R_r$, and names the clause that carries the weight: the contraction bound
$L_{\mathcal R} = \sup_{\mathcal R}\lVert D_k\Tmap\rVert < 1$ of Definition~\ref{def:theta}, read
along the equilibrium path, that is, as the supremum over $\vartheta \in \mathcal R_r$ of
$\lVert D_k\Tmap(k(\vartheta);\vartheta)\rVert$. That reading is the one used at Steps~1--4 and it
is the weaker of the two available; the stronger reading, a supremum over all of $\Theta$ as well,
is not consumed anywhere below. The differentiability clause of Assumption~\ref{asm:AGE} supplies
twice continuous differentiability of $\Tmap$ on a neighbourhood of the equilibrium graph
$\{(k(\vartheta),\vartheta) : \vartheta\in\mathcal R_r\}$, which is the weakest domain Steps~1--4
use; smoothness away from that graph is never called on. Hypothesis~(C-4) supplies twice continuous
differentiability of $\Dact$ in $(k,\kappa,r)$ on the same neighbourhood, and with it finiteness of
the derivatives named in \eqref{eq:BGE} and in $\bar k_{\kappa r}$. The integrability clause of
Assumption~\ref{asm:A2p} is not cited for that and cannot be: boundedness is not differentiability.
Finally, hypothesis~(C-2) makes $\mathcal R_r$ relatively open in both coordinates with
$\kappa\notin\{0,1\}$, so that every $\vartheta\in\mathcal R_r$ has a full two-dimensional
neighbourhood inside $\mathcal R_r$, which the implicit function theorem of Step~2, the mixed
partial of Step~5 and the sign argument of Step~7 all need, and which openness in $r$ alone would
not give, $\kappa$ ranging over $[0,1]$ in Definition~\ref{def:primitives}.

\emph{Step 1 (invertibility, without forming an inverse).} Fix $\vartheta\in\mathcal R_r$ and write
$A = D_k\Tmap(k(\vartheta);\vartheta)$, so $\lVert A\rVert \le L_{\mathcal R} < 1$ by~(C-1). By
submultiplicativity, $\lVert A^n\rVert \le L_{\mathcal R}^{\,n}$, so
$\sum_{n\ge0}\lVert A^n\rVert \le (1-L_{\mathcal R})^{-1} < \infty$. The space of
$(J-1)\times(J-1)$ matrices being finite dimensional and complete, the partial sums
$S_N = \sum_{n=0}^{N}A^n$ converge to a limit $S$ with
$\lVert S\rVert \le (1-L_{\mathcal R})^{-1}$; and from $(I-A)S_N = S_N(I-A) = I - A^{N+1}$ with
$\lVert A^{N+1}\rVert \le L_{\mathcal R}^{\,N+1} \to 0$, passing to the limit gives
$(I-A)S = S(I-A) = I$. Hence $I - D_k\Tmap$ is invertible at every point of the equilibrium graph,
with
\begin{equation}
\label{eq:C1-neumann}
\bigl\lVert (I-D_k\Tmap)^{-1} \bigr\rVert \;\le\; \frac{1}{1-L_{\mathcal R}} .
\end{equation}
The bound is a geometric sum of norms, so every derivative bound below is available from
$\lVert D_k\Tmap\rVert$ and one directional derivative of $\Tmap$, with no linear system solved.

\emph{Step 2 (the equilibrium branch is twice continuously differentiable).} Fix
$\vartheta\in\mathcal R_r$ and write $k^0 = k(\vartheta)$. Pick $L'$ with
$L_{\mathcal R} < L' < 1$. By the differentiability clause of Assumption~\ref{asm:AGE} the map
$(k,\vartheta')\mapsto D_k\Tmap(k;\vartheta')$ is continuous, and by~(C-3) and~(C-2) there are
$\delta,\delta_\vartheta > 0$ with the closed ball $\bar B(k^0,\delta)$ inside the interior of $\Theta$, the
closed ball $\bar B(\vartheta,\delta_\vartheta)$ inside $\mathcal R_r$ and inside the domain on which
Assumption~\ref{asm:AGE} and~(C-4) hold, and $\lVert D_k\Tmap(k';\vartheta')\rVert \le L'$
throughout $\bar B(k^0,\delta)\times\bar B(\vartheta,\delta_\vartheta)$. The ball being convex, the
mean-value inequality along the segment gives
$\lVert \Tmap(k_1;\vartheta') - \Tmap(k_2;\vartheta')\rVert \le L'\lVert k_1-k_2\rVert$ there.
Shrinking $\delta_\vartheta$ so that
$\lVert \Tmap(k^0;\vartheta') - \Tmap(k^0;\vartheta)\rVert \le (1-L')\delta$, and using
$\Tmap(k^0;\vartheta) = k^0$ from~(C-3), gives
$\lVert \Tmap(k_1;\vartheta') - k^0\rVert \le L'\delta + (1-L')\delta = \delta$ for every
$k_1\in\bar B(k^0,\delta)$ and $\vartheta'\in\bar B(\vartheta,\delta_\vartheta)$. So
$\Tmap(\cdot;\vartheta')$ maps the closed ball into itself and contracts it. That ball is a complete
metric space, so by the contraction mapping theorem each such $\vartheta'$ admits exactly one fixed
point $\hat k(\vartheta')$ of $\Tmap(\cdot;\vartheta')$ inside the ball, with
$\hat k(\vartheta) = k^0$; and the standard estimate
$\lVert\hat k(\vartheta')-\hat k(\vartheta'')\rVert \le (1-L')^{-1}\lVert
\Tmap(\hat k(\vartheta'');\vartheta') - \Tmap(\hat k(\vartheta'');\vartheta'')\rVert$ makes
$\hat k$ continuous. Local uniqueness and local continuity are consequences here, not hypotheses.

Now put $\Psi(k',\vartheta') := k' - \Tmap(k';\vartheta')$ on
$\bar B(k^0,\delta)\times\bar B(\vartheta,\delta_\vartheta)$. It is twice continuously differentiable by
Assumption~\ref{asm:AGE}, it vanishes at $(k^0,\vartheta)$ by~(C-3), and
$D_k\Psi = I - D_k\Tmap$ is invertible there by Step~1. The implicit function theorem in its twice
continuously differentiable form supplies neighbourhoods and a twice continuously differentiable map
$\vartheta'\mapsto\tilde k(\vartheta')$ solving $\Psi(\tilde k(\vartheta'),\vartheta') = 0$ with
$\tilde k(\vartheta) = k^0$, unique among solutions there; and since every such $\tilde k$ is a
fixed point of $\Tmap(\cdot;\vartheta')$ inside the ball, $\tilde k = \hat k$ where both are
defined. Hypothesis~(C-3)'s single-branch clause is consumed exactly here: the selected
$\vartheta'\mapsto k(\vartheta')$ does not jump to a different fixed point at nearby parameters, so
it stays in $\bar B(k^0,\delta)$ near $\vartheta$ and coincides there with $\hat k = \tilde k$. As
$\vartheta$ was an arbitrary point of the relatively open $\mathcal R_r$, the branch $k(\cdot)$ is
twice continuously differentiable on $\mathcal R_r$.

\emph{Step 3 (the first-derivative bounds).} Let $x\in\{\kappa,r\}$. Differentiating the identity
$k(\vartheta) = \Tmap(k(\vartheta);\vartheta)$ in $x$, which Step~2 licenses and which~(C-2) makes
meaningful in both coordinates, and applying the chain rule in the first argument,
$\partial_x k = D_k\Tmap\,\partial_x k + \partial_x\Tmap$, that is
$(I-D_k\Tmap)\,\partial_x k = \partial_x\Tmap$. Taking norms and using
$\lVert\partial_x k\rVert \le L_{\mathcal R}\lVert\partial_x k\rVert + \lVert\partial_x\Tmap\rVert$,
which may be rearranged because $\lVert\partial_x k\rVert$ is finite by Step~2 and
$L_{\mathcal R} < 1$,
\begin{equation}
\label{eq:C1-kbar}
\lVert\partial_\kappa k\rVert \;\le\; \bar k_\kappa = \frac{\lvert\partial_\kappa\Tmap\rvert}{1-L_{\mathcal R}},
\qquad
\lVert\partial_r k\rVert \;\le\; \bar k_r = \frac{\lvert\partial_r\Tmap\rvert}{1-L_{\mathcal R}},
\end{equation}
which are the $\bar k_x$ rows of Definition~\ref{def:theta}. The derivative on the right is the
\emph{partial} derivative of the outer map in the parameter at frozen cutoffs, not a total
derivative.

\emph{Step 4 (the cross-derivative bound).} Differentiate the $x = \kappa$ instance of Step~3 once
more in $r$, with every $\Tmap$ derivative evaluated at $(k(\vartheta),\vartheta)$ and the chain
rule applied through $k(\vartheta)$ in each slot. In components,
\[
\partial^2_{\kappa r}k^i
= \sum_j\Bigl[\Bigl(\sum_l\Tmap^i_{k_jk_l}\,\partial_r k^l + \Tmap^i_{k_jr}\Bigr)\partial_\kappa k^j
+ \Tmap^i_{k_j}\,\partial^2_{\kappa r}k^j\Bigr]
+ \sum_l\Tmap^i_{\kappa k_l}\,\partial_r k^l + \Tmap^i_{\kappa r},
\]
that is,
$(I-D_k\Tmap)\,\partial^2_{\kappa r}k = \Tmap_{\kappa r} + \Tmap_{\kappa k}[\partial_r k]
+ \Tmap_{rk}[\partial_\kappa k] + \Tmap_{kk}[\partial_r k,\partial_\kappa k]$. Taking norms on the
right with the pairings fixed above and the bounds \eqref{eq:C1-kbar}, then rearranging as in
Step~3,
\begin{equation}
\label{eq:C1-kkr}
\lVert\partial^2_{\kappa r}k\rVert
\;\le\; \bar k_{\kappa r}
= \frac{\lvert\Tmap_{\kappa r}\rvert + \lvert\Tmap_{\kappa k}\rvert\,\bar k_r
+ \lvert\Tmap_{rk}\rvert\,\bar k_\kappa + \lvert\Tmap_{kk}\rvert\,\bar k_\kappa\bar k_r}
{1-L_{\mathcal R}},
\end{equation}
which is the $\bar k_{\kappa r}$ row of Definition~\ref{def:theta}.

\emph{Step 5 (the exact cross-derivative decomposition).} By~(C-4) and Step~2 the composition
$\mathcal D$ is twice continuously differentiable on $\mathcal R_r$, and~(C-2) supplies the
two-dimensional domain the mixed partial needs. The chain rule in $\kappa$ gives
$\partial_\kappa\mathcal D = \Delta_k\cdot\partial_\kappa k + \Delta_\kappa$, and differentiating
that in $r$ term by term, with
$\partial_r(\Delta_k) = \Delta_{kk}[\,\cdot\,,\partial_r k] + \Delta_{kr}$ and
$\partial_r(\Delta_\kappa) = \Delta_{\kappa k}[\partial_r k] + \Delta_{\kappa r}$, gives the exact
identity
\begin{equation}
\label{eq:C1-decomp}
\frac{d^2\mathcal D}{d\kappa\,dr}
\;=\; \underbrace{\Delta_{\kappa r}}_{\text{fixed policy}}
\;+\; \underbrace{\Delta_{\kappa k}[\partial_r k]
\;+\; \bigl(\Delta_{kr} + \Delta_{kk}[\partial_r k]\bigr)\cdot\partial_\kappa k
\;+\; \Delta_k\cdot\partial^2_{\kappa r}k}_{\text{general-equilibrium remainder}} .
\end{equation}
Nothing is dropped and no inequality has been used: \eqref{eq:C1-decomp} is one identity in five
terms, grouped in four. Two features of it are where a miscount would hide. No term
$\Delta_{\kappa\kappa}$ or $\Delta_{rr}$ can appear, the derivative being mixed and each coordinate
differentiated once; and the last term is the only place a \emph{second} derivative of the
equilibrium map enters, which is why Step~4 is needed at all. The remainder is the object
Theorem~\ref{thm:T1} sets its equilibrium free of: the term
$(\partial_\kappa\Omega)(M_F - M_P)$ discarded at \eqref{eq:T1-three} under hypothesis~(T-7) is a
piece of $\Delta_k\cdot\partial_\kappa k$, and it reappears here.

\emph{Step 6 (the remainder is bounded by $\mathcal B_r^{GE}$).} Write $\mathcal E$ for the
remainder group of \eqref{eq:C1-decomp}, so that
$\mathcal E = d^2\mathcal D/d\kappa\,dr - \partial_{\kappa r}\Dact$. Apply the triangle inequality
to its four terms, bound each with the pairings fixed above (covector against vector through the
dual norm, bilinear form against two vectors) and insert \eqref{eq:C1-kbar} and
\eqref{eq:C1-kkr}, every factor being finite by~(C-4):
\[
\bigl\lvert\Delta_{\kappa k}[\partial_r k]\bigr\rvert \le \lvert\Delta_{\kappa k}\rvert\bar k_r,
\qquad
\bigl\lvert\bigl(\Delta_{kr}+\Delta_{kk}[\partial_r k]\bigr)\cdot\partial_\kappa k\bigr\rvert
\le \bigl(\lvert\Delta_{kr}\rvert + \lvert\Delta_{kk}\rvert\bar k_r\bigr)\bar k_\kappa,
\qquad
\bigl\lvert\Delta_k\cdot\partial^2_{\kappa r}k\bigr\rvert \le \lvert\Delta_k\rvert\bar k_{\kappa r}.
\]
Summing the three,
\begin{equation}
\label{eq:C1-remainder}
\bigl\lvert \mathcal E \bigr\rvert
\;\le\; \lvert\Delta_{\kappa k}\rvert\bar k_r
+ \bigl(\lvert\Delta_{kr}\rvert + \lvert\Delta_{kk}\rvert\bar k_r\bigr)\bar k_\kappa
+ \lvert\Delta_k\rvert\bar k_{\kappa r}
\;=\; \mathcal B_r^{GE},
\end{equation}
which is \eqref{eq:BGE}. The bookkeeping is worth auditing, since this is where the argument turns:
\eqref{eq:BGE} has three groups and $\mathcal E$ has four terms, mapped one to the first group, two
to the second, and one to the third; every cross-term of the mixed second derivative lands in
exactly one group, no group is fed by a term outside $\mathcal E$, and the only term of
\eqref{eq:C1-decomp} outside the groups is the fixed-policy cross-derivative that
\eqref{eq:C1-remainder} subtracts off. The only inequalities used are the triangle inequality and
the norm pairings. One consequence of \eqref{eq:C1-kbar} and \eqref{eq:C1-kkr} is worth naming: each
$\bar k_x$ carries one factor of $(1-L_{\mathcal R})^{-1}$ and $\bar k_{\kappa r}$ carries three, so
$\mathcal B_r^{GE} = O\bigl((1-L_{\mathcal R})^{-3}\bigr)$ as $L_{\mathcal R}\uparrow1$.

\emph{Step 7 (differentiating the equilibrium sensitivity).} Write $f(\vartheta') :=
d\mathcal D/d\kappa$ at $\vartheta'$, which is continuously differentiable on $\mathcal R_r$ by
Step~5. By hypothesis~(C-5), $f$ does not vanish on $\mathcal R_r$. Take a connected relatively open
neighbourhood $\mathcal N\subseteq\mathcal R_r$ of $\vartheta$, which~(C-2) supplies. A continuous
nowhere-zero real function on a connected set takes values of one sign: otherwise the intermediate
value theorem, applied along a path in $\mathcal N$ joining a point where $f>0$ to one where $f<0$,
would produce a zero of $f$ inside $\mathcal N$, contradicting~(C-5). So
$\operatorname{sgn} f$ is constant on $\mathcal N$; write $\mathfrak s\in\{-1,+1\}$ for that
constant. On $\mathcal N$,
$\mathcal S^{GE} = \lvert f\rvert = \mathfrak s\,f$, an identity of functions and not merely of
values at one point, and the right side is continuously differentiable as a constant multiple of a
continuously differentiable function. Hence $\mathcal S^{GE}$ is differentiable at $\vartheta$ with
\begin{equation}
\label{eq:C1-deriv}
\partial_r\mathcal S^{GE} \;=\; \mathfrak s\,\frac{d^2\mathcal D}{d\kappa\,dr} .
\end{equation}
What~(C-5) does here is remove the single point at which $x\mapsto|x|$ is not differentiable;
without it the conclusion is false, because at a parameter where the equilibrium liquidity
derivative vanishes $\mathcal S^{GE}$ has a kink and $\partial_r\mathcal S^{GE}$ does not exist. The
constancy of the sign is \emph{derived} on $\mathcal N$ from~(C-5) and~(C-2) and is not assumed.

\emph{Step 8 (splitting off the direct margin).} Substitute \eqref{eq:C1-decomp} into
\eqref{eq:C1-deriv} and use $\mathfrak s^2 = 1$. The fixed-policy group is read through
\eqref{eq:gPE}, whose sign factor is the sign of the \emph{equilibrium} liquidity derivative
$d\Dact/d\kappa$ evaluated along $k(\kappa,r)$, that is, $\mathfrak s$, while its
cross-derivative factor $\partial_{\kappa r}\Dact$ is the fixed-policy partial at
$(k(\vartheta),\vartheta)$. Hence $\mathfrak s\,\Delta_{\kappa r} = -g_r^{PE}$, and
\begin{equation}
\label{eq:C1-split}
\partial_r\mathcal S^{GE} \;=\; -\,g_r^{PE} \;+\; \mathfrak s\,\mathcal E ,
\end{equation}
every object in which is finite by~(C-4).

\emph{Step 9 (dominance closes the argument).} Multiplying by a factor of modulus one leaves the
modulus alone, so $\lvert \mathfrak s\,\mathcal E\rvert = \lvert\mathcal E\rvert \le
\mathcal B_r^{GE}$ by \eqref{eq:C1-remainder}; only the one-sided half
$\mathfrak s\,\mathcal E \le \mathcal B_r^{GE}$ is consumed. With \eqref{eq:C1-split},
\begin{equation}
\label{eq:C1-eta}
\partial_r\mathcal S^{GE}
\;\le\; -\,g_r^{PE} + \mathcal B_r^{GE}
\;=\; -\bigl(g_r^{PE} - \mathcal B_r^{GE}\bigr)
\;=\; -\,\eta_r ,
\end{equation}
with $\eta_r$ the slack of Definition~\ref{def:theta}; and $\eta_r > 0$ at every
$\vartheta\in\mathcal R_r$ by the strict dominance of hypothesis~(C-6). Therefore
$\partial_r\mathcal S^{GE} \le -\eta_r < 0$ at every $\vartheta\in\mathcal R_r$, which is
\eqref{eq:C1}. The argument never needs the remainder's sign, only its size, which is the point of
bounding the general-equilibrium channel rather than signing it.
\end{proof}

\medskip
\noindent\emph{Sign coherence, the norm, and what the argument above does not deliver.}
Three things are worth stating beside the argument rather than only beside the statement.

\emph{Sign coherence is not consumed.} Neither reading of the sign clause reaches \eqref{eq:C1}. The
constancy clause of Assumption~\ref{asm:AGE}, that the sign of the equilibrium liquidity
derivative is constant on the region, is not used: Step~7 derives constancy on a connected
neighbourhood from hypotheses~(C-5) and~(C-2), and \eqref{eq:C1-deriv} is a pointwise statement, so
the clause would matter only for naming one sign across a disconnected region. Nor is the coherence
of the fixed-policy sign with the equilibrium sign used. That coherence has one job and it is a job
of naming: at frozen $k$ the fixed-policy sensitivity $\mathcal S$ of Definition~\ref{def:premium}
has $r$-derivative $\operatorname{sgn}(\partial_\kappa\Dact)\,\partial_{\kappa r}\Dact$, and only
when that leading sign equals $\mathfrak s$ does the fixed-policy attenuation rate equal $-g_r^{PE}$
and the phrase ``the fixed-policy attenuation sign survives in equilibrium'' describe
\eqref{eq:C1}. Steps~1--9 are untouched either way, because \eqref{eq:C1-split} pulls out
$\mathfrak s$ and never the fixed-policy sign.

\emph{The hypotheses are relative to the norm.} Every quantity in \eqref{eq:C1-neumann} through
\eqref{eq:C1-eta} other than $\partial_{\kappa r}\Dact$ and $\mathfrak s$ depends on the norm fixed
by hypothesis~(C-1): $L_{\mathcal R}$, the bounds $\bar k_x$ and $\bar k_{\kappa r}$,
$\mathcal B_r^{GE}$, and hence $\eta_r$. A matrix with spectral radius below one has induced
operator norm below one in some norm but not necessarily in a given one, so both the hypothesis
$L_{\mathcal R} < 1$ and the slack in the conclusion are properties of the pair (region, norm), and
a rescaling of the cutoff coordinates can move a parameter vector across the boundary in either
direction. This is not a free lunch, since a change of norm also moves the dual norms inside
$\mathcal B_r^{GE}$; it is a reason to read the hypothesis existentially in the norm and to fix one
convention and keep to it, which is what~(C-1) does.

\emph{What is not delivered.} Nonemptiness of $\mathcal R_r$ is neither claimed nor provable here:
the region is a hypothesis, and whether any parameter vector satisfies (C-1)--(C-7) together is a
property of the calibration. Remark~\ref{rem:C1record} reports the dominance-and-contraction nodes
found by the numerical check, and what those nodes verify is narrow: the two pointwise
inequalities $L_{\mathcal R} < 1$ and $\eta_r > 0$ together with supporting diagnostics, and not the
full antecedent (C-1)--(C-7) and not a named nonempty region. Promoting nodes to a region needs a
modulus of continuity for $\eta_r$ on the hull of adjacent nodes and a genuine supremum for
$L_{\mathcal R}$ over the region, neither of which more grid nodes can supply. Finally, because
$\mathcal B_r^{GE}$ is a triangle-inequality bound rather than a tight one, $\eta_r \le 0$ carries
no information about the sign of $\partial_r\mathcal S^{GE}$: a failure of dominance is compatible
with equilibrium attenuation, with equilibrium amplification, and with equality.

\end{document}
